\section*{Working orientation for v\docversion}
\addcontentsline{toc}{section}{Working orientation for v\docversion}
\begin{quote}\small
\noindent\textbf{Reader's map.} The active reading path of the present document is: \emph{orientation and compact front extraction} $\to$ \emph{QAM formal layer} $\to$ \emph{tool and proof-compression layer} $\to$ \emph{open-question / live-burden layer} $\to$ \emph{computational and decoder layer}. The later appendix-style historical material remains available for proof-memory, genealogy, and local rebuild support, but it is not the governing active theorem core.
\end{quote}
This pre-v4 release proceeds directly from the abstract and contents to the nomenclature. The governing current-state claim is the groundwork-complete pre-OP frontier synthesis stated on the title page, in the abstract, and in the current-state audit blocks. Preserved release-roadmap notes, checkpoint transitions, and earlier front-matter control blocks remain part of the document as inherited baseline material for theorem-memory, genealogy, and reconstruction support; they are no longer to be read as the overriding current-working claim of v\docversion.

\paragraph{r06 bibliography hardening.}
A full-document literature scan was performed before this public split-bundle release. The scan found that the mathematical body repeatedly uses Teichmuller and quasiconformal language, Beltrami-style deformation intuition, Riemann--GUE/RMT spectral heuristics, Hopf/fibre-bundle topology, Dixon/octonionic and division-algebraic language, Dirac/Higgs/Yang--Mills and CKM/PMNS landmarks, Julia-boundary diagnostics, ZF/QAM coding language, and probe/cryptographic/computational analogies. Version~\docversion\ therefore adds a broader terminal bibliography and uses this paragraph as the compact external anchor surface: Teichmuller/quasiconformal background is anchored in \cite{Ahlfors2006,GardinerLakic2000,Hubbard2006,Sullivan1985}; Riemann-zeta, spectral, and GUE/RMT background in \cite{Riemann1859,Edwards1974,Titchmarsh1986,Conrey2003,Montgomery1973,Odlyzko1987,Mehta2004,KatzSarnak1999}; Hopf, fibre-bundle, and exceptional-geometry background in \cite{Hopf1931,Steenrod1951,Hatcher2002,Bryant1987,EguchiGilkeyHanson1980}; division-algebraic and octonionic background in \cite{Baez2002,Dixon1994,Furey2018}; QFT/gauge and mixing landmarks in \cite{Dirac1928,YangMills1954,EnglertBrout1964,Higgs1964,Cabibbo1963,KobayashiMaskawa1973,MakiNakagawaSakata1962}; complex-dynamical Julia background in \cite{Milnor2006}; ZF/set-theoretic foundations in \cite{Zermelo1908,Fraenkel1928,Jech2003,Kunen1980}; and the computational/probe/cryptographic side in \cite{Turing1936,Kasiski1863,RivestShamirAdleman1978,DiffieHellman1976,Norris1997,Glasserman2004,Bollerslev1986,Knuth1997,Rojas1998,Copeland2006}.

\paragraph{r08 MML/MMA compression map.}
The r07 local language interface is now used as a conservative compression and audit map for inherited HC/OP text. The purpose is not to delete historical proof material, but to mark repeated structural paragraphs as instances of typed operations such as \(\mathrm{CAPTURE}_{\mathrm{HC}}\), \(\mathrm{HANDOFF}_{\mathrm{HC}\to\mathrm{OP}}\), \(\mathrm{NORMALIZE}_{\mathrm{OP}}^{\mathrm{HC}}\), \(\mathrm{ROUTE}_{\mathrm{OP}}^{\mathrm{HC}}\), \(\mathrm{COMPARE}_{\mathrm{HC}}^{1/4\mid 5B}\), \(\mathrm{RESOLVE}\), \(\mathrm{RETURN}_{\mathrm{OP}}\), \(\mathrm{SCAN}\), \(\mathrm{RANK}\), and \(\mathrm{CERTIFY?}\). The compression map is therefore a reading and refactoring guide: it can shorten later exposition and make inherited blocks easier to audit, but it cannot manufacture certificates or remove the historical witness.

\paragraph{r10 first tool-contract application.}
The r09 tool library is now applied back to the active HC/OP proof spine. The purpose is deliberately modest: repeated paragraphs in the r01--r04 chain may be referenced by the contracts \(\mathfrak T_{\mathrm{lock}}\), \(\mathfrak T_{\mathrm{handoff}}\), \(\mathfrak T_{\mathrm{norm}}\), \(\mathfrak T_{\mathrm{route}}\), \(\mathfrak T_{\mathrm{cmp}}\), \(\mathfrak T_{\mathrm{resolve}}\), \(\mathfrak T_{\mathrm{cert}}\), and \(\mathfrak T_{\mathrm{frontier}}\) once their local hypotheses and guards have been stated. This is a citation and audit layer, not a shortening pass that deletes proof content.


\paragraph{r11 guarded reduction protocol.}
The present revision adds the first reduction discipline for using those tool contracts in future text. A paragraph is reducible only when it repeats an already stated safety, route, frontier, diagnostic, or no-false-certificate clause. It is not reducible when it introduces a new object, a first hypothesis, a local guard, a certificate test, a proof restriction, a theorematic status change, or historical witness material. Thus the r11 layer is the first practical simplifier produced by MML/MMA, but it remains conservative: it turns repeated prose into guarded citations and never turns an execution trace into an OP closure certificate.


\paragraph{r12 guarded-reduction application pass.}
The present revision applies the r11 discipline for the first time without deleting text. It installs a finite ledger of guarded citation sites along the active HC/OP proof spine: lock/capture reminders, handoff and normalization reminders, route-frontier reminders, paired-comparison status clauses, probe diagnostic clauses, and archive-to-tool extraction reminders. Each entry keeps its first occurrence and local guard explicit, records the relevant tool contract, and permits only later repeated tails to be cited in guarded normal form. Thus r12 is a practical audit layer: it tells the reader exactly where the new compression technology helps the existing text, while preserving theorematic status and the historical witness.


\paragraph{r13 guarded-reduction pilot normal form.}
The present revision turns the r12 ledger into a first usable rewriting pattern. It does not shorten the historical witness and does not replace first definitions. It only says that, after a first occurrence has stated its hypotheses and local guard, a later repeated HC/OP tail may be written in the guarded citation form
\[
  \mathrm{GRCite}_{\mathrm{HCOP}}(\ell;G,W,S),
\]
where \(\ell\) names the ledger entry, \(G\) points to the local guard, \(W\) points to the preserved witness or first occurrence, and \(S\) records the non-certificate status. This is the first concrete simplification normal form produced by the MML/MMA layer: it compresses repeated prose, not proof content.

\paragraph{r14 tail-normalized proof skeleton.}
The present revision organizes the r13 guarded-citation mechanism into a first tail-normalized proof skeleton for the active HC/OP spine. Its purpose is to show where the first explicit mathematical source block remains, where later repeated endings may be cited in guarded form, and how the dependency order of lock/capture, handoff, normalization, routing, paired comparison, and resolution is preserved. Thus r14 does not shorten the document by deletion. It gives future HC/OP text a compact reference skeleton in which repeated tails can be replaced by guarded citations only after their source block, local guard, witness, and theorematic status have already appeared.
\paragraph{r15 skeleton-to-workstep dispatch.}
The present revision turns the r14 tail skeleton into a first controlled dispatch object. Its purpose is not to claim an OP closure, but to say which finite next work packets may be launched from the ordered skeleton nodes. The admissible dispatch order remains lock/capture, handoff, normalization, routing, comparison, and resolution; each dispatched packet must keep its local guard, witness, route tag, and non-certificate status unless a separate explicit certificate test is later supplied.

\paragraph{r16 first dispatched lock/capture stabilization.}
The present revision executes the first dispatched work packet from the r15 menu. It acts only on the LOCK/CAPTURE node, records a finite lock-stability defect register, and returns one conservative status: handoff-ready stable source, bounded local repair source, or smaller frontier source. This is the first point at which the skeleton-to-workstep dispatch becomes an actual controlled mathematical work step; it still does not generate an OP closure certificate.


\paragraph{r07 QAM/MML/MMA local interface registry.}
Version~\docversion\ records the narrow local interface stamps for the machine-language layer in prose and uses symbolic local interface tags in formulas for the current HC-routed OP-entry chain. The historical computational chapter already contains the broader MML--MMA--MCP reserve architecture. The present r07 step does not activate that reserve as an independent proof engine. It only extracts the finite sublanguage needed to read the r01--r04 chain as a typed program:
\[
  \mathrm{QAM}_{\mathrm{HCOP}}^{\mathrm{loc}}
  \longrightarrow
  \mathrm{MML}_{\mathrm{HCOP}}^{\mathrm{loc}}
  \longrightarrow
  \mathrm{MMA}_{\mathrm{HCOP}}^{\mathrm{loc}}.
\]
Here \(\mathrm{QAM}_{\mathrm{HCOP}}^{\mathrm{loc}}\) is the local bridge tag over the inherited QAM-over-ZF object language, with interface stamp QAM-HC/OP 0.1 recorded as metadata; \(\mathrm{MML}_{\mathrm{HCOP}}^{\mathrm{loc}}\) is the finite instruction interface for capture, handoff, normalization, routing, comparison, and return, with interface stamp MML-HC/OP 0.2 recorded as metadata; and \(\mathrm{MMA}_{\mathrm{HCOP}}^{\mathrm{loc}}\) is the route-faithful execution semantics for those instructions on finite HC/OP states, with interface stamp MMA-HC/OP 0.2 recorded as metadata.


The following inherited front block should therefore be read as the carried-forward active baseline from the audited v3.30.0 line, with the following priorities and outcomes:
\begin{enumerate}[leftmargin=2em]
  \item \textbf{Preserve the stabilized closure/interface baseline.} The language/closure, branch/shell/routed, tensor, operator-family, translation, OP~3 shell-certificate, bundled interface compatibility, and terminal proof-use surface already exported in earlier public stages are treated as fixed baseline rather than reopened here.
  \item \textbf{Keep the partial-closure/no-second-carrier layer fixed.} The terminal-interface no-second-carrier theorem, the first partial OP~3/OP~5 closure theorem, and the restricted lift-back principle are treated as available structural infrastructure.
  \item \textbf{Keep the first retained-carrier optical approximation and machine-shadow baseline fixed.} Reduced prism readout, residual diagnostics, carrier-side fit/calibration/validation, inversion, refinement, reduced stability, and the bidirectional observer-side split \(\eta_+=\eta_{\mathrm{drive}}+\eta_{\mathrm{carrier}}\), \(\eta_-=\eta_{\mathrm{drive}}-\eta_{\mathrm{carrier}}\) are treated as the stabilized computational surface rather than rebuilt from scratch.
  \item \textbf{Inherit the completed tool/QAM/decoder layer as fixed infrastructure.} Channel-sorted total, imbalance, and bridge packets, finite QAM submenus, primary tags, registry assignments, tag-first reopening routes, the Monte-Carlo residual-signature routing layer, the action table, deterministic decoder syntax, transition function, observer/event normal form, evaluation map, and compact output tokens are treated as the completed route-control package inherited from v3.28.0.
  \item \textbf{Keep the prior post-closure burden reading visible as inherited baseline.} The v3.29.0--v3.30.0 burden picture is preserved here as inherited local baseline material from which the new pre-OP frontier synthesis is now read, not as an independent current public-claim layer that overrides the new groundwork status stated in the opening pages.
  \item \textbf{Keep the DOI logic normalized.} The concept DOI remains the stable title-page pointer of the Companion line. The governing latest public-release DOI in the chronology is now \href{https://doi.org/10.5281/zenodo.19750674}{10.5281/zenodo.19750674} for the public v3.35.0 release, while \href{https://doi.org/10.5281/zenodo.19746821}{10.5281/zenodo.19746821} remains the prior public release DOI for v3.34.0.
\end{enumerate}

\paragraph{Reading rule for the inherited active-front block.}
The inherited v3.30.0 front layer is kept intact as carried-forward baseline source material. It records the then-active retained-carrier optical and live-burden reading, including OP~1 and OP~5 Step~(B) closure language on the inherited admissible corridors, and should not be read as the governing present-release claim of v\docversion. In the present release final, those statements are preserved for theorem-memory and genealogy; the governing current-state claim, however, remains the frontier-certified pre-entry synthesis developed in the new post-v3.31.0 line. Accordingly, references below to live OP burden, route sharpening, kernel closure, or routed-balance obstruction inside the inherited front block should be read as preserved baseline source material rather than as the final current release claim of this release.

\subsection*{Compact public-shape branch/shell/routed extraction}
\begin{quote}\small
\noindent\textbf{Status.} \emph{Active front-facing compression block.}\par
\noindent\textbf{Block function.} This block gives the smallest front-facing statement of the public branch/shell/routed QAM package, so that a reader can see what is already stabilized before entering the fuller internal unfolding.
\end{quote}
The inherited public-shape QAM layer may be stated in compressed form before the full internal unfolding is read. In that form, one starts with the already public coded source/readout, branch/shell/routed, tensor, operator-family, and translation layers from v3.15.0--v3.17.0, then adds shell certificates, shell-defect witnesses, kernel-persistence statements, and shell-stable readout factorization. The guiding claim is that these added roles do not create a new closure criterion; they only make the OP~3 shell layer readable as one explicit coded closure mechanism.

\begin{definition}[Compact public-shape QAM package]
The \emph{compact public-shape QAM package} for the inherited public-shape v3.16.0 release layer consists of the following release-level data:
\begin{enumerate}[leftmargin=2em]
  \item tagged coded objects and minimal predicates for source states, layers, operators, readouts, branches, shell/refinement objects, admissibility witnesses, and tensor-coded composites;
  \item admissible coded transport together with coded closure-class preservation and closure-compatible coded readout;
  \item closure-admissible branch codes;
  \item shell-admissible refinement codes over closure-admissible branches;
  \item normalized admissible routed support over the admissible branch-shell support class.
\end{enumerate}
At this public-shape stage, tensor admissibility, operator-family/lambda-like maps, and support-selection or degeneracy refinements are not part of the compact release claim, even though they remain internally available below.
\end{definition}

\begin{theorem}[Compact public-shape preservation chain]\label{thm:compact-public-shape-preservation}
Let $s_{\mathrm{in}}$ be a coded source state, let $b$ be a closure-admissible branch code, let $\sigma$ be a shell-admissible refinement code on $b$, let $u$ be a normalized admissible routed support code over the admissible branch-shell support class, and let $r$ be a closure-compatible coded readout. Then the v3.16.0 public-shape layer preserves the same closure-relevant readout content along the full compressed chain
\[
s_{\mathrm{in}}
\longrightarrow
(b,\sigma)
\longrightarrow
u
\longrightarrow
r,
\]
provided the coded return realizations are taken inside the admissible branch-shell support sector.
\end{theorem}

\begin{proof}
At OP level, the same compression is built in stages: invariant source-return closure for the OP~1--OP~5 pair is fixed by Theorem~\ref{thm:op15-closure}; branchwise preservation is then isolated by Proposition~\ref{prop:op-branchwise-refinement}; shell-sensitive preservation at the OP~3 interface is isolated by Theorem~\ref{thm:op3-shell}; and normalized routed preservation is isolated by Proposition~\ref{prop:op-admissible-routing}. The QAM translation theorem (Theorem~\ref{thm:qam-translation}) carries exactly that OP chain into coded form. Later in the QAM section this becomes the admissible branch-shell theorem (Theorem~\ref{thm:qam-branch-shell}), the routed admissible support theorem (Theorem~\ref{thm:qam-routed-support}), and the visible/structural separation principle (Corollary~\ref{cor:qam-visible-structural}). The present statement records that entire OP-to-QAM compression as one front-facing release theorem.
\end{proof}

\begin{remark}[Public-shape reading guide]\label{rem:public-shape-reading-guide}
For the purposes of the current release, the later QAM section should be read with the following dictionary in mind:
\begin{enumerate}[leftmargin=2em]
  \item ``coded source state + closure-compatible coded readout'' is the coded front-face of the OP~1/OP~5 source-return pair;
  \item ``closure-admissible branch code'' is the coded branchwise refinement role;
  \item ``shell-admissible refinement code'' is the coded OP~3 shell role;
  \item ``normalized admissible routed support'' is the coded routed aggregation role over the admissible branch-shell support sector;
  \item ``closure-relevant readout content'' always means readout agreement after passage through the associated closure-factor map.
\end{enumerate}
This dictionary is the intended reading bridge between the compact front theorem and the later theorem-level unpacking in Theorem~\ref{thm:qam-branch-shell}, Theorem~\ref{thm:qam-routed-support}, Theorem~\ref{thm:qam-translation}, and Corollary~\ref{cor:qam-visible-structural}.
\end{remark}

\begin{corollary}[Release-level visible/structural rule]\label{cor:release-visible-structural-rule}
Within the compact public-shape QAM package, visible coded mismatch in packet, phase-lock, or $\varepsilon$-screening data is not by itself a closure failure. At the inherited public-shape level, genuine closure failure begins only when branch admissibility or shell admissibility is lost.
\end{corollary}

\begin{remark}[Structural provenance of the public-shape extraction]
The v3.16.0 front block should be read as the coded shadow of one already stabilized OP chain and not as a freestanding new criterion. In that reading, coded source state and coded readout come from the closed OP~1--OP~5 source-return pair of Theorem~\ref{thm:op15-closure}; closure-admissible branch codes come from the branchwise refinement layer of Proposition~\ref{prop:op-branchwise-refinement}; shell-admissible refinement codes come from the OP~3 shell theorem of Theorem~\ref{thm:op3-shell}; normalized admissible routed support comes from the routed aggregation layer of Proposition~\ref{prop:op-admissible-routing}; and the release-level visible/structural rule is the QAM-over-ZF front-face of the OP screening logic, where packet/phase/epsilon data remain preselection data unless they witness loss of admissibility. The public-shape extraction therefore does not add a second closure notion. It republishes the already existing source-return closure architecture in the smaller coded alphabet needed for the current QAM-over-ZF line.
\end{remark}

\begin{remark}
This compact block is intentionally stronger in readability than in scope. It exposes the third public-shape QAM step as one coherent public claim, but it does not yet export the full internal support-selection or wider compression machinery. In other words, v3.17.0 was prepared as the controlled public extraction of the tensor/operator layer, not as the publication of the entire internal QAM architecture.
\end{remark}

\begin{remark}[Deferred kernel-first outlook]
The front theorem should also be read as the visible place where any later zero-kernel migration would have to pay off. The deferred kernel notes below do not introduce a second active tuple language, but they already identify the only later proof-compression route that would matter on the present public surface: first kernel-first transport reduction (Proposition~\ref{prop:qam-kernel-first-transport}), then local branch-shell kernel reduction (Proposition~\ref{prop:qam-kernel-first-branch-shell}), and only afterwards routed family compression (Proposition~\ref{prop:qam-kernel-first-routed}). Read together with the roadmap-fit remark, the later checkpoint summary, the activation criterion of Theorem~\ref{thm:qam-zero-kernel-activation-criterion}, and the later remarks on the staged activation roadmap, the criterion-to-roadmap reading rule, and the compressed synthesis of the deferred zero-kernel path, this deferred path should therefore be understood as a later proof-theoretic refinement of the same front chain rather than as a competing current-release line.
\end{remark}

\begin{remark}[QAM-over-ZF entry layer]
A next formal development line of the Companion is the introduction of QAM not as a replacement for Zermelo--Fraenkel set theory, but as a structure-first coded language over ZF. The point of this layer is to give the current OP/branch/shell architecture a cleaner formal interface for states, layers, operators, readouts, admissibility, closure preservation, and tensor-coded composition while remaining, in principle, back-translatable into set-theoretic form.

At the present integrated stage, the admissible OP fragment remains the first translation target of this QAM line: the OP~1/OP~5 source-return closure architecture, admissible branch transport, OP~3 shell-admissible refinement, admissible routed support, tensor composition over admissible coded objects, and the visible packet/phase-lock/epsilon screening layer. Relative to the public v3.16.0 release, the next public-shape export in that historical sequence was aimed at the tensor and operator/lambda layer rather than at a further broadening of the branch/shell/routed block.

The reserved minor-release path remains the following. First, v3.15.0 fixed the minimal QAM language layer over ZF, including tagged set-codes for states, layers, operators, readouts, branches, shell/refinement objects, admissibility witnesses, and tensor-coded composites. Second, v3.16.0 extracted the branch/shell/routed translation into public-shape form, including closure-admissible branch codes, shell-admissible refinement codes, admissible routed support, and the first coded branch-shell transport results. Third, v3.17.0 realized the tensor/operator public line; fourth, v3.18.0 exported the first public OP~3/QAM closure package; fifth, v3.19.0 hardened the witness-grade reopening discipline of that bundled channel; and the present v\docversion\ release line now pilots a closure-liquid-crystal / holographic Julia-boundary re-reading above that compressed surface while support-selection refinements remain below as later-stage technical support.

The present integrated continuation therefore exposes the compressed public-shape QAM core together with its branch/shell/routed baseline, its tensor/operator extension, its public OP~3/QAM closure package, its bundled witness-grade compression surface, and the released v3.20.0 crystal--shadow architecture: tagged coded objects and minimal predicates, admissible transport with coded closure-class preservation, closure-compatible coded readout, closure-admissible branch codes, shell-admissible refinement codes, admissible routed support, tensor admissibility, tensor closure, first operator-family or lambda-like translation statements, terminal proof-use compression, addressed crystal packets/corridors/classes, theorem-level shadows, and the terminal crystal--shadow interface. The broader internal QAM package remains available below as technical support and does not alter the mathematical claims already frozen in the public v3.20.0 release.

\textbf{QAM-coded object preview.} The minimal intended object types are states, layers, operators, readouts, branches, shell/refinement objects, admissibility/closure witnesses, and tensor-coded composites. A set $x$ is intended to count as a QAM-coded object exactly when it has one of the tagged forms
\[
\begin{aligned}
&\langle 0,a\rangle,\quad
\langle 1,\ell,D\rangle,\quad
\langle 2,\ell_1,\ell_2,G\rangle,\quad
\langle 3,\ell,W,R\rangle,\\
&\langle 4,b,B\rangle,\quad
\langle 5,r,s,T\rangle,\quad
\langle 6,u\rangle,\quad
\langle 7,p,q\rangle.
\end{aligned}
\]
If $p$ and $q$ are QAM-coded objects, the tensor-coded composite is written as
\[
p \otimes q := \langle 7,p,q\rangle.
\]
At the entry stage, these tags and tensor codes are only coding devices that distinguish structural roles inside ZF and reserve tensor syntax for later internal hardening.
\end{remark}

\begin{remark}[Dependency guide for the internal QAM package]
The internal QAM-over-ZF line should be read in the following order: coded objects and predicates first, admissible transport and closure-class schemata second, branch/shell/routed admissibility third, tensor admissibility and tensor closure rules fourth, faithful translation of the admissible OP fragment fifth, operator-family/lambda-like structural maps sixth, and support-selection/degeneracy handling only as a seventh-layer refinement on top of the already admissible fragment.
\end{remark}

\subsection*{How to read the internal QAM package}
The internal QAM package should be read in two layers. At the compact reader level, one should keep only the following chain in view: coded objects and predicates, admissible transport and closure, branch/shell/routed admissibility, tensor closure rules, faithful translation of the admissible OP fragment, and only afterwards operator/lambda and support-selection refinements. At the technical level, the detailed sections below provide the full internal unfolding of these same layers in theorematic form.

\begin{remark}[QAM normalization guide]
The current internal package is not meant to be enlarged indiscriminately. Its next intended use is normalization and later public extraction. In particular, later work should prefer compression, dependency control, and clean appendix-style grouping over uncontrolled addition of parallel QAM subsystems.
\end{remark}

\subsection*{Roadmap and extraction criteria for the v3.17.0 public-shape preparation}
The internal QAM package is now far enough along that one may formulate a controlled extraction path toward the next public QAM-oriented line. The intended purpose of this extraction is not to publish the entire internal QAM package at once, but to expose its stable and readable core in stages.

The current preferred extraction sequence is the following.
\begin{enumerate}[leftmargin=2em]
  \item \textbf{First public QAM line (realized in v3.15.0): language and closure interface.} Publish a compact QAM-over-ZF entry layer together with the minimal predicate package, admissible transport, coded closure-class preservation, closure-compatible readout, and a first translation statement for the admissible OP fragment.
  \item \textbf{Second public QAM line (realized in v3.16.0):} \textbf{branch/shell/routed translation.} Extend the public QAM line by adding closure-admissible branch codes, shell-admissible refinement codes, admissible routed support, and the first coded branch-shell transport results.
  \item \textbf{Third public QAM line (realized in v3.17.0): tensor and operator/lambda extension.} Add tensor admissibility and tensor closure rules together with the first operator-family or lambda-like structural maps and explicit translation or round-trip statements, while keeping selection/support handling internal for now.
  \item \textbf{Fourth public QAM line (realized in v3.18.0): first public OP~3/QAM closure package.} Expose shell certificates, shell-defect localization, kernel persistence under the bridge hypotheses, shell-stable readout factorization, and the bundled OP~3$\to$OP~5/global interface in compact public shape.
  \item \textbf{Fifth public QAM line (realized in v3.19.0): witness-grade hardening and local reopening discipline.} Add witness-grade proof-use control for the bundled channel together with a first narrow licensed kernel-first rewrite window, a first corridor-level activation lemma, a first corridor-to-bundle propagation rule, a first bundle-level reopening compression principle, a first profile-to-run compression principle, a first word-determined run-class step, a first run-class concatenation principle, a first finite run-class normal-form principle, a first normal-form-determined bundled proof-use class step, a first proof-use-class composition principle, a first associative proof-use-class normal-form principle, a first contraction-stable class-normal-form principle, and a terminal proof-use surface.
  \item \textbf{Sixth public QAM line (realized in v3.20.0): closure-liquid-crystal, addressed role layer, theorem-level shadows, and terminal crystal--shadow interface.} Re-read OP~3 as local cell compatibility, OP~5 as defect transport on the same cellular carrier, visible sectors as liquid-crystal domains of one common bulk order, Julia-type diagnostics as holographic boundary readout, stabilize a first base36/Q-word role addressing with reserved 0-band transport support, and connect the resulting terminal crystal-readout surface to a bundled theorem-level shadow package and a terminal crystal--shadow interface.
\end{enumerate}

At the present internal stage, the extraction criteria for the sixth public QAM line should remain strict. A QAM block should be promoted into a public release only if it satisfies all of the following conditions:
\begin{enumerate}[leftmargin=2em]
  \item \textbf{compression:} the public block can be read in a compact main-text form without forcing the full internal package into the public line;
  \item \textbf{dependency clarity:} the block has a clean dependency guide and does not rely on hidden internal lemmas;
  \item \textbf{ZF-groundedness:} the coded language remains explicitly readable as a definitional overlay over ZF rather than as an unexplained new ontology;
  \item \textbf{translation contact:} the public block makes clear what part of the current OP/branch/shell architecture it actually translates;
  \item \textbf{non-overreach:} the public block does not yet claim more than the present coded layer genuinely supports.
\end{enumerate}

\begin{remark}
The current internal package already satisfies much of this extraction logic, but it is still broader than what should be exported immediately. The next internal passes should therefore favor tensor/operator compression, translation clarity, and clean staging over further rapid enlargement.
\end{remark}

\begin{remark}[How the present release-pass checkpoint fits the public-shape roadmap]\label{rem:qam-roadmap-checkpoint-fit}
The deferred kernel-first checkpoint developed later in this file should be read as an internal support block for the same v3.17.0 public-shape preparation path, not as a second release claim. Its current function is to prepare a later proof-compression interface for the compact tensor/operator chain by isolating the future kernel vocabulary, bridge data, staged stress tests, and activation criteria needed for a controlled migration. In particular, the checkpoint helps explain why the third public QAM line should remain narrow at visible level while still allowing later internal compression work to accumulate behind it.
\end{remark}

\subsection*{Compressed public-shape QAM branch-shell core}
The second public QAM line should expose only the smallest next coded layer needed to carry the current admissible OP fragment and its branch/shell/routed refinement into a ZF-grounded formal interface. At compressed level, this public-shape branch-shell core should be read through six ingredients:
\begin{enumerate}[leftmargin=2em]
  \item tagged QAM-coded objects and the minimal predicate layer;
  \item admissible transport together with coded closure-class preservation;
  \item closure-compatible coded readout;
  \item closure-admissible branch codes;
  \item shell-admissible refinement codes;
  \item admissible routed support together with first coded branch-shell transport statements.
\end{enumerate}

\begin{theorem}[Compressed public-shape QAM branch-shell theorem]
Assume that coded QAM objects are read as a definitional overlay over ZF, that admissible coded transport preserves coded closure class, and that the relevant coded readout is closure-compatible. Assume further that closure-admissible branch codes and shell-admissible refinement codes preserve coded closure class, that normalized admissible routed supports use only such branch-shell pairs, and that the admissible OP fragment is translated faithfully into the coded QAM layer. Then the translated source/return closure content of the admissible OP fragment, together with its admissible branchwise, shellwise, and routed realizations, is preserved at coded level, and closure-relevant readout agreement is retained after passage through the closure-factor map.
\end{theorem}


\begin{remark}
This compressed theorem is the intended mathematical core of the second public QAM line. It remains deliberately weaker than the full internal QAM package: tensor rules, operator/lambda maps, and support-selection layers remain available internally, but they need not all be exported at once in order to make the next public QAM step mathematically meaningful.
\end{remark}

\subsection*{Structural systems view: QAM as a virtual machine with a ZFS-style theorem store}
\begin{quote}\small
\noindent\textbf{Status.} \emph{Active explanatory overlay, not a second mathematical criterion.}\par
\noindent\textbf{Block function.} This block explains why the QAM objects, kernel states, snapshots, and safe side branches belong to one architecture, without changing the theorematic content.
\end{quote}
The present QAM export should also be read through a systems-level analogy. The visible coded tuple language plays the role of a readable virtual-machine interface over the ZF-grounded base, while the later exposed payload kernels and kernel normal forms behave like the machine-state identities that remain stable beneath several visible tuple surfaces. In that sense, the public QAM line is the user-facing instruction and state surface, whereas the deeper kernel-first line is the later OS-level state-management discipline. This reading is not introduced as a replacement for the mathematics, but as a compact way to explain why branch/shell/routed transport, kernel normal forms, search keys, snapshots, and safe side branches fit together as one contiguous architecture. It also aligns with the earlier Zuse Z3 comparison already used in the project: the current point is to interpret the coded QAM layer not merely as notation, but as a structured machine view with auditable state persistence.

\begin{remark}[VM/OS reading of the current QAM line]\label{rem:qam-vm-os-reading}
For the present document, the systems analogy should be read in five layers.
\begin{enumerate}[leftmargin=2em]
  \item \textbf{Visible QAM views.} The active tuple surfaces such as $\langle 0,a\rangle$, $\langle 0,a,b\rangle$, branch codes, shell/refinement codes, and routed return codes form the readable public interface of the machine.
  \item \textbf{Kernel / VM state.} The exposed payload kernels, their normal forms, and the resulting class-level closure/readout data form the later machine-state identity beneath that visible interface.
  \item \textbf{OS-level orchestration.} Admissible transport, local branch/shell return, and routed family aggregation act like the current system-level transition and scheduling layer that moves visible states while preserving closure-relevant content.
  \item \textbf{ZFS-style theorem store.} Search keys, snapshots, clones, audits, and later proof-compression branches belong to an internal persistence layer that manages state identity, safe side branches, and deduplication of repeated kernel classes.
  \item \textbf{Pool/group structure.} Sphere objects may later be treated as logical devices, while trioles or routed-support families may be grouped as internal device families. At the present stage, mirror-style reuse of equal kernel classes is the clearest fit; parity-like reconstruction remains exploratory and should be treated only as a later structural option.
\end{enumerate}
\end{remark}

\begin{figure}[H]
\centering
\resizebox{0.94\linewidth}{!}{%
\begin{tikzpicture}[
  font=\small,
  mainbox/.style={draw, rounded corners=3pt, align=center, minimum width=8.0cm, minimum height=1.05cm, inner sep=6pt},
  sidebox/.style={draw, rounded corners=3pt, align=center, minimum width=3.9cm, minimum height=1.0cm, inner sep=5pt},
  arr/.style={-{Latex[length=2.6mm]}, line width=0.5pt}
]

\node[mainbox, fill=blue!6]   (visible)  at (0,0)   {\textbf{Visible QAM views}\\Tuple surfaces, branch/shell codes, routed returns};
\node[sidebox, fill=yellow!10] (public)   at (7.0,0)   {\textbf{Public surface}\\Readable theorem interface};
\node[mainbox, fill=green!6]  (kernel)   at (0,-1.9) {\textbf{Kernel / VM state}\\Payload kernels, normal forms, closure/readout identity};
\node[mainbox, fill=orange!8] (os)       at (0,-3.8) {\textbf{OS-level orchestration}\\Admissible transport, branch/shell return, routed aggregation};
\node[mainbox, fill=purple!7] (store)    at (0,-5.7) {\textbf{ZFS-style theorem store}\\Search keys, snapshots, clones, audits, safe side branches};
\node[sidebox, fill=yellow!10] (internal) at (7.0,-5.7) {\textbf{Internal persistence}\\State reuse and branch hygiene};
\node[mainbox, fill=gray!12]  (pool)     at (0,-7.6) {\textbf{Pool / group structure}\\Sphere objects as logical devices; trioles/routed families as grouped device families};

\draw[arr] (visible) -- (kernel);
\draw[arr] (kernel) -- (os);
\draw[arr] (os) -- (store);
\draw[arr] (store) -- (pool);
\draw[arr] (visible.east) -- (public.west);
\draw[arr] (store.east) -- (internal.west);

\node[draw, rounded corners=5pt, fit=(visible)(kernel)(os)(store)(pool), inner sep=9pt, line width=0.6pt] (stackfit) {};
\node[above=0.15cm of stackfit.north, fill=white, inner sep=2pt] {\textbf{QAM as VM/OS/ZFS stack}};
\end{tikzpicture}%
}\caption{Graphical reading of the current QAM export as a virtual-machine/OS stack with a ZFS-style structural store. Visible tuple surfaces provide the public theorem interface; exposed payload kernels and kernel normal forms provide the later machine-state identity; snapshots, clones, search keys, and audits manage the internal proof store without changing the visible public QAM syntax.}
\label{fig:qam-vm-os-zfs}
\end{figure}

\begin{remark}[Why the ZFS analogy is structurally useful here]\label{rem:qam-zfs-utility}
The ZFS-style reading is not meant as decorative metaphor. It isolates four practical gains already visible in the present file. First, equal kernel classes can later be treated as mirror-like state reuse rather than as repeatedly independent visible tuple surfaces. Second, the deferred kernel-first path already behaves like copy-on-write: experimental compression statements are developed beside the active tuple proofs without rewriting the active line. Third, checkpoints, audits, and activation criteria behave like snapshots and scrub rules for the theorem store, since they record exactly which deferred branches are stable and which still remain non-authoritative. Fourth, once kernel normal forms and search keys are in place, repeated branch/shell/routed return families can be indexed and deduplicated as one structural graph beneath a readable tree-shaped theorem presentation. This is precisely why the systems analogy deserves a more prominent place next to the public QAM extraction itself.
\end{remark}

\begin{remark}[Current scope of the VM/OS/ZFS view]\label{rem:qam-vm-os-scope}
In the inherited v3.16.0 public-shape layer this systems view remains explanatory and organizational. It does not introduce a second formal semantics, and it does not yet activate RAID-like reconstruction as an active theorem rule. The present mathematically justified gains at that inherited layer are the mirror-like reuse of equal kernel classes, the snapshot/clone logic of active versus deferred proof branches, and the content-addressed search-key discipline attached to kernel normal forms. The v3.34.0 extension below activates only a guarded store-and-array vocabulary over the already available QAM/MML/MMA interface. It still does not turn mirror, stripe, parity, or repair language into an OP-closure certificate.
\end{remark}


\subsection*{QAM/MML/MMA store layer and Sphaera array extension}
\begin{quote}\small
\noindent\textbf{Status.} \emph{Active v3.34.0 infrastructure block; no OP closure claim.}\par
\noindent\textbf{Block function.} This block internalizes the earlier ZFS/RAID analogy. QAM supplies the object store, MML supplies the store and array instructions, MMA supplies the execution/readout semantics, and finite Sphaera arrays supply guarded multi-object redundancy, comparison, and residual/parity witnesses.
\end{quote}

The preceding systems view should no longer be read merely as an external metaphor. At the present local layer it becomes a controlled internal interface. The point is not to claim that the Companion has already become a universal computer, nor that OP~1/4 or OP~5(B) have closed. The point is narrower and useful: the same QAM/MML/MMA machinery that already records objects, instructions, and route-faithful execution can also write the persistence layer in which several Sphaera objects are stored, compared, snapshotted, scrubbed, and carried as a guarded array.

\paragraph{Human reading.}
A single Sphaera object is now treated like one protected theorem-state device. Several Sphaera objects may be grouped into an array. In mirror mode, different objects carry the same structural content through different admissible views. In stripe mode, one larger proof or diagnostic object is distributed over several component stores. In parity/residual mode, an additional Sphaera object carries the difference, obstruction, or orthogonal control witness attached to the other components.

\paragraph{LLM/audit reading.}
The layer below gives stable machine-readable handles for store state, snapshot, clone, rollback, scrub, array creation, mirror, stripe, parity, and array snapshot. These handles are not theorematic shortcuts. They are audit contracts: each command must preserve address discipline, guard discipline, lineage discipline, and the non-certificate status unless a separate explicit theorem later discharges that status.

\begin{figure}[H]
\centering
\begin{tikzpicture}[
  node distance=8mm and 10mm,
  box/.style={draw, rounded corners, align=center, inner sep=5pt, text width=0.2\textwidth},
  bigbox/.style={draw, rounded corners, align=center, inner sep=6pt, text width=0.28\textwidth},
  >=Latex
]
\node[box] (qam) {\textbf{QAM}\\Object store\\addresses, objects, guards, lineage};
\node[box, right=of qam] (mml) {\textbf{MML}\\Instruction language\\store and array commands};
\node[box, right=of mml] (mma) {\textbf{MMA}\\Execution semantics\\guard-faithful transitions};
\node[bigbox, below=12mm of mml] (array) {\textbf{Sphaera array layer}\\mirror \quad stripe \quad parity/residual\\snapshot \quad clone \quad rollback \quad scrub};
\draw[->] (qam) -- (mml);
\draw[->] (mml) -- (mma);
\draw[->] (qam.south) -- (array.north west);
\draw[->] (mml.south) -- (array.north);
\draw[->] (mma.south) -- (array.north east);
\end{tikzpicture}
\caption{Reader-facing systems map for the QAM/MML/MMA store-array extension. QAM supplies the stored object layer, MML supplies the allowed store and array commands, MMA supplies the guarded execution semantics, and the Sphaera array layer packages mirror, stripe, and parity/residual modes together with snapshot, clone, rollback, and scrub operations.}
\label{fig:qam-mml-mma-store-array-map}
\end{figure}

\begin{table}[H]
\centering
\small
\begin{tabular}{>{\raggedright\arraybackslash}p{0.18\textwidth} >{\raggedright\arraybackslash}p{0.26\textwidth} >{\raggedright\arraybackslash}p{0.26\textwidth} >{\raggedright\arraybackslash}p{0.20\textwidth}}
\toprule
\textbf{Layer} & \textbf{Human reading} & \textbf{LLM/audit reading} & \textbf{Typical outputs}\\
\midrule
QAM & protected object store & typed address/object/guard/ledger state & store state, snapshot anchor\\
MML & allowed store and array actions & finite instruction alphabet & \texttt{STORE}, \texttt{SNAPSHOT}, \texttt{PARITY}\\
MMA & how commands really act & deterministic guarded transition semantics & continue / halt / reject / guard\\
Sphaera array & several Sphaera objects coupled as one working array & finite guarded family with cross-sphere guards and parity witnesses & mirror, stripe, parity/residual, scrub\\
\bottomrule
\end{tabular}
\caption{Compact overview table for the store-array layer. This separates reader intuition from the audit contract without changing the theorematic status of the block.}
\label{tab:qam-mml-mma-store-array-overview}
\end{table}


\begin{definition}[QAM store state]\label{def:qam-store-state}
A \emph{QAM store state} is a finite guarded tuple
\[
  \mathcal S_n
  :=
  (\mathcal A_n,\mathcal O_n,\mathcal R_n,\mathcal G_n,\mathcal L_n),
\]
where \(\mathcal A_n\) is the finite address/tag domain, \(\mathcal O_n\) is the finite family of stored QAM objects, \(\mathcal R_n\) is the finite relation/link register between stored objects, \(\mathcal G_n\) is the finite guard register, and \(\mathcal L_n\) is the finite lineage/audit ledger. The state is \emph{admissible} if every ledger pointer lands in \(\mathcal A_n\), every relation in \(\mathcal R_n\) uses stored addresses, and every active guard in \(\mathcal G_n\) is satisfied by the displayed objects and relations.
\end{definition}

\begin{definition}[Store instruction alphabet]\label{def:mml-store-instructions}
The local MML store alphabet is the finite instruction family
\[
\begin{split}
  \mathcal I_{\mathrm{store}}
  :=\{&\texttt{STORE},\texttt{LINK},\texttt{SNAPSHOT},\texttt{CLONE},\texttt{ROLLBACK},\texttt{SCRUB},\texttt{MERGE\_CAND},\texttt{EXPORT}\}.
\end{split}
\]
Here \(\texttt{STORE}\) writes an object to an address, \(\texttt{LINK}\) writes a relation, \(\texttt{SNAPSHOT}\) freezes an admissible state, \(\texttt{CLONE}\) opens a guarded side branch from a snapshot, \(\texttt{ROLLBACK}\) returns to a recorded snapshot, \(\texttt{SCRUB}\) checks the guards and lineage, \(\texttt{MERGE\_CAND}\) proposes a guarded merge candidate, and \(\texttt{EXPORT}\) exposes a readout surface. None of these instructions is allowed to manufacture a theorematic certificate by syntax alone.
\end{definition}

\begin{definition}[MMA store transition semantics]\label{def:mma-store-transition}
The local MMA store semantics is a partial deterministic transition map
\[
  \delta_{\mathrm{store}}
  :
  \mathcal S_n\times\mathcal I_{\mathrm{store}}
  \dashrightarrow
  \mathcal S_{n+1}\times
  \{\mathsf{continue},\mathsf{halt},\mathsf{reject},\mathsf{guard}\}.
\]
It is \emph{guard-faithful} if every defined transition either preserves admissibility of the resulting store or returns the status \(\mathsf{guard}\) together with a defect entry in the ledger. It is \emph{lineage-faithful} if each defined transition records enough ledger data to identify its source state, instruction, guard status, and output state.
\end{definition}

\begin{definition}[Guarded snapshot and clone]\label{def:qam-store-snapshot}
For an admissible QAM store \(\mathcal S_n\), a \emph{guarded snapshot} is a tuple
\[
  \operatorname{snap}(\mathcal S_n)
  :=
  (\mathcal S_n,\chi_n,\lambda_n),
\]
where \(\chi_n\) is a finite integrity witness for the address, relation, guard, and ledger registers, and \(\lambda_n\) is a lineage anchor. A \emph{clone} of \(\operatorname{snap}(\mathcal S_n)\) is a new working state \(\mathcal S'_n\) with the same stored content and a new branch tag in the ledger. The clone is not a new theorem; it is a new guarded working branch from the same stored state.
\end{definition}

\begin{proposition}[Snapshot stability of admissible QAM stores]\label{prop:qam-store-snapshot-stability}
Every admissible finite QAM store state admits a guarded snapshot in the sense of Definition~\ref{def:qam-store-snapshot}. If the MMA store semantics is guard-faithful and lineage-faithful, then \(\texttt{SNAPSHOT}\), \(\texttt{CLONE}\), and \(\texttt{ROLLBACK}\) preserve the distinction between stored mathematical objects and release/audit metadata.
\end{proposition}

\begin{proof}
Since the address, object, relation, guard, and ledger registers are finite, their displayed state can be recorded as a finite tuple together with a finite integrity witness and a lineage anchor. This gives \(\operatorname{snap}(\mathcal S_n)\). Guard-faithfulness prevents a failed guard from being silently stored as an admissible state, and lineage-faithfulness records the source and branch data. Thus snapshot, clone, and rollback may move between recorded store states, but they do not promote a release label, branch label, or audit label into the mathematical object language.
\end{proof}

\begin{definition}[Sphaera array]\label{def:sphaera-array}
A \emph{Sphaera array} is a finite guarded family
\[
  \mathbb S_n
  :=
  \bigl((\mathcal S_n^{(1)},\ldots,\mathcal S_n^{(k)});\Gamma_n,\Pi_n,\Lambda_n\bigr),
  \qquad k\ge 1,
\]
where each \(\mathcal S_n^{(i)}\) is an admissible QAM store state, \(\Gamma_n\) is the cross-sphere guard register, \(\Pi_n\) is the parity/residual witness register, and \(\Lambda_n\) is the array-level lineage ledger. The array is admissible if all component stores are admissible and all cross-sphere guards in \(\Gamma_n\) are satisfied.
\end{definition}

\begin{definition}[Array instruction alphabet]\label{def:sphaera-array-instructions}
The local MML array alphabet is
\[
\begin{aligned}
  \mathcal I_{\mathrm{array}}
  :=\{&\texttt{ARRAY\_CREATE},\texttt{MIRROR},\texttt{STRIPE},\texttt{PARITY},\\
       &\texttt{SCRUB\_ARRAY},\texttt{REPAIR\_CAND},\texttt{SNAPSHOT\_ARRAY}\}.
\end{aligned}
\]
The command \(\texttt{REPAIR\_CAND}\) intentionally returns a candidate repair surface, not an automatic mathematical repair theorem. The command \(\texttt{PARITY}\) records a residual or orthogonal control witness, not a proof that the residual vanishes.
\end{definition}

\begin{definition}[Mirror, stripe, and parity modes]\label{def:sphaera-array-modes}
Let \(\mathbb S_n\) be an admissible Sphaera array.
\begin{enumerate}[leftmargin=2em]
  \item In \emph{mirror mode}, two component stores \(\mathcal S_n^{(i)}\) and \(\mathcal S_n^{(j)}\) are coupled by a cross-guard requiring equality of the relevant closure-compatible readout, even if the visible store presentations differ.
  \item In \emph{stripe mode}, an object \(X\) is stored as a finite guarded decomposition \(X\rightsquigarrow(X_1,\ldots,X_k)\), with each stripe \(X_i\) stored in a specified component and with a reconstruction witness recorded in \(\Pi_n\).
  \item In \emph{parity/residual mode}, one component or register stores a residual witness \(P\) for a guarded comparison of other components. The condition \(P=0\), if ever needed, must be proved separately; it is not asserted by the existence of the parity register.
\end{enumerate}
\end{definition}

\begin{definition}[Triadic Sphaera array]\label{def:triadic-sphaera-array}
A \emph{triadic Sphaera array} is an admissible Sphaera array of the form
\[
  \mathbb S_n^{ABC}
  :=
  (\mathcal S_A,\mathcal S_B,\mathcal S_C;\Gamma_{ABC},\Pi_{ABC},\Lambda_{ABC})
\]
in which \(\mathcal S_A\) and \(\mathcal S_B\) carry two active readout or probe-facing stores, while \(\mathcal S_C\) carries the orthogonal residual/parity witness. The guard register records the structural condition that \(\mathcal S_C\) is read orthogonally to the active \(A/B\) coupling channel and that the HC/MM readout of the array preserves the inherited route tag.
\end{definition}

\begin{proposition}[Guarded array snapshot]\label{prop:sphaera-array-snapshot}
Every admissible finite Sphaera array admits a guarded array snapshot
\[
  \operatorname{snap}_{\mathrm{arr}}(\mathbb S_n)
  :=
  \bigl(\mathbb S_n,(\chi_n^{(1)},\ldots,\chi_n^{(k)}),\chi_n^{\Gamma},\chi_n^{\Pi},\lambda_n^{\mathrm{arr}}\bigr),
\]
preserving component stores, cross-sphere guards, parity/residual witnesses, and array lineage.
\end{proposition}

\begin{proof}
Apply Proposition~\ref{prop:qam-store-snapshot-stability} to each finite component store. Since the cross-guard, parity/residual, and array-ledger registers are also finite, record their integrity witnesses and one array lineage anchor. The resulting tuple is a guarded array snapshot. It records the array state without adding an OP-closure certificate and without changing any component object into metadata.
\end{proof}

\begin{proposition}[Finite array scrub]\label{prop:sphaera-array-scrub}
For every finite Sphaera array there is a deterministic scrub procedure
\[
  \operatorname{scrub}_{\mathrm{arr}}(\mathbb S_n)
  \in
  \{\mathsf{pass},\mathsf{fail}\}\times\mathcal D_n,
\]
where \(\mathcal D_n\) is a finite defect register. The scrub passes exactly when all component-store guards, cross-sphere guards, parity/residual registrations, and lineage pointers pass their finite checks.
\end{proposition}

\begin{proof}
The component store registers, cross-guard register, parity/residual register, and lineage ledger are finite by definition. A deterministic scrub checks them in a fixed order and records every failed check in \(\mathcal D_n\). If no failed check is recorded, the result is \(\mathsf{pass}\); otherwise it is \(\mathsf{fail}\) with the finite defect register. No repair or closure claim is produced by this test.
\end{proof}

\begin{theorem}[QAM/MML/MMA Sphaera array layer]\label{thm:qam-mml-mma-sphaera-array-layer}
The QAM/MML/MMA stack admits a conservative Sphaera array layer in the following sense. Given finite admissible QAM stores, the store instruction alphabet of Definition~\ref{def:mml-store-instructions}, the array instruction alphabet of Definition~\ref{def:sphaera-array-instructions}, and guard-faithful lineage-faithful MMA transition semantics, one obtains finite guarded snapshots, finite array snapshots, and deterministic array scrub surfaces. These surfaces preserve stored mathematical objects, cross-sphere guards, residual/parity witnesses, and audit lineage without manufacturing an OP-closure certificate.
\end{theorem}

\begin{proof}
The store part follows from Definitions~\ref{def:qam-store-state}--\ref{def:qam-store-snapshot} and Proposition~\ref{prop:qam-store-snapshot-stability}. The array part follows by grouping finitely many admissible stores as in Definition~\ref{def:sphaera-array}, assigning finite cross-guards, finite parity/residual witnesses, and a finite lineage ledger, and then applying Propositions~\ref{prop:sphaera-array-snapshot} and~\ref{prop:sphaera-array-scrub}. The last clause follows from the explicit non-certificate restrictions in Definitions~\ref{def:mml-store-instructions}, \ref{def:sphaera-array-instructions}, and \ref{def:sphaera-array-modes}: mirror, stripe, parity, scrub, and repair-candidate language records infrastructure and diagnostic data only. It does not assert vanishing residuals or OP closure.
\end{proof}

\begin{corollary}[RAID-like reading without RAID overclaim]\label{cor:sphaera-array-raid-reading}
The Sphaera array layer supports a RAID-like reading: mirror mode corresponds to redundant closure-compatible readout, stripe mode corresponds to guarded distribution of a larger object over several component stores, and parity/residual mode corresponds to an explicit diagnostic witness for cross-component mismatch. This reading is internal to QAM/MML/MMA, but it remains weaker than technical RAID reconstruction and weaker than any OP-closure theorem.
\end{corollary}

\begin{proof}
This is only a translation of Definition~\ref{def:sphaera-array-modes} into systems language. The mathematical content remains the guarded store, snapshot, array, parity/residual, and scrub structure already proved above.
\end{proof}

\begin{remark}[Tool extraction from the Sphaera array layer]\label{rem:sphaera-array-tool-extraction}
The new layer extracts four reusable tools for later Companion revisions.
\begin{enumerate}[leftmargin=2em]
  \item \textbf{Snapshot tool.} Freeze one admissible QAM store without changing the mathematical object language.
  \item \textbf{Clone/rollback tool.} Open or return from a guarded branch while preserving lineage and non-certificate status.
  \item \textbf{Array-scrub tool.} Test several Sphaera stores, cross-guards, and residual witnesses as one finite diagnostic surface.
  \item \textbf{Triadic parity tool.} Use a third, orthogonal Sphaera object as residual/parity witness for a two-store active coupling.
\end{enumerate}
These tools are deliberately infrastructural. Their value is that later OP-facing work can compare, branch, and roll back several coupled Sphaera objects without confusing state persistence with theorematic closure.
\end{remark}

\begin{remark}[Human/LLM interface gain]\label{rem:sphaera-array-human-llm-gain}
For a human reader, the store-array layer explains why the Companion keeps snapshots, historical witnesses, audit ledgers, and branch restrictions so visibly: they are the theorem-store equivalent of protected state persistence. For an LLM reader, the same layer supplies exact handles: \(\mathcal S_n\) for one store, \(\mathbb S_n\) for an array, \(\Gamma_n\) for cross-guards, \(\Pi_n\) for parity/residual witnesses, and \(\Lambda_n\) for lineage. This is the intended compromise between human readability and machine precision.
\end{remark}


\subsection*{Pre-v4 QAM Hilbert-shift and guarded mixed readouts}
\begin{quote}\small
\noindent\textbf{Status.} \emph{Pre-v4 preparation layer only.} The inherited QAM tuple tags are not globally migrated in Version~\docversion{}. The block below records the conservative Hilbert-shift plan, the mixed-readout package discipline, and an optional binary mask registry for future readout classification.
\end{quote}

The index-0 issue is now treated explicitly. In the inherited QAM tuple namespace, tag 0 occurs as an inhabited constructor tag, for instance on visible tuple surfaces of the form $\langle 0,a\rangle$. In a Boolean mask or IPv4-like readout registry, however, value 0 must mean no active kind. The next major QAM migration should therefore not identify these two uses. The preparation layer below records the planned Hilbert-hotel shift that can free tag 0 in the future while preserving the payloads and proofs.

\begin{definition}[QAM Hilbert-shift plan]\label{def:qam-hilbert-shift-plan}
Let $\mathcal K_{\mathrm{old}}$ be the inherited finite QAM tuple-tag namespace used by the current public syntax, and let $\mathcal K_{\mathrm{new}}$ be a shifted namespace with tag 0 reserved for the empty/no-active-kind marker. The planned Hilbert-shift is the constructor-level map
\[
  H_{\mathrm{QAM}}(\langle i,x_1,\ldots,x_m\rangle)
  :=
  \langle i+1,x_1,\ldots,x_m\rangle.
\]
It shifts only the constructor tag. It does not change the stored payload $(x_1,\ldots,x_m)$, the intended constructor distinction, or any mathematical content carried by the payload.
\end{definition}

\begin{definition}[Pre-v4 reserved-zero discipline]\label{def:pre-v4-reserved-zero-discipline}
The preparation line distinguishes two uses of zero.
\begin{enumerate}[leftmargin=2em]
  \item In the inherited v3.x tuple syntax, tag 0 remains an inhabited tuple tag wherever the old QAM definitions still use it.
  \item In the planned v4 shifted namespace and in every binary mask registry, tag or mask value 0 is reserved for empty, undefined, or no active kind.
\end{enumerate}
No Version~\docversion{} statement may silently replace the first convention by the second. Any global replacement is reserved for a later v4 migration theorem.
\end{definition}

\begin{proposition}[Conservative shape of the QAM Hilbert-shift]\label{prop:qam-hilbert-shift-conservative-shape}
Suppose a finite QAM sublanguage is translated by $H_{\mathrm{QAM}}$ and every constructor predicate, admissibility predicate, readout predicate, and closure-relevant predicate is translated by the same tag shift. Then the translated sublanguage is a conservative reindexing of the old one: payloads, constructor distinctions, admissibility status, readout compatibility, and closure-relevant truth values are preserved under the translation.
\end{proposition}

\begin{proof}
The map $H_{\mathrm{QAM}}$ is injective on constructor tags and leaves all payload coordinates unchanged. A constructor predicate in the shifted namespace is defined to hold exactly when the corresponding old predicate holds before applying the shift. The same translation is imposed on admissibility, readout, and closure-relevant predicates. Thus every proof step that used only constructor distinction, payload data, admissibility, readout compatibility, or closure-relevant class data has an isomorphic shifted proof step. The proposition records the required conservativity shape for a later v4 migration; it does not perform the global migration here.
\end{proof}

\begin{definition}[Guarded mixed QAM readout]\label{def:guarded-mixed-qam-readout}
A \emph{guarded mixed QAM readout} is a finite derived package
\[
  R_{\mathrm{mix}}
  :=
  (r_1,\ldots,r_m;\Gamma_{\mathrm{mix}},\mu,J_{\mathrm{mix}},\sigma_{\mathrm{mix}}),
\]
where each $r_i$ is an already admissible component readout on its own allowed QAM kind, $\Gamma_{\mathrm{mix}}$ is a finite guard package, $\mu$ is a finite aggregation map on component readout values, $J_{\mathrm{mix}}$ is the closure-factor map for the aggregate, and $\sigma_{\mathrm{mix}}$ records the theorematic status. It is admissible only when all component readouts are admissible and the aggregation map does not identify non-equivalent closure-relevant classes.
\end{definition}

\begin{definition}[Binary mask registry for mixed readouts]\label{def:binary-mask-registry-mixed-readouts}
A \emph{binary mask registry} for mixed QAM readouts assigns to each inhabited readout kind a nonzero power-of-two bit
\[
  \beta_j = 2^j,
  \qquad j\ge 0,
\]
and reserves mask value 0 for no active kind. A mixed-readout mask is a finite Boolean OR of active kind bits. Boolean AND tests whether a component kind is present, XOR records symmetric difference between two masks, and complement is allowed only relative to a fixed finite mask universe $\Omega$.
\end{definition}

\begin{proposition}[Mixed readout factorization]\label{prop:mixed-readout-factorization}
Let $R_{\mathrm{mix}}$ be a guarded mixed QAM readout whose component readouts $r_i$ are closure-compatible. If the aggregation map $\mu$ factors through the component closure-factor maps and $J_{\mathrm{mix}}$ factors through the resulting aggregate class, then $R_{\mathrm{mix}}$ is closure-compatible.
\end{proposition}

\begin{proof}
Closure-compatible component readouts identify no closure-relevant distinction beyond their component class maps. Since $\mu$ factors through those class maps, it can only aggregate already class-level data. Since $J_{\mathrm{mix}}$ factors through that aggregate class, the mixed readout cannot introduce a new closure-relevant distinction that was invisible to the component closure classes. Therefore the mixed package is closure-compatible. The argument uses only factorization and guards; it does not introduce primitive mixed tensors.
\end{proof}

\begin{remark}[v4 migration trigger]\label{rem:v4-qam-migration-trigger}
The full replacement of the inherited tuple tags by the shifted tags is intentionally not a Version~\docversion{} operation. It changes the visible QAM object-language namespace and therefore belongs to a major version. Version~\docversion{} only records why the shift is useful, what the conservative proof obligation looks like, and how mixed readouts can already be treated as derived packages without forcing the migration.
\end{remark}


% ============================================================================
% Quantum Sphaera Companion v3.35.0-r03 insertion patch
% Two-sided structural tensor notation for QAM/MML/MMA readouts
%
% Intended insertion point:
%   In the active-core source, place this block after the short QAM/MML/MMA
%   reader-orientation / pre-v4 QAM preparation material, and before any
%   broad theorematic QAM continuation.  In older sources the closest safe
%   anchor is after the remark "QAM normalization guide".
%
% Scope:
%   Conservative notation/compression layer only.  This is not the separate
%   pedagogical reader manual paper; that manual remains a future standalone
%   companion/guide document.
% ============================================================================

\subsection*{Two-sided structural tensor notation for QAM/MML/MMA readouts}
\begin{quote}\small
\noindent\textbf{Status.} \emph{Active notation and compression layer; not a new tensor calculus and not a new closure criterion.}\par
\noindent\textbf{Block function.} This block records the narrow two-sided tensor grammar used to separate readout-side, structure-side, comparison-side, and mask-side roles in later QAM/MML/MMA expressions.  It is included here as a conservative technical layer only.  A fuller pedagogical reader manual for the notation is intentionally kept as a separate future paper or companion guide.
\end{quote}

The QAM/MML/MMA line repeatedly alternates between four roles: frame-side readout, structure-side action, comparison of active channels, and compatibility testing of readout or mask modes.  The following notation is introduced only to make those roles visible in the formula syntax.  It does not replace classical tensor calculus and it does not turn every displayed expression into a new primitive tensor object.

\begin{definition}[Two-sided structural tensor interface]\label{def:two-sided-structural-tensor-interface}
A \emph{two-sided structural tensor display} is a tensor-like expression in which a structural object may carry front-side and back-side indices, written schematically as
\[
{}^{A}{}_{B}\mathcal{T}^{i}{}_{j}.
\]
The front indices \({}^{A}{}_{B}\) record the left-sided or readout-facing interface: frame, observer, language, readout, or external addressing data.  The back indices \({}^{i}{}_{j}\) record the right-sided or structure-facing interface: action, transport, field, operator, or internal structural data.  Thus the display is read as a typed structural interface rather than as a mere component array.
\end{definition}

\begin{definition}[Admissible contractions and comparisons]\label{def:two-sided-admissible-contractions}
The two-sided notation obeys the following contraction discipline.
\begin{enumerate}[leftmargin=2em]
  \item \textbf{Classical action/readout.} Ordinary upper--lower contraction remains unchanged.  If
  \[
  A_\mu\,{}^\mu\mathcal{T}_\nu\,B^\nu
  \]
  is well-defined, it is read as a classical action, evaluation, or sandwich readout: first apply the structure-facing operation to \(B\), then read the result by \(A\).
  \item \textbf{Active--active comparison.} Two upper indices are not contracted as a raw Einstein sum.  They may only be compared through an explicitly declared comparison tensor, for instance
  \[
  G_{\alpha\beta}X^\alpha Y^\beta .
  \]
  This means active structural comparison, such as angle, distance, coupling strength, orthogonality, channel similarity, or residual size.
  \item \textbf{Passive--passive comparison.} Two lower indices are likewise not contracted raw.  They may only be compared through an explicitly declared dual comparison tensor, for instance
  \[
  G^{\alpha\beta}\omega_\alpha\eta_\beta .
  \]
  This means comparison of readout forms, observer modes, measurement channels, or passive projections.
  \item \textbf{Mask and selector compatibility.} Boolean, finite, or mask-like compatibility tests must be carried by an explicit selector tensor, for instance
  \[
  \chi_{mn}R^mS^n .
  \]
  This expression says that readout modes \(m\) and \(n\) are first tested for admissible compatibility by \(\chi\), and only compatible contributions are admitted to the mixed readout.
\end{enumerate}
\end{definition}

\begin{remark}[Conservative reading rule]\label{rem:two-sided-conservative-reading-rule}
The two-sided notation is conservative because the only unchanged contraction is the classical upper--lower action/readout contraction.  Same-variance pairings are not silently added to the calculus; they are admitted only when a comparison tensor or selector tensor has been declared.  Hence every two-sided display must be expandable into ordinary QAM-over-ZF role predicates together with explicit comparison or compatibility data.
\end{remark}

\begin{definition}[Guarded mixed QAM readout]\label{def:guarded-two-sided-mixed-qam-readout}
Let \(R^m\) and \(S^n\) be QAM readout modes and let \(\chi_{mn}\) be an explicit admissibility or compatibility selector.  A \emph{guarded mixed QAM readout} is a derived readout expression of the form
\[
\operatorname{Mix}_{\chi}(R,S)(x)
\quad\leadsto\quad
\chi_{mn}R^m(x)S^n(x),
\]
where the right-hand display is read with the contraction discipline of Definition~\ref{def:two-sided-admissible-contractions}.  The mixed readout is therefore not a new primitive QAM object.  It is a compatibility-filtered package of already admitted readout modes.
\end{definition}

\begin{remark}[Residual comparison grammar]\label{ex:two-sided-residual-comparison}
The live residual triple
\[
\rho=(\rho_{\mathrm{CL}},\rho_{\Theta},\rho_A)
\]
may be compressed into a residual-channel vector \(\rho^\alpha\).  Its total comparison size may then be written as
\[
\|\rho\|_G^2=G_{\alpha\beta}\rho^\alpha\rho^\beta .
\]
This is not a raw contraction of two upper indices.  It is the comparison of two active residual-channel entries through the declared comparison tensor \(G_{\alpha\beta}\).
\end{remark}

\begin{remark}[HC/MM sandwich grammar]\label{ex:two-sided-hc-mm-sandwich}
An HC/MM-mediated readout may be displayed schematically as
\[
{}^{R}{}_{\Omega}\mathcal{M}^{i}{}_{j},
\]
where the front side records the readout or observer-side interface and the back side records the structure-side transport or matrix action.  When paired with ordinary upper--lower indices, the expression remains an ordinary sandwich readout.  When two readout channels or two structure channels are compared, the comparison must be mediated by an explicitly named comparison tensor or selector.
\end{remark}

\begin{proposition}[Conservative expansion of two-sided displays]\label{prop:two-sided-conservative-expansion}
Every well-formed two-sided structural tensor display in the sense of Definitions~\ref{def:two-sided-structural-tensor-interface} and~\ref{def:two-sided-admissible-contractions} admits a conservative expansion into the inherited QAM-over-ZF language: front-side indices expand to readout/frame/language role data, back-side indices expand to structure/action/transport role data, classical upper--lower contractions remain ordinary action or evaluation clauses, and every same-variance pairing expands to an explicit comparison or compatibility predicate.
\end{proposition}

\begin{proof}
Erase the visual sidedness notation and retain its role data as ordinary QAM codes.  Front-side symbols become readout, frame, observer, or language tags.  Back-side symbols become structure, action, transport, or field tags.  Classical upper--lower contractions are left unchanged.  A same-variance display is well formed only if it contains a named comparison tensor or selector, and so it expands to the corresponding comparison or compatibility predicate.  No new primitive closure relation, admissibility relation, or theorematic certificate is introduced by this erasure.  The display is therefore a conservative abbreviation of already typed QAM-over-ZF data.
\end{proof}

\begin{remark}[LLM-facing role stability]\label{rem:two-sided-llm-role-stability}
For machine reading and LLM-assisted proof navigation, the notation functions as a syntactic guardrail.  The front/back split marks whether a symbol belongs to the readout-facing side or the structure-facing side, while the contraction discipline marks whether a formula is an action, a readout, a comparison, or a compatibility test.  This reduces the risk that a readout is mistaken for an operator, that a diagnostic comparison is mistaken for a proof of closure, or that a guarded mixed readout is mistaken for a new primitive tensor object.
\end{remark}

\begin{remark}[Use in the present release line]\label{rem:two-sided-present-release-use}
In the present v3.35.x line this notation should be used locally and sparingly: first for guarded mixed QAM readouts, residual/drift comparisons, HC/MM sandwich expressions, triadic orthogonality comparisons, and Boolean or mask-like compatibility tests.  A global rewrite of the whole Companion into the two-sided notation is not part of this release line.  The separate pedagogical reader manual may later explain the notation at a slower pace without enlarging the active theorem core.
\end{remark}


% ============================================================================
% Quantum Sphaera Companion v3.35.0-r03 insertion patch
% Two-sided tensor stability, mixed-readout normal form, and Boolean masks
%
% Scope:
%   Stabilization layer only.  This block makes the r02 notation explicitly
%   expandable back into QAM-over-ZF role data and records the mask/readout
%   grammar that will later help the v4 QAM reindexing.  It is not the separate
%   pedagogical reader-manual paper.
% ============================================================================

\subsection*{r03 stability layer: two-sided expansion, mixed-readout normal form, and Boolean masks}
\begin{quote}\small
\noindent\textbf{Status.} \emph{Conservative stabilization of the r02 notation layer; not a global rewrite, not a new tensor calculus, and not a closure theorem.}\par
\noindent\textbf{Block function.} This block fixes the exact erasure route from two-sided displays back to ordinary QAM-over-ZF role data, normalizes guarded mixed readouts, and states the Boolean selector grammar used by IPv4-like readout masks.  The separate reader manual remains a future standalone paper or guide.
\end{quote}

The r02 notation made readout-side, structure-side, comparison-side, and mask-side roles visible in formulas.  The present r03 step turns that notation into a safer tool by recording the normal forms that an audit reader or LLM reader should recover whenever a two-sided display appears.  The rule is intentionally asymmetric: notation may compress role data, but every compressed display must expand back into explicit role data before it can be used as proof content.

\begin{definition}[Two-sided role-erasure map]\label{def:r03-two-sided-role-erasure-map}
Let \(E\) be a well-formed two-sided structural tensor display in the sense of Definitions~\ref{def:two-sided-structural-tensor-interface} and~\ref{def:two-sided-admissible-contractions}.  Its \emph{role-erasure expansion}
\[
\mathfrak{E}_{\mathrm{role}}(E)
\]
is the ordinary QAM-over-ZF package obtained by replacing:
\begin{enumerate}[leftmargin=2em]
  \item every front-side index by an explicit readout, frame, observer, or language-role tag;
  \item every back-side index by an explicit structure, action, field, transport, or operator-role tag;
  \item every classical upper--lower contraction by the corresponding ordinary action/evaluation clause;
  \item every same-variance comparison by a named comparison predicate carrying its declared comparison tensor;
  \item every Boolean or selector occurrence by a finite compatibility predicate over the declared mask universe.
\end{enumerate}
This expansion is an erasure of notation only.  It does not erase guards, selectors, comparison tensors, branch restrictions, or provenance data.
\end{definition}

\begin{proposition}[Role-erasure normal form]\label{prop:r03-role-erasure-normal-form}
Every well-formed two-sided display \(E\) used in the v3.35.x line has a role-erasure expansion \(\mathfrak{E}_{\mathrm{role}}(E)\) whose clauses are built from already admitted QAM role predicates, admitted comparison data, admitted selector data, and ordinary finite tuple coding over ZF.  In particular, the two-sided display cannot by itself introduce a new primitive tensor kind, a new admissibility relation, or a new theorematic certificate.
\end{proposition}

\begin{proof}
Proceed by structural inspection of the display grammar.  Front/back sidedness contributes only role tags.  Ordinary upper--lower contraction is already part of the inherited action/readout grammar.  A same-variance pairing is well formed only when it contains an explicitly named comparison tensor or selector, hence it expands to a comparison or compatibility predicate rather than to an untyped Einstein contraction.  All finite lists of such clauses are QAM tuple-coded over the existing ZF substrate.  Therefore the two-sided syntax is conservative under role erasure.
\end{proof}

\begin{definition}[Guarded mixed-readout normal form]\label{def:r03-guarded-mixed-readout-normal-form}
A guarded mixed QAM readout is in \emph{r03 normal form} when it is presented as a finite package
\[
\mathcal R_{\mathrm{mix}}
=
\bigl(I,\{R_i\}_{i\in I},\Gamma_{\mathrm{mix}},\chi,\mu,J_{\mathrm{mix}}\bigr),
\]
where \(I\) is a finite index set, each \(R_i\) is an already admitted component readout, \(\Gamma_{\mathrm{mix}}\) is the inherited guard package, \(\chi\) is a finite selector or mask predicate over \(I\), \(\mu\) is an aggregation map defined only on selector-admissible component families, and \(J_{\mathrm{mix}}\) is the declared mixed readout surface.  A two-sided display such as
\[
\chi_{mn}R^mS^n
\]
is only an abbreviation for this finite package in the binary case \(I=\{m,n\}\).
\end{definition}

\begin{proposition}[Mixed-readout normal-form expansion]\label{prop:r03-mixed-readout-normal-form-expansion}
Every guarded mixed readout admitted by Definition~\ref{def:guarded-two-sided-mixed-qam-readout} expands to an r03 guarded mixed-readout normal form in the sense of Definition~\ref{def:r03-guarded-mixed-readout-normal-form}.  Conversely, every r03 guarded mixed-readout normal form may be displayed in two-sided shorthand only after its guards, selector, aggregation map, and output surface have been declared.
\end{proposition}

\begin{proof}
The binary shorthand \(\chi_{mn}R^mS^n\) already names the two component readouts and the selector.  The surrounding QAM context supplies the inherited guard package, the admissible aggregation rule, and the declared output surface.  Collecting these data gives the normal-form tuple.  Conversely, starting from the normal-form tuple, one may suppress the finite package into two-sided shorthand only because the suppressed guard, selector, aggregation, and output data remain recoverable.  Thus the shorthand and the normal form are not different mathematical objects; the former is a compact notation for the latter.
\end{proof}

\begin{definition}[Boolean selector and mask grammar]\label{def:r03-boolean-selector-mask-grammar}
Fix a finite readout-kind universe \(\Omega=\{k_1,\dots,k_N\}\).  A \emph{Boolean selector mask} is a function
\[
\chi:\Omega\longrightarrow\{0,1\}
\]
or, equivalently, a finite bit mask whose value 0 denotes no active kind and whose nonzero bits denote active readout kinds.  The admissible Boolean operations are:
\begin{align*}
\chi_A\vee\chi_B &\quad\text{for union of active readout kinds,}\\
\chi_A\wedge\chi_B &\quad\text{for intersection or compatibility testing,}\\
\chi_A\oplus\chi_B &\quad\text{for diagnostic symmetric difference,}\\
\neg_{\Omega}\chi_A &\quad\text{for relative complement inside the fixed universe }\Omega.
\end{align*}
No absolute complement is used without first declaring \(\Omega\).  A selector expression may filter a readout package, but it does not alter the payload of any component readout.
\end{definition}

\begin{proposition}[Mask-filter conservativity]\label{prop:r03-mask-filter-conservativity}
Let \(\mathcal R_{\mathrm{mix}}\) be an r03 guarded mixed-readout normal form and let \(\chi\) be a Boolean selector mask over its finite readout-kind universe.  Applying \(\chi\) restricts the active component menu of \(\mathcal R_{\mathrm{mix}}\), but it does not change the underlying component payloads, the inherited guards, or the closure status of any component readout.
\end{proposition}

\begin{proof}
A selector mask is a finite predicate on the index universe.  It chooses which already declared component clauses are active in the aggregation domain.  It neither rewrites the component readouts nor changes their guard packages.  Hence any closure-compatible status used by the mixed package still has to come from the component readouts and the declared factorization map, not from the mask itself.
\end{proof}

\begin{remark}[LLM/audit role table]\label{rem:r03-llm-audit-role-table}
For later machine reading, the following table fixes the intended role parsing.  It is a local audit aid inside the Companion, not the separate pedagogical reader manual.
\begin{center}
\renewcommand{\arraystretch}{1.12}
\begin{tabular}{|p{0.28\textwidth}|p{0.34\textwidth}|p{0.25\textwidth}|}
\hline
\textbf{Displayed form} & \textbf{Intended role} & \textbf{Audit guard} \\
\hline
front-side indices & readout, frame, observer, language, or addressing side & expand to explicit QAM role tags \\
\hline
back-side indices & structure, action, field, transport, or operator side & expand to explicit QAM role tags \\
\hline
upper--lower contraction & classical action, evaluation, or sandwich readout & ordinary tensor/QAM readout clause \\
\hline
upper--upper via \(G_{\alpha\beta}\) & active structural comparison & no raw same-variance contraction \\
\hline
lower--lower via \(G^{\alpha\beta}\) & passive readout comparison & no raw same-variance contraction \\
\hline
\(\chi_{mn}\) or a mask bit & compatibility or activation selector & filter only; no new payload \\
\hline
\end{tabular}
\end{center}
\end{remark}

\begin{corollary}[r03 proof-neutrality of the tensor layer]\label{cor:r03-proof-neutrality-of-tensor-layer}
The r03 tensor-stability layer may be cited to shorten repeated readout/action/comparison/mask passages, but it may not be cited as an independent proof of closure, OP resolution, primitive mixed-tensor existence, or global QAM namespace migration.
\end{corollary}

\begin{proof}
By Propositions~\ref{prop:r03-role-erasure-normal-form}, \ref{prop:r03-mixed-readout-normal-form-expansion}, and~\ref{prop:r03-mask-filter-conservativity}, every r03 display reduces to already declared role data, normal-form mixed-readout data, and finite selector predicates.  The reduction supplies notation control only.  It does not add the missing theorematic hypotheses needed for closure, OP resolution, primitive mixed tensors, or v4 reindexing.
\end{proof}

\begin{remark}[Placement of the future reader manual]\label{rem:r03-reader-manual-placement}
The present block deliberately stays short and formal.  A slower reader manual may later explain the same grammar with diagrams, examples, and human-facing intuition, but that manual should remain a separate paper or guide so that the Active Core remains the governing proof line.
\end{remark}


% Quantum Sphaera Companion v3.35.0-r04 release-candidate insertion
% Release-candidate freeze/audit layer only.  No mathematical closure claim.

\subsection*{r04 release-candidate freeze audit for the pre-v4 QAM/tensor layer}
\addcontentsline{toc}{subsection}{r04 release-candidate freeze audit for the pre-v4 QAM/tensor layer}
\label{subsec:r04-release-candidate-freeze-audit-qam-tensor}

\noindent\textbf{Status.} \emph{Release-candidate audit layer only; no new theorematic certificate, no global QAM namespace migration, no primitive mixed tensor, and no OP closure.}\par
\noindent\textbf{Block function.} This r04 block freezes the r02/r03 notation and mixed-readout preparation layer for release-candidate review.  It records what must remain true when the Active Core is read together with the separate Historical Witness PDF.

\begin{definition}[r04 pre-v4 QAM/tensor release-candidate package]
\label{def:r04-pre-v4-qam-tensor-release-candidate-package}
The \emph{r04 pre-v4 QAM/tensor release-candidate package} consists of the following finite data:
\begin{enumerate}[label=(\roman*)]
  \item the inherited v3.34.0 public split-bundle baseline and its concept/version DOI record;
  \item the v3.35.0 short-abstract and extended-orientation architecture;
  \item the pre-v4 Hilbert-shift plan reserving global QAM tag migration for a later major version;
  \item the guarded mixed-readout discipline and optional binary mask registry;
  \item the r02 two-sided structural tensor notation;
  \item the r03 role-erasure, mixed-readout normal form, Boolean selector grammar, and proof-neutrality layer;
  \item the present r04 audit/handoff metadata saying that the Active Core is the governing proof line and that the Historical Witness PDF remains preserved proof memory.
\end{enumerate}
It is a packaging and audit object, not a closure object.
\end{definition}

\begin{proposition}[r04 release-candidate conservativity]
\label{prop:r04-release-candidate-conservativity}
The r04 pre-v4 QAM/tensor release-candidate package is conservative over the r03 tensor-stability layer with respect to theorematic content.  It may change release-candidate metadata, routing language, audit surfaces, and handoff wording, but it may not create any new OP certificate, global Hilbert-shift theorem, primitive mixed-tensor axiom, or replacement of the inherited QAM tuple namespace.
\end{proposition}

\begin{proof}
Each mathematical role introduced by r02 is erased by Definition~\ref{def:r03-two-sided-role-erasure-map} and normalized by Proposition~\ref{prop:r03-role-erasure-normal-form}.  Each mixed readout is expanded by Proposition~\ref{prop:r03-mixed-readout-normal-form-expansion}, and each Boolean mask action is restricted by Proposition~\ref{prop:r03-mask-filter-conservativity}.  The r04 package adds only release-candidate metadata and bundle-handoff statements around these already guarded objects.  Hence no new theorematic content is introduced.
\end{proof}

\begin{corollary}[r04 freeze-readiness condition]
\label{cor:r04-freeze-readiness-condition}
If the build audit reports no duplicate labels, no undefined references, no undefined citations, and no fatal LaTeX errors for both the Active Core and the Historical Witness, then Version~\docversion{} is structurally ready to be treated as a freeze candidate for the v3.35.0 pre-v4 QAM/tensor preparation release, subject only to the external release decision and any assigned Zenodo version DOI.
\end{corollary}

\begin{remark}[Reader-manual separation at r04]
\label{rem:r04-reader-manual-separation}
The pedagogical reader manual for the two-sided notation remains outside the Active Core.  It may later be written as a separate guide with diagrams, examples, and slower explanations.  The present Companion keeps only the compact formal grammar and the proof-neutrality guards needed by the governing proof line.
\end{remark}


% Quantum Sphaera Companion v3.35.0 release-final DOI insertion
% Release metadata only. No mathematical object is indexed by the DOI.

\subsection*{v3.35.0 DOI-bearing release freeze}
\addcontentsline{toc}{subsection}{v3.35.0 DOI-bearing release freeze}
\label{subsec:v3350-doi-bearing-release-freeze}

\noindent\textbf{Status.} \emph{Release-final metadata layer only; no new theorematic certificate, no global QAM namespace migration, no primitive mixed tensor, and no OP closure.}\par
\noindent\textbf{Version DOI.} The public version-specific Zenodo DOI for this release is
\[
  \href{https://doi.org/10.5281/zenodo.19750674}{10.5281/zenodo.19750674}.
\]
The Companion concept DOI remains \href{https://doi.org/10.5281/zenodo.19210728}{10.5281/zenodo.19210728}. The version DOI is release metadata only and is not a mathematical subscript, superscript, route tag, or proof object.

\begin{definition}[v3.35.0 release-final package]
\label{def:v3350-release-final-package}
The \emph{v3.35.0 release-final package} consists of the r04 release-candidate package of Definition~\ref{def:r04-pre-v4-qam-tensor-release-candidate-package}, the assigned public version DOI above, the Active Core PDF, the preserved Historical Witness PDF, the release manifest, checksums, audit report, Zenodo description, and release notes.  It is a release object, not a theorematic closure object.
\end{definition}

\begin{proposition}[v3.35.0 release-final conservativity]
\label{prop:v3350-release-final-conservativity}
The v3.35.0 release-final package is conservative over the r04 release-candidate package with respect to theorematic content.  DOI insertion, release-note generation, manifest creation, checksum creation, and Zenodo-description preparation do not create any new OP certificate, global QAM Hilbert-shift theorem, primitive mixed-tensor axiom, or replacement of the inherited QAM tuple namespace.
\end{proposition}

\begin{proof}
By Proposition~\ref{prop:r04-release-candidate-conservativity}, the r04 package is conservative over r03.  The present layer adds only release metadata and publication packaging: the version DOI, release file names, manifest/checksum records, audit report, Zenodo description, and release notes.  None of these objects changes the payload of any QAM role, tensor grammar, mask selector, mixed readout, or OP-facing proof obligation.  Hence theorematic content is unchanged.
\end{proof}


\begin{remark}[How the systems view reads the present visual diagnostics]\label{rem:qam-vm-os-visuals}
The same stack view also explains why Appendix~C may carry highly visible packet, Julia-boundary, and frame-flicker diagnostics without turning those figures into primary closure axioms. In the systems reading, such figures live on the visible-view layer: they are auditable readout surfaces of one and the same transported channel, not independent kernel identities. This is exactly why the OP~5 packet/breathing/frame-flicker block and the Julia-boundary block belong naturally to the same document architecture as the deferred kernel store. The figures make visible what the current theorem store is tracking, while closure itself still remains controlled at the structural layer beneath the pictures.
\end{remark}


\begin{remark}[QAM/systems crosswalk for the present document]\label{rem:qam-vm-os-crosswalk}
The present document uses the systems analogy in a disciplined layer-by-layer way. The following crosswalk fixes the intended reading.
\begin{center}
\renewcommand{\arraystretch}{1.12}
\begin{tabular}{|p{0.29\textwidth}|p{0.28\textwidth}|p{0.31\textwidth}|}
\hline
\textbf{QAM/document layer} & \textbf{Systems reading} & \textbf{Present mathematical status} \\
\hline
Visible coded tuples, packet diagnostics, Julia views, routed return surfaces & User-facing VM/view layer & Actively visible and auditable, but not by themselves closure-generating \\
\hline
Exposed payload kernels and later zero-kernel re-readings & VM-state identity layer & Deferred and prospective only; not yet activated on the public tuple surface \\
\hline
Admissible transport, branch/shell return, routed support organization & OS-level orchestration layer & Active theorem content in the present QAM export \\
\hline
Kernel normal forms, search keys, snapshots, clones, audits & ZFS-style theorem-store layer & Partly active as organizational discipline, partly deferred as later compression infrastructure \\
\hline
Pool/group structure of sphere objects and support families & Device / vdev / pool layer & Explanatory architecture only at the present stage \\
\hline
\end{tabular}
\end{center}
\end{remark}

\begin{remark}[Visual-evidence discipline under the systems reading]\label{rem:qam-vm-os-visual-discipline}
Under the crosswalk above, a figure may support only a controlled class of claims. It may expose \emph{where} a visible readout departs from a rigid carrier picture, \emph{which} packet or channel decomposition is plausible on the view layer, and \emph{why} a later kernel-first compression path is worth isolating. It may \emph{not} by itself create a new closure axiom, replace an active transport/branch/shell/routed theorem, or identify a kernel class solely from appearance. Any later proof use must therefore pass through the active theorem layer or through a deferred bridge that has itself already been discharged.
\end{remark}


\noindent\textbf{v3.16.0 checkpoint (deferred kernel-first path).} The deferred kernel-first material added in the present local v3.16.0 line should be read as one controlled checkpoint rather than as a second public language block. Its current role is to prepare a later proof-compression path while leaving the active coded tuple predicates unchanged.

\begin{remark}[Checkpoint summary for the deferred kernel-first path in the v3.16.0 public release]\label{rem:qam-zero-kernel-checkpoint}
At the present stage, the deferred kernel-first path contributes six concrete pieces of structure.
\begin{enumerate}[leftmargin=2em]
  \item It introduces a prospective recursive zero-kernel vocabulary in which the visible state surface $\langle 0,a\rangle$ may later be re-read through exposed payload kernels, kernel normal forms, and class-level readout factorization.
  \item It records a normal-form and search-key discipline so that later graph/tree based proof search can be tied to canonical kernel representatives rather than to repeated visible tuple surfaces.
  \item It isolates the missing bridge datum from active state surfaces to canonical kernels and identifies the direct visible payload as the simplest current kernel candidate on the present tuple surface.
  \item It tests the resulting kernel-first viewpoint on the three most relevant compression stages of the current QAM line: admissible transport, local branch-shell return, and routed family compression.
  \item It formulates an explicit activation package and a staged migration roadmap, making clear both what would have to be proved before activation and in which order the relevant closure/readout factorization work should be completed.
  \item It remains strictly deferred: no active public predicate, closure theorem, transport theorem, branch-shell theorem, routed-support theorem, or translation theorem is replaced by the kernel language at the present stage.
\end{enumerate}
\end{remark}

\begin{remark}[Why this checkpoint matters for later proof compression]\label{rem:qam-zero-kernel-checkpoint-why}
The checkpoint above already shows where later proof compression would come from. Transport is the first chain-level reduction, branch-shell return the first local surface reduction, routed support the first family reduction, and the translation theorem then becomes the first clause-level reduction at theorem scale. This does not yet shorten the active proofs automatically, but it isolates a concrete path by which set-theoretic comparison load could later be shifted from repeated visible tuple surfaces to kernel normal forms and one final class-level readout step.
\end{remark}

\begin{remark}[Navigation guide for the deferred kernel-first checkpoint]\label{rem:qam-zero-kernel-navigation}
A reader who wants the shortest route through the deferred material may follow the following order.
\begin{enumerate}[leftmargin=2em]
  \item \textbf{Purpose and fit.} Start with the roadmap-fit remark, the checkpoint summary, and the companion remark on why the checkpoint matters in order to see why the deferred path belongs to the inherited public-shape v3.16.0 release layer without changing the public branch/shell/routed extraction.
  \item \textbf{Vocabulary.} Then read the terminological convention for the deferred kernel path together with Definitions~\ref{def:qam-zero-kernel}--\ref{def:qam-kernel-key} for the later kernel vocabulary, primitive zero-objects, recursive reuse, aspect projections, normal forms, and internal search keys.
  \item \textbf{Compatibility.} Next read the axiom-coherence audit together with the current-limitation remark to see why this vocabulary remains deferred and does not compete with the active tuple layer.
  \item \textbf{Bridge.} Then read Definition~\ref{def:qam-kernel-bridge-datum}, Theorem~\ref{thm:qam-kernel-bridge}, Definition~\ref{def:qam-visible-payload}, and Theorem~\ref{thm:qam-direct-kernel-constructor} for the exact interface between visible state surfaces and later kernel-first comparison.
  \item \textbf{Stress tests.} Only afterwards read Propositions~\ref{prop:qam-kernel-first-transport}, \ref{prop:qam-kernel-first-branch-shell}, and \ref{prop:qam-kernel-first-routed}, which test the bridge on transport, local return, and routed family compression.
  \item \textbf{Activation and use.} Finally read Definition~\ref{def:qam-zero-kernel-activation-package}, Theorem~\ref{thm:qam-zero-kernel-activation-criterion}, the later activation-roadmap and reading-guide remarks, and then the dependency-map, minimal-proof-obligation, local-admissibility, and rewrite-schema remarks for the activation threshold, staged migration order, dependency logic, and later local rewrite test.
\end{enumerate}
In that order, the deferred block reads as one controlled proof-compression path rather than as a loose collection of auxiliary notes.
\end{remark}

\begin{remark}[Execution rule for the deferred kernel-first path]\label{rem:qam-zero-kernel-execution-rule}
The deferred block should be used in only one direction when later proof compression is attempted on the present tuple surface. Start with the visible state surface, expose the payload kernel, pass to kernel normal form, prove closure or readout agreement at class level, and lift back to the visible surface only at the end. If the later local admissibility test fails at any point, then the argument should remain on the active tuple surface rather than forcing a premature kernel-first rewrite.
\end{remark}

\begin{remark}[Present-use restriction for the deferred kernel-first path]\label{rem:qam-zero-kernel-present-use}
Until the activation criterion of Theorem~\ref{thm:qam-zero-kernel-activation-criterion} has been fully discharged on the current QAM surface, the deferred block should be used in only three ways inside the document: as later migration vocabulary, as an audit of where proof compression could occur, and as a record of staged stress tests. It should not be cited as a replacement for any active tuple predicate, transport theorem, branch-shell theorem, routed-support theorem, or translation theorem. In particular, whenever both a visible-surface proof and a deferred kernel-first restatement are available, the visible-surface proof remains authoritative on the inherited public-shape v3.16.0 layer.
\end{remark}


\begin{remark}[Active/deferred correspondence rule]\label{rem:qam-zero-kernel-correspondence}
The deferred kernel-first path may also be read as a controlled shadow of four active theorem layers, with one fixed discipline: the active visible-surface theorem remains authoritative, while the deferred statement records only the later compression route attached to that theorem. Concretely, Theorem~\ref{thm:qam-first-op-transport} is paired with Proposition~\ref{prop:qam-kernel-first-transport} as the chain-level compression test; Theorem~\ref{thm:qam-branch-shell} is paired with Proposition~\ref{prop:qam-kernel-first-branch-shell} as the local return-surface compression test; Theorem~\ref{thm:qam-routed-support} is paired with Proposition~\ref{prop:qam-kernel-first-routed} as the family-level compression test; and Theorem~\ref{thm:qam-translation} is paired with the later local admissibility reading of the translation theorem as the clause-level compression test. This correspondence is only a reading rule for later migration work. It does not transfer proof authority from the active theorem to its deferred companion; it only says where a later kernel-first rewrite would have to pay off if the activation package were eventually discharged.
\end{remark}

\subsection*{Prospective recursive zero-kernel scheme (deferred architecture note)}

The following block records a possible later re-reading of the current $\langle 0,a\rangle$-entry as a recursive order-$0$ readout kernel. It is included here only as a controlled design note for future transfer work. It is \emph{not} activated as part of the present public QAM tag system and does not modify the inherited public-shape coded predicates from v3.16.0.

\begin{remark}[Terminological convention for the deferred kernel path]\label{rem:qam-zero-kernel-terminology}
Throughout the deferred block, four layers should be kept distinct and named consistently. The expression \emph{visible state surface} refers to the active public state-code presentation $s=\langle 0,a\rangle$ on the current tuple level. The expression \emph{exposed payload kernel} refers to the payload extracted from such a visible state surface by the deferred kernel constructor. The expression \emph{kernel normal form} refers to the canonical representative selected from the admissible kernel equivalence class. Finally, \emph{class-level readout factorization} means that closure-relevant readout depends only on that kernel normal-form class and is lifted back to the visible state surface only at the end. This convention is purely terminological, but it keeps the active tuple language and the deferred kernel language aligned.
\end{remark}

\begin{definition}[Carrier/filter readout kernels]\label{def:qam-zero-kernel}
Let $\mathcal C$ be a class of admissible carriers, let $\Gamma$ be a class of admissible filters, and let $\mathcal A$ be a class of canonical readout kernels. For each $\gamma\in\Gamma$, assume a partial readout map
\[
\rho_\gamma:\mathcal C \dashrightarrow \mathcal A,
\]
whose defined values are written
\[
a=\rho_\gamma(T)\in\mathcal A.
\]
We read $a$ as the minimal canonical readout kernel of the carrier $T$ relative to the filter $\gamma$.
\end{definition}

\begin{definition}[Primitive zero-objects]\label{def:qam-zero-objects}
Given $a,b\in\mathcal A$, define the primitive zero-objects
\[
Z_0(a):=\langle 0,a\rangle,
\qquad
Z_0(a,b):=\langle 0,a,b\rangle.
\]
The intended reading is that $Z_0(a)$ is the unary collapsed readout attached to the kernel $a$, while $Z_0(a,b)$ is the binary collapsed readout attached to the kernel pair $(a,b)$.
\end{definition}

\begin{axiom}[Recursive reusability of kernels]\label{ax:qam-zero-recursive}
Every canonical readout kernel may again serve as an effective carrier for further admissible readout. Concretely, for admissible $\eta\in\Gamma$, one may again form
\[
\rho_\eta(a)\in\mathcal A
\]
whenever the corresponding readout is defined. Thus the kernel class $\mathcal A$ is recursively reusable as an effective carrier layer.
\end{axiom}

\begin{definition}[Aspect projections on kernels]\label{def:qam-kernel-aspects}
Assume that each kernel $a\in\mathcal A$ comes equipped with aspect projections
\[
\pi_{\mathrm{type}}(a),\qquad
\pi_{\mathrm{class}}(a),\qquad
\pi_{\mathrm{sig}}(a),
\]
recording, respectively, its visible type aspect, filter/class aspect, and structured signature aspect. These are not treated as three competing objects; they are three admissible readings of the same kernel.
\end{definition}

\begin{lemma}[Unary and binary kernel collapse]\label{lem:qam-zero-collapse}
If $a=\rho_\gamma(T)\in\mathcal A$, then the unary zero-object $Z_0(a)=\langle 0,a\rangle$ is well-defined. If additionally $b=\rho_\eta(U)\in\mathcal A$, then the binary zero-object $Z_0(a,b)=\langle 0,a,b\rangle$ is well-defined.
\end{lemma}

\begin{proof}
The unary claim follows directly from Definition~\ref{def:qam-zero-objects} once $a\in\mathcal A$ is given by Definition~\ref{def:qam-zero-kernel}. The binary claim is identical after adjoining a second admissible kernel $b\in\mathcal A$.
\end{proof}

\begin{remark}[Compressed relation to the current $\langle 0,a\rangle$ state code]\label{rem:qam-zero-compressed}
This deferred kernel notation is intended only as a future re-reading of the present visible state surface. In a later binary-tag or order-sensitive QAM line, one may reinterpret the current $\langle 0,a\rangle$ notation as a compressed public surface for a filtered order-$0$ readout kernel. For the present version, however, the active public predicates remain exactly those of the current tagged tuple language.
\end{remark}

\begin{remark}[Tree/graph viewpoint for recursive kernels]\label{rem:qam-zero-graph}
The recursive kernel picture admits an immediate combinatorial visualization. One may regard repeated kernel readout
\[
T \xrightarrow{\rho_\gamma} a \xrightarrow{\rho_\eta} a' \xrightarrow{\rho_\theta} a''
\]
as a rooted readout tree whose vertices are kernels and whose edges are admissible filter applications. Equally, one may organize the same data as a directed multigraph: unary zero-objects $Z_0(a)$ mark collapsed kernel vertices, while binary zero-objects $Z_0(a,b)$ record elementary visible relations between kernel vertices. This suggests a later graph-based search or proof-compression interface without changing the current public tuple system.
\end{remark}


\subsection*{Prospective kernel normal forms and search keys (deferred architecture note)}

The same deferred kernel language also suggests a later search discipline in which recursive kernel proofs are first normalized and only then indexed. This block is likewise prospective only: it does not alter the active inherited public-shape tuple predicates.

\begin{definition}[Kernel equivalence and canonical normal form]\label{def:qam-kernel-normal-form}
Assume an admissible equivalence relation $\sim$ on $\mathcal A$ such that equivalent kernels have the same admissible aspect projections and the same closure-relevant readout content. A \emph{kernel normal form} is a partial map
\[
\operatorname{nf}:\mathcal A \dashrightarrow \mathcal A
\]
selecting, whenever defined, one canonical representative from each admissible equivalence class. Thus
\[
a\sim b \Longrightarrow \operatorname{nf}(a)=\operatorname{nf}(b)
\]
whenever both normal forms are defined.
\end{definition}

\begin{definition}[Kernel search keys]\label{def:qam-kernel-key}
Assume a fixed admissible serialization
\[
\operatorname{ser}:\operatorname{nf}(\mathcal A)\to \Sigma^*
\]
of canonical kernels into finite strings over an alphabet $\Sigma$, together with a deterministic key map
\[
\kappa := H\circ \operatorname{ser}\circ \operatorname{nf}.
\]
The value $\kappa(a)$ is called the \emph{kernel search key} of $a$ whenever the normal form is defined. It is intended as an internal search/index key rather than as part of the visible mathematical surface.
\end{definition}

\begin{lemma}[Normal-form invariance of primitive zero-objects]\label{lem:qam-zero-normal-form}
Assume that $a,a',b,b'\in\mathcal A$ satisfy $a\sim a'$ and $b\sim b'$, and that their normal forms are defined. Then the unary zero-objects $Z_0(a)$ and $Z_0(a')$ determine the same canonical unary kernel surface, while the binary zero-objects $Z_0(a,b)$ and $Z_0(a',b')$ determine the same canonical binary kernel surface after passage to kernel normal form.
\end{lemma}

\begin{proof}
By Definition~\ref{def:qam-kernel-normal-form}, equivalence implies equality of canonical normal forms. Applying the same normal-form representative at the unary or binary zero-object level yields the same canonical collapsed surface in either case.
\end{proof}

\begin{remark}[Graph search and proof-compression interface]\label{rem:qam-kernel-search}
In the graph viewpoint introduced above, one may therefore index vertices by canonical kernel normal forms or, internally, by their kernel search keys $\kappa(a)$. Repeated readout branches that land in the same canonical kernel class can then be merged, memoized, or searched only once. This is exactly the kind of mechanism by which recursive kernel language could later compress set-theoretic proof search without changing the visible public theorem statements.
\end{remark}


\subsection*{Axiom-coherence audit and prospective proof reductions (deferred audit note)}

The deferred zero-kernel language should be read together with the active public QAM tuple layer under the following discipline.

\begin{remark}[Axiom-coherence audit]\label{rem:qam-zero-axiom-audit}
At the present stage, the deferred zero-kernel layer is formally compatible with the active QAM tuple layer for four reasons.
\begin{enumerate}
\item The active tuple predicates remain unchanged: in particular,
\[
\mathrm{QState}(x)\Longleftrightarrow \exists a\,(x=\langle 0,a\rangle)
\]
continues to define the public visible state surface.
\item The deferred zero-kernel definitions do not introduce a second active state predicate. They only record a possible later re-reading of the same public surface, as stated in Remark~\ref{rem:qam-zero-compressed}.
\item The recursive kernel axiom (Axiom~\ref{ax:qam-zero-recursive}) and the kernel normal-form discipline (Definition~\ref{def:qam-kernel-normal-form}) are prospective only; no active theorem in the inherited public-shape QAM chain depends on them as hypotheses.
\item Consequently, the current closure-class, admissible-transport, branch/shell, routed-support, and translation theorems remain exactly those of the active tuple language. The deferred kernel layer adds a controlled future interpretation, not a rival public axiom system.
\end{enumerate}
\end{remark}

\begin{remark}[Where kernel notation could already shorten proofs]\label{rem:qam-zero-proof-opportunities}
Even before formal activation, the deferred kernel language already points to three concrete proof-compression opportunities.
\begin{enumerate}
\item In the first QAM transport theorem, the chain
\[
s_{\mathrm{in}},\ s_2,\ s_3,\ s_4,\ t_{\mathrm{fwd}},\ t_{\mathrm{bwd}}
\]
is currently carried at full state-code level. A later kernel reduction could first pass each state to a canonical readout kernel and then prove closure agreement on kernel classes before lifting the result back to coded states.
\item In the branch-shell theorem and the routed admissible support theorem, each return realization currently enters through full branch/shell state data. A kernel-first restatement could instead reduce every admissible branch-shell realization to unary kernel surfaces $Z_0(a)$ and compare routed sectors only through their canonical kernel classes.
\item In closure-compatible readout arguments, one currently factors raw readouts through the class-target map $J$ after proving $\mathrm{QClEq}(s,t)$. A later kernel language could absorb much of this step by treating the closure-relevant image itself as the primitive kernel surface attached to the state.
\end{enumerate}
In all three cases, the expected gain is not a new theorem but a shorter proof path: reduce to kernels first, perform the closure comparison there, and only then return to the visible tuple layer.
\end{remark}

\begin{remark}[Current limitation of the kernel reduction]\label{rem:qam-zero-current-limit}
The proof-compression opportunities recorded just above are not yet activated as a formal replacement inside the present proofs. The missing step is a theorem-level bridge showing that every active coded state relevant to the admissible OP fragment carries a canonical kernel representative with the same closure-relevant content. Until such a bridge is proved, the zero-kernel language should be treated as a planning interface rather than as a proof rule.
\end{remark}



\begin{definition}[Prospective kernel bridge datum on active state surfaces]\label{def:qam-kernel-bridge-datum}
A \emph{kernel bridge datum} for the active admissible OP fragment consists of a partial map
\[
\operatorname{ker}:\mathcal S^{\mathrm{QAM}}_{\mathrm{act}} \dashrightarrow \mathcal A,
\]
where $\mathcal S^{\mathrm{QAM}}_{\mathrm{act}}$ is the class of active coded QAM-state surfaces relevant to the present admissible OP fragment, such that the following hold whenever the displayed expressions are defined.
\begin{enumerate}
\item \textbf{surface realization:}
\[
Z_0(\operatorname{ker}(s))=s;
\]
\item \textbf{closure reflection:} if $\mathrm{QClEq}(s,t)$, then
\[
\operatorname{nf}(\operatorname{ker}(s))=\operatorname{nf}(\operatorname{ker}(t));
\]
\item \textbf{closure-readout factorization:} for every closure-compatible coded readout $r$, there exists a map
\[
\mathcal R_r: \operatorname{nf}(\mathcal A) \to \operatorname{Im}(J_r)
\]
into the closure-relevant target of the associated class-factor map $J_r$ such that whenever $\mathrm{QRead}(r,s,w)$ holds,
\[
J_r(w)=\mathcal R_r(\operatorname{nf}(\operatorname{ker}(s))).
\]
\end{enumerate}
\end{definition}

\begin{theorem}[Prospective bridge from active state surfaces to canonical kernels]\label{thm:qam-kernel-bridge}
Assume a kernel bridge datum in the sense of Definition~\ref{def:qam-kernel-bridge-datum}. Then every active closure comparison in the admissible OP fragment may be reduced to canonical kernel comparison before returning to the visible state surface. In particular:
\begin{enumerate}
\item if $\mathrm{QAdmTransport}(o,s,t)$, then
\[
\operatorname{nf}(\operatorname{ker}(s))=\operatorname{nf}(\operatorname{ker}(t));
\]
\item if $t_{b,\sigma}$ is a branch-shell coded return realization as in Theorem~\ref{thm:qam-branch-shell}, then
\[
\operatorname{nf}(\operatorname{ker}(s_{\mathrm{in}}))=\operatorname{nf}(\operatorname{ker}(t_{b,\sigma}));
\]
\item if $u$ is a normalized admissible routed support code as in Theorem~\ref{thm:qam-routed-support}, then for every closure-compatible coded readout $r$ the closure-relevant image of the routed return is determined solely by the canonical kernel class
\[
\operatorname{nf}(\operatorname{ker}(s_{\mathrm{in}})).
\]
\end{enumerate}
Consequently, the present transport, branch/shell, and routed-support proofs admit a future kernel-first proof skeleton: first pass to $\operatorname{nf}(\operatorname{ker}(\cdot))$, then compare closure there, and only afterwards return to the visible QAM state surface.
\end{theorem}

\begin{proof}
For (1), admissible transport implies $\mathrm{QClEq}(s,t)$ by the active closure-preservation axiom. The closure-reflection clause of Definition~\ref{def:qam-kernel-bridge-datum} therefore yields
\[
\operatorname{nf}(\operatorname{ker}(s))=\operatorname{nf}(\operatorname{ker}(t)).
\]
For (2), Theorem~\ref{thm:qam-branch-shell} gives $\mathrm{QClEq}(s_{\mathrm{in}},t_{b,\sigma})$, so the same closure-reflection clause yields equality of canonical kernel representatives. For (3), Theorem~\ref{thm:qam-routed-support} states that the routed return has the same closure-relevant image as the coded source state under every closure-compatible coded readout. By clause (3) of Definition~\ref{def:qam-kernel-bridge-datum}, those closure-relevant images factor through the same canonical kernel normal form. Hence the routed readout depends only on $\operatorname{nf}(\operatorname{ker}(s_{\mathrm{in}}))$. The final sentence is simply the proof reading order extracted from (1)--(3).
\end{proof}

\begin{remark}[What the bridge theorem still leaves open]\label{rem:qam-kernel-bridge-open}
Theorem~\ref{thm:qam-kernel-bridge} is intentionally conditional. It does not yet prove that a kernel bridge datum exists inside the active public tuple system; it isolates exactly the missing interface that would be needed to turn the deferred zero-kernel language into an actual proof tool. This also identifies the next genuine activation step: construct $\operatorname{ker}$ canonically from active coded states and verify clauses (1)--(3) without leaving the current OP/QAM semantics.
\end{remark}


\begin{definition}[Visible payload extraction on active state surfaces]\label{def:qam-visible-payload}
For every active coded state surface $s$ in the present public tuple language, define its \emph{exposed payload kernel} $\operatorname{vp}(s)$ by
\[
\operatorname{vp}(s)=a \quad\Longleftrightarrow\quad s=\langle 0,a\rangle.
\]
Equivalently, $\operatorname{vp}$ is the unique projection that removes the visible state tag $0$ and exposes the payload already carried by the active state surface.
\end{definition}

\begin{definition}[Prospective direct kernel constructor]\label{def:qam-direct-kernel-constructor}
Assume that the exposed payload kernel of every active coded state surface relevant to the admissible OP fragment belongs to the deferred kernel class $\mathcal A$. Then define the \emph{direct kernel constructor}
\[
\operatorname{ker}_{\mathrm{dir}} : \mathcal S^{\mathrm{QAM}}_{\mathrm{act}} \dashrightarrow \mathcal A,
\qquad
\operatorname{ker}_{\mathrm{dir}}(s):=\operatorname{vp}(s).
\]
Thus, on the visible state surface, the simplest possible kernel candidate is obtained by forgetting only the active state tag and keeping the payload itself as the deferred kernel representative.
\end{definition}

\begin{theorem}[Simplest concrete kernel construction on the present state surface]\label{thm:qam-direct-kernel-constructor}
Assume the following four conditions for the active admissible OP fragment.
\begin{enumerate}
\item \textbf{payload admissibility:} every active coded state surface $s$ relevant to the fragment has the form $s=\langle 0,a\rangle$ with $a\in\mathcal A$;
\item \textbf{visible zero-surface realization:} for every such $a$,
\[
Z_0(a)=\langle 0,a\rangle;
\]
\item \textbf{closure reflection on exposed payload kernels:} whenever $\mathrm{QClEq}(\langle 0,a\rangle,\langle 0,b\rangle)$,
\[
\operatorname{nf}(a)=\operatorname{nf}(b);
\]
\item \textbf{class-level readout factorization:} for every closure-compatible coded readout $r$, there exists a map
\[
\widetilde{\mathcal R}_r:\operatorname{nf}(\mathcal A)\to \operatorname{Im}(J_r)
\]
such that whenever $\mathrm{QRead}(r,\langle 0,a\rangle,w)$,
\[
J_r(w)=\widetilde{\mathcal R}_r(\operatorname{nf}(a)).
\]
\end{enumerate}
Then $\operatorname{ker}_{\mathrm{dir}}$ is a kernel bridge datum in the sense of Definition~\ref{def:qam-kernel-bridge-datum}. Consequently Theorem~\ref{thm:qam-kernel-bridge} applies with
\[
\operatorname{ker}=\operatorname{ker}_{\mathrm{dir}}.
\]
\end{theorem}

\begin{proof}
Clause (1) ensures that $\operatorname{ker}_{\mathrm{dir}}$ is defined on every active state surface relevant to the fragment. By Definition~\ref{def:qam-visible-payload}, if $s=\langle 0,a\rangle$, then $\operatorname{ker}_{\mathrm{dir}}(s)=a$. Clause (2) therefore gives
\[
Z_0(\operatorname{ker}_{\mathrm{dir}}(s))=Z_0(a)=\langle 0,a\rangle=s,
\]
which is exactly the surface-realization clause of Definition~\ref{def:qam-kernel-bridge-datum}. For closure reflection, if $\mathrm{QClEq}(s,t)$ with $s=\langle 0,a\rangle$ and $t=\langle 0,b\rangle$, clause (3) gives
\[
\operatorname{nf}(\operatorname{ker}_{\mathrm{dir}}(s))=\operatorname{nf}(a)=\operatorname{nf}(b)=\operatorname{nf}(\operatorname{ker}_{\mathrm{dir}}(t)).
\]
Finally, clause (4) is exactly the required closure-readout factorization after substituting $\operatorname{ker}_{\mathrm{dir}}(\langle 0,a\rangle)=a$. Hence all three clauses of Definition~\ref{def:qam-kernel-bridge-datum} hold.
\end{proof}

\begin{remark}[What this construction shows]\label{rem:qam-direct-kernel-constructor-meaning}
Theorem~\ref{thm:qam-direct-kernel-constructor} shows that the difficult part of the bridge is not the extraction step itself: on the current public state surface, the most economical candidate is simply the exposed payload of the active state surface. The real work is pushed into the three semantic checks: that the payload already lies in the admissible kernel class, that closure equivalence is visible at payload-normal-form level, and that closure-relevant readout factors through that same normal form. If these three checks can later be discharged, then many current state-level proofs may be shortened by passing immediately from $\langle 0,a\rangle$ to $a$ before performing the closure comparison.
\end{remark}

\subsection*{QAM minimal predicates and coded relations}
\begin{quote}\small
\noindent\textbf{Status.} \emph{Active formal-language layer.}\par
\noindent\textbf{Block function.} This block introduces the smallest coded predicate/relation vocabulary needed to read the present admissible OP fragment in QAM-over-ZF form.
\end{quote}

We now pass from the entry-layer preview to the first explicit predicate layer of QAM-over-ZF. At this stage, the QAM language still remains conservative over ZF: the predicates below are intended as definable recognition or relation clauses on tagged set-codes rather than as commitments to a new ontology.

\begin{definition}[Primitive QAM predicates]
We introduce the primitive predicates
\[
\begin{aligned}
&\mathrm{QState}(x),\quad \mathrm{QLayer}(x),\quad \mathrm{QOp}(x),\quad \mathrm{QReadout}(x),\\[0.15em]
&\mathrm{QBranch}(x),\quad \mathrm{QShell}(x),\quad \mathrm{QWitness}(x),\quad \mathrm{QTensor}(x).
\end{aligned}
\]
with the intended ZF-grounded readings
\[
\mathrm{QState}(x)\;:\Longleftrightarrow\;\exists a\,(x=\langle 0,a\rangle),
\]
\[
\mathrm{QLayer}(x)\;:\Longleftrightarrow\;\exists \ell\,\exists D\,(x=\langle 1,\ell,D\rangle),
\]
\[
\mathrm{QOp}(x)\;:\Longleftrightarrow\;\exists \ell_1\,\exists \ell_2\,\exists G\,(x=\langle 2,\ell_1,\ell_2,G\rangle),
\]
\[
\mathrm{QReadout}(x)\;:\Longleftrightarrow\;\exists \ell\,\exists W\,\exists R\,(x=\langle 3,\ell,W,R\rangle),
\]
\[
\mathrm{QBranch}(x)\;:\Longleftrightarrow\;\exists b\,\exists B\,(x=\langle 4,b,B\rangle),
\]
\[
\mathrm{QShell}(x)\;:\Longleftrightarrow\;\exists r\,\exists s\,\exists T\,(x=\langle 5,r,s,T\rangle),
\]
\[
\mathrm{QWitness}(x)\;:\Longleftrightarrow\;\exists u\,(x=\langle 6,u\rangle),
\]
\[
\mathrm{QTensor}(x)\;:\Longleftrightarrow\;\exists p\,\exists q\,(x=\langle 7,p,q\rangle).
\]
\end{definition}

\begin{remark}
At this stage, these predicates are still recognition clauses rather than a full independent axiom system. Their purpose is to let the Companion speak explicitly about coded states, layers, operators, readouts, branches, shell/refinement objects, witnesses, and tensor composites inside one ZF-grounded language layer.
\end{remark}

\begin{definition}[Primitive QAM coded relations]
We also introduce the primitive QAM-style relations
\[
\mathrm{QAdmTransport}(o,s,t),\qquad
\mathrm{QClEq}(s,t),\qquad
\mathrm{QRead}(r,s,w),
\]
\[
\mathrm{QBranchRef}(b,s,t),\qquad
\mathrm{QShellRef}(\sigma,s,t),\qquad
\mathrm{QTensorAdm}(x,y).
\]
Their intended readings are as follows:
\begin{enumerate}
\item $\mathrm{QAdmTransport}(o,s,t)$ means that the operator-code $o$ transports the coded state $s$ admissibly to the coded state $t$;
\item $\mathrm{QClEq}(s,t)$ means that $s$ and $t$ belong to the same coded closure class;
\item $\mathrm{QRead}(r,s,w)$ means that the readout-code $r$ reads the coded state $s$ into the coded observable value $w$;
\item $\mathrm{QBranchRef}(b,s,t)$ means that the branch-code $b$ refines the coded state $s$ to the coded state $t$;
\item $\mathrm{QShellRef}(\sigma,s,t)$ means that the shell/refinement-code $\sigma$ refines the coded state $s$ to the coded state $t$;
\item $\mathrm{QTensorAdm}(x,y)$ means that the tensor-coded composite $x\otimes y$ is admitted as a legitimate coded QAM object.
\end{enumerate}
\end{definition}

\begin{remark}
These relations are the first minimal bridge between the current OP/branch/shell architecture and its coded QAM-over-ZF form. In particular, admissibility and closure-class preservation already appear here as primitive language-level notions rather than as only informal motivation.
\end{remark}

\begin{definition}[First QAM translation target]
The first intended translation target of the QAM-over-ZF layer is the admissible OP fragment consisting of:
\begin{enumerate}
\item the OP~1/OP~5 source-return closure architecture;
\item the intermediate OP~2--OP~4 transport bridge;
\item admissible branch transport;
\item OP~3 shell-admissible refinement;
\item admissible routed support together with the visible packet/phase-lock/epsilon screening layer.
\end{enumerate}
\end{definition}

\begin{remark}
The significance of the present predicate layer is that QAM is no longer only a roadmap note. Even at this first stage, the language already distinguishes object-kinds and relation-kinds closely enough to serve as a faithful coded entry interface for the current admissible OP fragment.
\end{remark}



\subsection*{QAM transport, closure, and readout schemata}
\begin{quote}\small
\noindent\textbf{Status.} \emph{Active theorem-core block inside the QAM formal layer.}\par
\noindent\textbf{Block function.} This block turns the predicate language into actual transport, closure-class, and readout principles, so it is the first place where the QAM interface becomes theorem-usable.
\end{quote}

We now pass from the minimal predicate layer to the first theorematic QAM transport/closure layer. The point of this section is still conservative: it does not yet claim a fully independent QAM foundation, but it does begin to formulate admissible transport, closure-class preservation, and closure-compatible readout directly in the coded QAM-over-ZF language.

\begin{definition}[QAM-coded transport graph]
A QAM-operator code
\[
o=\langle 2,\ell_1,\ell_2,G\rangle
\]
is called a \emph{coded transport graph} if $G$ is a set of ordered pairs coding a functional transport rule from coded states on layer $\ell_1$ to coded states on layer $\ell_2$.
\end{definition}

\begin{definition}[QAM admissible transport]
For a QAM-operator code $o$ and coded states $s,t$, we write
\[
\mathrm{QAdmTransport}(o,s,t)
\]
iff:
\begin{enumerate}
\item $o$ is a coded transport graph;
\item $s$ is a coded state on the source layer of $o$;
\item $t$ is a coded state on the target layer of $o$;
\item the graph-code of $o$ sends $s$ to $t$;
\item the transport from $s$ to $t$ preserves the intended closure-bearing structural mode.
\end{enumerate}
\end{definition}

\begin{remark}
At this stage, the final clause is still schematic. Its precise intended mathematical content is supplied by the coded closure-class relation below.
\end{remark}

\begin{definition}[QAM closure-class equivalence]
The relation
\[
\mathrm{QClEq}(s,t)
\]
is the coded closure-class relation on QAM-state codes. It is intended to mean that $s$ and $t$ represent the same closure-bearing structural mode, even if they differ by admissible transport, admissible frame-shifted realization, admissible branch refinement, or admissible shell refinement.
\end{definition}

\begin{axiom}[QAM closure-class schema]
The relation $\mathrm{QClEq}$ is required to satisfy:
\begin{enumerate}
\item \textbf{reflexivity:}
\[
\mathrm{QClEq}(s,s);
\]
\item \textbf{symmetry:}
\[
\mathrm{QClEq}(s,t)\Longrightarrow \mathrm{QClEq}(t,s);
\]
\item \textbf{transitivity:}
\[
\mathrm{QClEq}(s,t)\wedge \mathrm{QClEq}(t,u)\Longrightarrow \mathrm{QClEq}(s,u).
\]
\end{enumerate}
\end{axiom}

\begin{remark}
Thus $\mathrm{QClEq}$ is the coded QAM analogue of the native closure carrier-class relation used in the current OP architecture.
\end{remark}

\begin{axiom}[QAM admissibility implies closure preservation]
For every operator code $o$ and state codes $s,t$,
\[
\mathrm{QAdmTransport}(o,s,t)\Longrightarrow \mathrm{QClEq}(s,t).
\]
\end{axiom}

\begin{remark}
This is the first central QAM principle. Admissibility is not defined merely by existence of transport, but by preservation of closure class.
\end{remark}

\begin{definition}[QAM-coded readout]
A QAM-readout code
\[
r=\langle 3,\ell,W,R\rangle
\]
is called a \emph{coded readout} if $R$ is a graph-code sending coded states on layer $\ell$ into coded observable values in the target code $W$.
\end{definition}

\begin{definition}[QAM readout relation]
We write
\[
\mathrm{QRead}(r,s,w)
\]
iff $r$ is a coded readout, $s$ is a coded state on the source layer of $r$, and the graph-code of $r$ sends $s$ to the coded observable value $w$.
\end{definition}

\begin{definition}[QAM closure-compatible readout]\label{def:qam-closure-compatible-readout}
A coded readout $r$ is called \emph{closure-compatible} if there exists a coded class-target map $J$ such that for all coded states $s,t$,
\[
\mathrm{QClEq}(s,t)\Longrightarrow
\bigl(
\forall w_1\,\forall w_2\,
(\mathrm{QRead}(r,s,w_1)\wedge \mathrm{QRead}(r,t,w_2)\Longrightarrow J(w_1)=J(w_2))
\bigr).
\]
\end{definition}

\begin{remark}
This is the coded QAM version of readout factorization through closure class. Raw readouts may differ, but their closure-relevant image under $J$ must agree whenever the source states belong to the same closure class.
\end{remark}

\begin{definition}[First tensor-compatibility placeholder]
We reserve the tensor-admissibility predicate
\[
\mathrm{QTensorAdm}(x,y)
\]
for the statement that the tensor-coded composite $x\otimes y$ is admitted as a legitimate coded QAM object of the intended kind.
\end{definition}

\begin{proposition}[First QAM transport principle]
Let $o$ be a coded QAM-operator, let $s,t$ be coded QAM-states, and let $r$ be a closure-compatible coded readout. If
\[
\mathrm{QAdmTransport}(o,s,t),
\]
then:
\begin{enumerate}
\item \[
\mathrm{QClEq}(s,t);
\]
\item the closure-relevant image of any readout of $s$ agrees with the closure-relevant image of any readout of $t$ under the factor map $J$ associated with $r$.
\end{enumerate}
\end{proposition}

\begin{proof}
The first claim is exactly the admissibility-implies-closure-preservation schema. The second follows from closure-compatibility of the readout: since $\mathrm{QClEq}(s,t)$ holds, any coded readouts of $s$ and $t$ are identified after passage through the closure-relevant class-target map $J$.
\end{proof}

\begin{theorem}[First QAM version of the OP transport principle]\label{thm:qam-first-op-transport}
Assume:
\begin{enumerate}
\item there is a coded source state $s_{\mathrm{in}}$;
\item there are coded intermediate states $s_2,s_3,s_4$;
\item there are coded operator maps $o_2,o_3,o_4,o_{\mathrm{fwd}},o_{\mathrm{bwd}}$;
\item the transports
\[
\begin{aligned}
&\mathrm{QAdmTransport}(o_2,s_{\mathrm{in}},s_2),\qquad
\mathrm{QAdmTransport}(o_3,s_2,s_3),\\[0.15em]
&\mathrm{QAdmTransport}(o_4,s_3,s_4),\qquad
\mathrm{QAdmTransport}(o_{\mathrm{fwd}},s_{\mathrm{in}},t_{\mathrm{fwd}}),\\[0.15em]
&\mathrm{QAdmTransport}(o_{\mathrm{bwd}},s_{\mathrm{in}},t_{\mathrm{bwd}}).
\end{aligned}
\]
all hold;
\item the relevant coded readout is closure-compatible.
\end{enumerate}
Then all coded states
\[
s_{\mathrm{in}},\ s_2,\ s_3,\ s_4,\ t_{\mathrm{fwd}},\ t_{\mathrm{bwd}}
\]
belong to the same coded closure class. Consequently, all closure-relevant readouts of these states agree after passage through the closure-factor map.
\end{theorem}

\begin{proof}
Each admissible transport preserves closure class by the admissibility schema. Therefore
\[
\mathrm{QClEq}(s_{\mathrm{in}},s_2),\qquad
\mathrm{QClEq}(s_2,s_3),\qquad
\mathrm{QClEq}(s_3,s_4).
\]
By transitivity,
\[
\mathrm{QClEq}(s_{\mathrm{in}},s_4).
\]
Likewise,
\[
\mathrm{QClEq}(s_{\mathrm{in}},t_{\mathrm{fwd}})
\qquad\text{and}\qquad
\mathrm{QClEq}(s_{\mathrm{in}},t_{\mathrm{bwd}})
\]
follow from admissible forward and backward transport. Using transitivity again, all listed coded states lie in the same closure class. Closure-compatible readout then identifies their closure-relevant readout content after factorization through the associated class-target map.
\end{proof}

\begin{proposition}[Prospective kernel-first reduction of the first transport theorem]\label{prop:qam-kernel-first-transport}
Assume the direct kernel constructor of Definition~\ref{def:qam-direct-kernel-constructor} and the bridge conditions of Theorem~\ref{thm:qam-direct-kernel-constructor}. Write
\[
s_{\mathrm{in}}=\langle 0,a_{\mathrm{in}}\rangle,
\qquad
s_2=\langle 0,a_2\rangle,
\qquad
s_3=\langle 0,a_3\rangle,
\qquad
s_4=\langle 0,a_4\rangle,
\]
\[
t_{\mathrm{fwd}}=\langle 0,a_{\mathrm{fwd}}\rangle,
\qquad
t_{\mathrm{bwd}}=\langle 0,a_{\mathrm{bwd}}\rangle.
\]
Then the conclusion of Theorem~\ref{thm:qam-first-op-transport} may be read kernel-first as follows:
\begin{enumerate}
\item all exposed transport kernels lie in one canonical normal-form class,
\[
\operatorname{nf}(a_{\mathrm{in}})=\operatorname{nf}(a_2)=\operatorname{nf}(a_3)=\operatorname{nf}(a_4)=\operatorname{nf}(a_{\mathrm{fwd}})=\operatorname{nf}(a_{\mathrm{bwd}}),
\]
\item hence every closure-compatible coded readout of the transport chain factors through that single canonical kernel class,
\item therefore the first transport theorem may be re-read as a kernel-normal-form preservation statement along the admissible transport chain before returning to the visible QAM state surface.
\end{enumerate}
\end{proposition}

\begin{proof}
Theorem~\ref{thm:qam-first-op-transport} states that all six visible coded states belong to the same coded closure class. Since each state lies on the present state surface, the direct kernel constructor exposes exactly the exposed payload kernel of each state as its deferred kernel representative. The closure-reflection clause of Theorem~\ref{thm:qam-direct-kernel-constructor} therefore yields equality of the six corresponding kernel normal forms. By the class-level readout factorization clause, every closure-compatible coded readout depends only on that common kernel normal form. Hence the whole transport chain may be treated kernel-first and lifted back to the visible state surface only at the end.
\end{proof}

\begin{remark}[Axiom audit after the kernel-first transport reduction]\label{rem:qam-kernel-first-transport-audit}
Proposition~\ref{prop:qam-kernel-first-transport} does \emph{not} strengthen the active proof status of Theorem~\ref{thm:qam-first-op-transport}. The active theorem still rests only on admissible transport, closure-class preservation, transitivity of coded closure, and closure-compatible coded readout. The kernel-first proposition adds force only under the deferred payload-level bridge hypotheses isolated in Definition~\ref{def:qam-kernel-bridge-datum} and Theorem~\ref{thm:qam-direct-kernel-constructor}. Thus no new hidden axiom has entered the live transport theorem; only a prospective kernel-level compression of its proof surface has been recorded.
\end{remark}

\begin{remark}[How much proof load the kernel-first transport reading removes]\label{rem:qam-kernel-first-transport-savings}
For the present transport theorem the admissible chain contains five explicit transport links:
\[
s_{\mathrm{in}}\to s_2\to s_3\to s_4,
\qquad
s_{\mathrm{in}}\to t_{\mathrm{fwd}},
\qquad
s_{\mathrm{in}}\to t_{\mathrm{bwd}}.
\]
So the gain is again \emph{not} a reduction in the number of transport links, but a reduction in the size of the object carried through each link. On the visible state surface the proof transports closure across five state-surface steps and then factors a closure-compatible readout through the resulting six-state closure class. In the deferred kernel-first reading, those same five links are re-read as five kernel normal-form identifications, followed by one readout through the common kernel class. If one measures a visible-state transport comparison by cost $C_{\mathrm{surf}}$ and a kernel normal-form comparison by cost $C_{\mathrm{ker}}$, then the conceptual load changes from roughly
\[
5\,C_{\mathrm{surf}} + C_{\mathrm{read}}
\]
to
\[
5\,C_{\mathrm{ker}} + C_{\mathrm{class}},
\]
with intended gain $C_{\mathrm{ker}}<C_{\mathrm{surf}}$ and typically $C_{\mathrm{class}}\ll C_{\mathrm{read}}$. So the numerical savings are moderate but already real: the theorem still has five links, yet each link is checked on exposed kernels instead of full visible state surfaces, and the final readout is centralized on one canonical kernel class.
\end{remark}

\begin{remark}
This theorem is only the first coded QAM shadow of the current OP transport architecture. It does not yet include the full branch-shell routed layer, but it already captures the source/intermediate/return logic in a ZF-grounded coded form.
\end{remark}



\subsection*{QAM branch, shell, and routed support layer}
\begin{quote}\small
\noindent\textbf{Status.} \emph{Active theorem-core branch-shell extension.}\par
\noindent\textbf{Block function.} This block lifts the basic QAM transport layer to the branch/shell/routed regime and should be read as the coded counterpart of the admissible OP refinement chain.
\end{quote}

We now extend the QAM-over-ZF language from basic transport and readout to the branch/shell/routed layer. The point of this section is to translate the current admissible-branch, OP~3 shell-admissibility, and admissible routed-support architecture into a single coded QAM fragment without yet claiming a fully independent QAM foundation. For direct continuity with the compact front extraction, the present subsection should be read in the same order as the release dictionary of Remark~\ref{rem:public-shape-reading-guide}: branch admissibility first, shell admissibility second, admissible branch-shell support third, normalized routed support fourth, and only then the visible-versus-structural screening rule.

\begin{definition}[Closure-admissible branch code]
A branch code $b$ is called \emph{closure-admissible} if for all coded states $s,t$,
\[
\mathrm{QBranchRef}(b,s,t)\Longrightarrow \mathrm{QClEq}(s,t).
\]
Equivalently, branch refinement by $b$ preserves the coded closure class.
\end{definition}

\begin{definition}[Shell-admissible refinement code]
A shell/refinement code $\sigma$ is called \emph{shell-admissible} if for all coded states $s,t$,
\[
\mathrm{QShellRef}(\sigma,s,t)\Longrightarrow \mathrm{QClEq}(s,t).
\]
Equivalently, shell-sensitive or residual-sensitive refinement by $\sigma$ preserves the coded closure class.
\end{definition}

\begin{remark}
These are the direct coded QAM analogues of closure-admissible branches and OP~3 shell-admissible refinements in the current Companion line.
\end{remark}

\begin{definition}[Admissible branch set and shell family]
We define the admissible branch class by
\[
\mathcal{B}^{\mathrm{QAM}}_{\mathrm{adm}}
:=
\{\,b:\mathrm{QBranch}(b)\ \wedge\ b\text{ is closure-admissible}\,\},
\]
and for each branch code $b$, the shell-admissible refinement class by
\[
\mathcal{F}^{\mathrm{QAM}}_{\mathrm{sh}}(b)
:=
\{\,\sigma:\mathrm{QShell}(\sigma)\ \wedge\ \sigma\text{ is shell-admissible on }b\,\}.
\]
\end{definition}

\begin{remark}
At this stage, the phrase ``shell-admissible on $b$'' may be read minimally as: shell refinement by $\sigma$ acts on branchwise states associated with $b$ without leaving the coded closure class.
\end{remark}

\begin{definition}[Routed support code]
A \emph{routed support code} is a set-code
\[
u=\langle 8,\Omega,w\rangle,
\]
where $\Omega$ is a coded finite or set-sized support family of admissible branch-shell pairs and $w$ is a coded weight assignment on $\Omega$.
\end{definition}

\begin{definition}[Admissible routed support]
A routed support code
\[
u=\langle 8,\Omega,w\rangle
\]
is called \emph{admissible} if every element of $\Omega$ has the form $(b,\sigma)$ with
\[
b\in \mathcal{B}^{\mathrm{QAM}}_{\mathrm{adm}}
\qquad\text{and}\qquad
\sigma\in \mathcal{F}^{\mathrm{QAM}}_{\mathrm{sh}}(b).
\]
\end{definition}

\begin{remark}
Thus admissible routed support is the coded QAM version of the admissible branch-shell support sector in the current OP architecture.
\end{remark}

\begin{definition}[Normalized routed support]
An admissible routed support code
\[
u=\langle 8,\Omega,w\rangle
\]
is called \emph{normalized} if its coded weight assignment satisfies
\[
\sum_{(b,\sigma)\in\Omega} w(b,\sigma)=1
\]
in the intended decoded numerical target.
\end{definition}

\begin{remark}
The normalization condition is stated semiformally at this stage. In the fully coded version, the sum relation would itself be represented by a suitable arithmetic coding layer over ZF.
\end{remark}

\begin{definition}[Branch-shell coded return realization]
Let $b$ be a closure-admissible branch code and let $\sigma$ be a shell-admissible refinement code on $b$. A \emph{branch-shell coded return realization} is a coded state
\[
t_{b,\sigma}
\]
such that:
\begin{enumerate}
\item $t_{b,\sigma}$ is reached from the coded source state through admissible branch refinement by $b$;
\item the OP~3-level shell refinement by $\sigma$ is shell-admissible;
\item the resulting return-side coded state is read through a closure-compatible coded readout.
\end{enumerate}
\end{definition}

\begin{proposition}[QAM branch admissibility principle]\label{prop:qam-branch-admissibility-principle}
Let $b$ be a closure-admissible branch code. If
\[
\mathrm{QBranchRef}(b,s,t),
\]
then
\[
\mathrm{QClEq}(s,t).
\]
Consequently, every closure-compatible readout identifies the closure-relevant image of $s$ and $t$ after passage through the associated class-factor map.
\end{proposition}

\begin{proof}
The first claim is just the definition of closure-admissibility. The readout claim follows from closure-compatibility of coded readouts: closure-equivalent coded states have the same closure-relevant image after factorization.
\end{proof}

\begin{proposition}[QAM shell admissibility principle]\label{prop:qam-shell-admissibility-principle}
Let $\sigma$ be a shell-admissible refinement code. If
\[
\mathrm{QShellRef}(\sigma,s,t),
\]
then
\[
\mathrm{QClEq}(s,t).
\]
Consequently, shell-sensitive refinement does not alter closure-relevant readout content as long as shell-admissibility is preserved.
\end{proposition}

\begin{proof}
This is the direct content of shell-admissibility together with closure-compatible readout factorization.
\end{proof}

\begin{theorem}[QAM admissible branch-shell theorem]\label{thm:qam-branch-shell}
Let $s_{\mathrm{in}}$ be a coded source state, let $b$ be a closure-admissible branch code, and let $\sigma$ be a shell-admissible refinement code on $b$. Assume there exists a branch-shell coded return realization
\[
t_{b,\sigma}
\]
associated with $(b,\sigma)$, and let $r$ be a closure-compatible coded readout. Then
\[
\mathrm{QClEq}(s_{\mathrm{in}},t_{b,\sigma}).
\]
Hence every closure-compatible readout of $s_{\mathrm{in}}$ and $t_{b,\sigma}$ agrees after passage through its associated closure-factor map.
\end{theorem}

\begin{proof}
By closure-admissibility of $b$, branch refinement preserves the coded closure class. By shell-admissibility of $\sigma$, the shell-sensitive refinement preserves the same coded closure class at the OP~3 layer. Hence the final coded return realization $t_{b,\sigma}$ remains closure-equivalent to the coded source state $s_{\mathrm{in}}$. Closure-compatible readout then identifies the closure-relevant content of both states after factorization.
\end{proof}



\subsection*{Shell certificates and the first OP~3/QAM closure guide}
\begin{quote}\small
\noindent\textbf{Status.} \emph{Active closure-guide layer.}\par
\noindent\textbf{Block function.} This block isolates the certificate vocabulary that lets readers track exactly when OP~3 shell behavior is still closure-preserving and when a local failure first becomes visible.
\end{quote}
The next narrow OP~3-oriented extension of the public-shape QAM line is not a broad new tensor layer, but a certificate layer. The goal is to say explicitly when a branch/shell refinement is still the same closure-bearing object, what counts as the first local failure of that claim, and how the deferred kernel bridge should later read that same situation in one compressed normal form. Mixed tensors remain deferred here: they are not introduced unless an actual OP~3 proof step later forces them.

\begin{center}
\fbox{\begin{minipage}{0.93\textwidth}
\small
\textbf{QAM reading key for the OP~3 closure block.}\\[0.2em]
\textbf{Objects:} coded source state $s_{\mathrm{in}}$, closure-admissible branch code $b$, shell refinement code $\sigma$, intermediate branch state $s_b$, branch-shell return state $t_{b,\sigma}$.\\
\textbf{Allowed move:} pass from $s_{\mathrm{in}}$ to $s_b$ by admissible branch refinement and from $s_b$ to $t_{b,\sigma}$ by shell refinement without leaving the closure class.\\
\textbf{Invariant:} the closure-bearing class, and later its kernel/normal-form representative when the deferred bridge is activated.\\
\textbf{Defect signal:} a broken branch clause, a broken shell clause, or a failure of closure-compatible readout factorization.\\
\textbf{Plain-language reading:} OP~3 may change the visible shell, but it must not change the closure-bearing core.
\end{minipage}}
\end{center}

\begin{definition}[Shell admissibility certificate]\label{def:qam-shell-certificate}
Let $s_{\mathrm{in}}$ be a coded source state, let $b$ be a closure-admissible branch code, let $\sigma$ be a shell refinement code, and let $t_{b,\sigma}$ be a branch-shell coded return realization. A \emph{shell admissibility certificate} for $(s_{\mathrm{in}},b,\sigma,t_{b,\sigma})$ consists of coded data witnessing the following four clauses:
\begin{enumerate}[leftmargin=2em]
  \item there exists a coded branch state $s_b$ with $\mathrm{QBranchRef}(b,s_{\mathrm{in}},s_b)$;
  \item the branch code $b$ is closure-admissible;
  \item one has $\mathrm{QShellRef}(\sigma,s_b,t_{b,\sigma})$ and the refinement code $\sigma$ is shell-admissible on $b$;
  \item the return state $t_{b,\sigma}$ is read through a closure-compatible coded readout.
\end{enumerate}
\end{definition}

\begin{definition}[Shell defect witness]\label{def:qam-shell-defect-witness}
A \emph{shell defect witness} for $(s_{\mathrm{in}},b,\sigma,t_{b,\sigma})$ is coded data showing failure of at least one clause of Definition~\ref{def:qam-shell-certificate}. Thus a shell defect witness may record failure of admissible branch passage, failure of shell admissibility, or failure of closure-compatible readout factorization.
\end{definition}

\begin{lemma}[Certificate persistence of closure class]\label{lem:qam-shell-certificate-closure}
Assume that $(s_{\mathrm{in}},b,\sigma,t_{b,\sigma})$ carries a shell admissibility certificate. Then
\[
\mathrm{QClEq}(s_{\mathrm{in}},t_{b,\sigma}).
\]
Hence every closure-compatible coded readout identifies the same closure-relevant image on $s_{\mathrm{in}}$ and $t_{b,\sigma}$ after passage through the associated class-factor map.
\end{lemma}

\begin{proof}
By clause (1) of Definition~\ref{def:qam-shell-certificate}, there exists $s_b$ with
\[
\mathrm{QBranchRef}(b,s_{\mathrm{in}},s_b).
\]
By clause (2), the branch code $b$ is closure-admissible, so Proposition~\ref{prop:qam-branch-admissibility-principle} yields
\[
\mathrm{QClEq}(s_{\mathrm{in}},s_b).
\]
By clause (3), one has
\[
\mathrm{QShellRef}(\sigma,s_b,t_{b,\sigma}),
\]
and the refinement code $\sigma$ is shell-admissible on $b$. Proposition~\ref{prop:qam-shell-admissibility-principle} therefore gives
\[
\mathrm{QClEq}(s_b,t_{b,\sigma}).
\]
Transitivity of the coded closure relation yields
\[
\mathrm{QClEq}(s_{\mathrm{in}},t_{b,\sigma}).
\]
The final sentence is then the usual closure-compatible readout factorization claim.
\end{proof}

\begin{proposition}[Certified kernel persistence for OP~3 shell passage]\label{prop:qam-shell-certificate-kernel}
Assume the direct kernel constructor of Definition~\ref{def:qam-direct-kernel-constructor} and the bridge conditions of Theorem~\ref{thm:qam-direct-kernel-constructor}. Let
\[
s_{\mathrm{in}}=\langle 0,a_{\mathrm{in}}\rangle
\]
be the coded source state, and let
\[
t_{b,\sigma}=\langle 0,a_{b,\sigma}\rangle
\]
be a branch-shell coded return realization. If $(s_{\mathrm{in}},b,\sigma,t_{b,\sigma})$ carries a shell admissibility certificate, then
\[
\operatorname{nf}(a_{b,\sigma})=\operatorname{nf}(a_{\mathrm{in}}).
\]
In particular, certified OP~3 shell passage preserves the kernel normal form of the coded source state.
\end{proposition}

\begin{proof}
By Lemma~\ref{lem:qam-shell-certificate-closure}, one has
\[
\mathrm{QClEq}(s_{\mathrm{in}},t_{b,\sigma}).
\]
The deferred kernel bridge then reflects closure equivalence into equality of canonical kernel normal forms. Since
\[
s_{\mathrm{in}}=\langle 0,a_{\mathrm{in}}\rangle,
\qquad
 t_{b,\sigma}=\langle 0,a_{b,\sigma}\rangle,
\]
this yields
\[
\operatorname{nf}(a_{b,\sigma})=\operatorname{nf}(a_{\mathrm{in}}).
\]
\end{proof}

\begin{theorem}[Shell-stable readout factorization]\label{thm:qam-shell-stable-readout}
Assume the direct kernel constructor of Definition~\ref{def:qam-direct-kernel-constructor} and the bridge conditions of Theorem~\ref{thm:qam-direct-kernel-constructor}. Let
\[
s_{\mathrm{in}}=\langle 0,a_{\mathrm{in}}\rangle,
\qquad
 t_{b,\sigma}=\langle 0,a_{b,\sigma}\rangle,
\]
let $r$ be a closure-compatible coded readout, and assume that $(s_{\mathrm{in}},b,\sigma,t_{b,\sigma})$ carries a shell admissibility certificate. Then for every pair of coded observable values $w_{\mathrm{in}},w_{b,\sigma}$ with
\[
\mathrm{QRead}(r,s_{\mathrm{in}},w_{\mathrm{in}}),
\qquad
\mathrm{QRead}(r,t_{b,\sigma},w_{b,\sigma}),
\]
one has
\[
J_r(w_{\mathrm{in}})
=
\widetilde{\mathcal R}_r(\operatorname{nf}(a_{\mathrm{in}}))
=
\widetilde{\mathcal R}_r(\operatorname{nf}(a_{b,\sigma}))
=
J_r(w_{b,\sigma}).
\]
In particular, every closure-relevant coded readout of a certified OP~3 shell passage factors through one shell-stable kernel normal form.
\end{theorem}

\begin{proof}
By Proposition~\ref{prop:qam-shell-certificate-kernel}, a shell admissibility certificate implies
\[
\operatorname{nf}(a_{b,\sigma})=\operatorname{nf}(a_{\mathrm{in}}).
\]
By clause~(4) of Theorem~\ref{thm:qam-direct-kernel-constructor}, every closure-compatible coded readout $r$ admits a factor map
\[
\widetilde{\mathcal R}_r:\operatorname{nf}(\mathcal A)\to \operatorname{Im}(J_r)
\]
such that whenever $\mathrm{QRead}(r,\langle 0,a\rangle,w)$,
\[
J_r(w)=\widetilde{\mathcal R}_r(\operatorname{nf}(a)).
\]
Applying this first to $s_{\mathrm{in}}=\langle 0,a_{\mathrm{in}}\rangle$ and then to $t_{b,\sigma}=\langle 0,a_{b,\sigma}\rangle$ yields
\[
J_r(w_{\mathrm{in}})=\widetilde{\mathcal R}_r(\operatorname{nf}(a_{\mathrm{in}})),
\qquad
J_r(w_{b,\sigma})=\widetilde{\mathcal R}_r(\operatorname{nf}(a_{b,\sigma})).
\]
Since the two kernel normal forms agree, the displayed chain follows. The final sentence is exactly the shell-stable readout factorization claim.
\end{proof}

\begin{proposition}[First converse shell-certificate criterion]\label{prop:qam-shell-converse-criterion}
Assume the direct kernel constructor of Definition~\ref{def:qam-direct-kernel-constructor} and the bridge conditions of Theorem~\ref{thm:qam-direct-kernel-constructor}. Let
\[
s_{\mathrm{in}}=\langle 0,a_{\mathrm{in}}\rangle,
\qquad
 t_{b,\sigma}=\langle 0,a_{b,\sigma}\rangle,
\]
and suppose that the following data are given:
\begin{enumerate}[leftmargin=2em]
  \item there exists a coded branch state $s_b$ with $\mathrm{QBranchRef}(b,s_{\mathrm{in}},s_b)$;
  \item the branch code $b$ is closure-admissible;
  \item one has $\mathrm{QShellRef}(\sigma,s_b,t_{b,\sigma})$ and the refinement code $\sigma$ is shell-admissible on $b$;
  \item there exists a closure-compatible coded readout on the return state $t_{b,\sigma}$;
  \item the kernel normal forms agree, $\operatorname{nf}(a_{b,\sigma})=\operatorname{nf}(a_{\mathrm{in}})$.
\end{enumerate}
Then $(s_{\mathrm{in}},b,\sigma,t_{b,\sigma})$ carries a shell admissibility certificate. Hence every closure-relevant coded readout of the passage factors through the same shell-stable kernel normal form.
\end{proposition}

\begin{proof}
The first four assumptions are exactly the four clauses required in Definition~\ref{def:qam-shell-certificate}, so they furnish a shell admissibility certificate for $(s_{\mathrm{in}},b,\sigma,t_{b,\sigma})$. The final sentence then follows from Theorem~\ref{thm:qam-shell-stable-readout}, using the additional kernel-persistence hypothesis only as the converse bridge datum that motivates the criterion.
\end{proof}

\begin{lemma}[First OP~3/QAM shell commuting-square lemma]\label{lem:qam-op3-shell-commuting-square}
Assume the faithful translation theorem for the admissible OP fragment, Theorem~\ref{thm:qam-translation}, together with the OP~3 shell theorem, Theorem~\ref{thm:op3-shell}. Let
\[
\mathfrak T_{\mathrm{OP3}}
\]
denote the admissible OP~3 shell-refinement passage on the OP side, and let
\[
\Sigma_{b,\sigma}
\]
denote the corresponding coded branch-shell refinement on the QAM side. Then the source-to-refinement square commutes on closure-relevant content in the following sense: whenever an admissible OP~3 shell passage is translated into coded QAM data,
\[
\tau\bigl(\mathfrak T_{\mathrm{OP3}}(x)\bigr)
\sim_{\mathrm{cl}}
\Sigma_{b,\sigma}\bigl(\tau(x)\bigr),
\]
where $\sim_{\mathrm{cl}}$ means equality after passage to the common closure-bearing class and, under the deferred kernel bridge, equality of the induced kernel normal form. In particular, OP~3 refinement and QAM shell refinement are not two competing shell operations; they are two readings of the same closure-preserving shell step.
\end{lemma}

\begin{proof}
By Theorem~\ref{thm:qam-translation}, the admissible OP fragment transports into coded QAM data without changing the source/return closure logic. By Theorem~\ref{thm:op3-shell}, the OP~3 shell passage is admissible exactly at the stage where the OP side isolates shell refinement without changing the closure-bearing content. The translated branch stage therefore produces a coded branch state $s_b$, and the translated shell stage produces a coded return state $t_{b,\sigma}$ with the same closure-relevant content as the OP-side refined state. Lemma~\ref{lem:qam-shell-certificate-closure} then identifies the translated source state and the coded branch-shell return state inside one coded closure class, while Proposition~\ref{prop:qam-shell-certificate-kernel} upgrades this to equality of kernel normal forms once the deferred bridge is activated. That is exactly the asserted commuting-square statement.
\end{proof}

\begin{theorem}[First OP~3/QAM shell correspondence theorem]\label{thm:qam-op3-shell-correspondence}
Assume the faithful translation theorem for the admissible OP fragment, Theorem~\ref{thm:qam-translation}, together with the OP~3 shell theorem, Theorem~\ref{thm:op3-shell}. Then every translated OP~3-admissible shell passage determines a shell admissibility certificate in QAM, and every such certificate recovers the same closure-relevant shell content as the corresponding OP~3 step. Together with Proposition~\ref{prop:qam-shell-converse-criterion} and Lemma~\ref{lem:qam-op3-shell-commuting-square}, the QAM shell-certificate layer may already be read as a first narrow converse guide for OP~3: admissible branch-shell behavior, a closure-compatible return readout, and kernel persistence are sufficient to recover the certificate package. Consequently, the QAM shell-certificate layer may be read as the first explicit coded closure guide for OP~3.
\end{theorem}

\begin{proof}
Theorem~\ref{thm:qam-translation} translates the admissible OP fragment into the coded QAM layer while preserving its source/return closure logic. Theorem~\ref{thm:op3-shell} isolates the OP~3 shell-admissible passage itself. Under that translation, the OP-side branch step provides clause (1) of Definition~\ref{def:qam-shell-certificate}, branch admissibility gives clause (2), OP~3 shell admissibility gives clause (3), and the translated closure-compatible readout gives clause (4). Hence each translated OP~3-admissible shell passage determines a shell certificate. Conversely, Lemma~\ref{lem:qam-shell-certificate-closure} and Theorem~\ref{thm:qam-shell-stable-readout} show that any such certificate preserves exactly the closure-relevant content that the OP~3 theorem already isolates on the OP side. Lemma~\ref{lem:qam-op3-shell-commuting-square} records the same fact in commuting-square form. Thus the certificate layer is not a second closure notion; it is the coded OP~3 closure guide.
\end{proof}

\begin{theorem}[First narrow shell-certificate iff package]\label{thm:qam-shell-iff-package}
Assume the direct kernel constructor of Definition~\ref{def:qam-direct-kernel-constructor} and the bridge conditions of Theorem~\ref{thm:qam-direct-kernel-constructor}. Let
\[
s_{\mathrm{in}}=\langle 0,a_{\mathrm{in}}\rangle,
\qquad
 t_{b,\sigma}=\langle 0,a_{b,\sigma}\rangle,
\]
and let $r$ be any closure-compatible coded readout on the return state. Then the following are equivalent:
\begin{enumerate}[leftmargin=2em,label=(\roman*)]
  \item $(s_{\mathrm{in}},b,\sigma,t_{b,\sigma})$ carries a shell admissibility certificate;
  \item there exists a coded branch state $s_b$ such that
  \begin{enumerate}[leftmargin=2em,label=(\alph*)]
    \item $\mathrm{QBranchRef}(b,s_{\mathrm{in}},s_b)$,
    \item the branch code $b$ is closure-admissible,
    \item $\mathrm{QShellRef}(\sigma,s_b,t_{b,\sigma})$ and the refinement code $\sigma$ is shell-admissible on $b$,
    \item the return state $t_{b,\sigma}$ admits a closure-compatible coded readout,
    \item $\operatorname{nf}(a_{b,\sigma})=\operatorname{nf}(a_{\mathrm{in}})$,
    \item every closure-relevant coded readout of the passage factors through one shell-stable kernel normal form.
  \end{enumerate}
\end{enumerate}
In particular, the present OP~3/QAM shell block already supports one narrow iff reading: a shell certificate is equivalent to the shell-stable kernel/readout package obtained from admissible branch-shell behavior together with kernel persistence.
\end{theorem}

\begin{proof}
Assume \emph{(i)}. Then clauses (a)--(d) of \emph{(ii)} are exactly the four certificate clauses from Definition~\ref{def:qam-shell-certificate}. Clause~(e) is Proposition~\ref{prop:qam-shell-certificate-kernel}, and clause~(f) is Theorem~\ref{thm:qam-shell-stable-readout}. Hence \emph{(ii)} follows.

Assume conversely \emph{(ii)}. Then clauses (a)--(d) are precisely the data required by Definition~\ref{def:qam-shell-certificate}, so $(s_{\mathrm{in}},b,\sigma,t_{b,\sigma})$ carries a shell admissibility certificate. This is exactly Proposition~\ref{prop:qam-shell-converse-criterion}. Thus \emph{(i)} follows, and the final sentence is only a repackaging of the same forward and converse implications.
\end{proof}

\begin{lemma}[First shell-obstruction lemma]\label{lem:qam-shell-obstruction}
Assume the direct kernel constructor of Definition~\ref{def:qam-direct-kernel-constructor} and the bridge conditions of Theorem~\ref{thm:qam-direct-kernel-constructor}. Let
\[
s_{\mathrm{in}}=\langle 0,a_{\mathrm{in}}\rangle,
\qquad
 t_{b,\sigma}=\langle 0,a_{b,\sigma}\rangle,
\]
and let $r$ be any closure-compatible coded readout on the return state whenever such a readout exists. Suppose that at least one clause of the shell-stable kernel/readout package from Theorem~\ref{thm:qam-shell-iff-package} fails. Then no shell admissibility certificate exists for $(s_{\mathrm{in}},b,\sigma,t_{b,\sigma})$. Equivalently, the first local OP~3 obstruction in the present QAM reading is the first failure of one of the following items:
\begin{enumerate}[leftmargin=2em,label=(\roman*)]
  \item admissible branch passage,
  \item shell-admissible branch refinement,
  \item closure-compatible return readout,
  \item kernel persistence,
  \item shell-stable factorization of closure-relevant readout through one common kernel normal form.
\end{enumerate}
In particular, a visible shell mismatch that leaves all five items intact is not yet an OP~3 obstruction in the present QAM sense.
\end{lemma}

\begin{proof}
By Theorem~\ref{thm:qam-shell-iff-package}, a shell admissibility certificate is equivalent to the shell-stable kernel/readout package. Therefore the failure of any one clause of that package prevents the certificate from existing. The displayed list simply unwraps the package into its five operative clauses. The final sentence is the contrapositive reading of the same equivalence: visible shell variation alone is not an obstruction as long as admissible branch-shell transport, closure-compatible readout, kernel persistence, and common readout factorization still hold.
\end{proof}

\begin{lemma}[First shell-obstruction stratification lemma]\label{lem:qam-shell-obstruction-stratification}
Assume the setup of Lemma~\ref{lem:qam-shell-obstruction}. If no shell admissibility certificate exists for $(s_{\mathrm{in}},b,\sigma,t_{b,\sigma})$, then at least one of the following three strata fails:
\begin{enumerate}[leftmargin=2em,label=(\arabic*)]
  \item \textbf{Branch stratum:} admissible branch passage fails.
  \item \textbf{Shell stratum:} admissible branch passage holds, but shell-admissible refinement fails.
  \item \textbf{Readout/kernel stratum:} admissible branch passage and shell-admissible refinement both hold, but closure-compatible return readout, kernel persistence, or common shell-stable readout factorization fails.
\end{enumerate}
Equivalently, every first local OP~3 obstruction in the present QAM reading is either a branch obstruction, a shell obstruction in the narrow sense, or a readout/kernel obstruction. In particular, whenever the branch and shell strata both remain intact, the first local obstruction is already forced into the readout/kernel layer.
\end{lemma}

\begin{proof}
Lemma~\ref{lem:qam-shell-obstruction} says that the first local obstruction is exactly the failure of one of five operative clauses. Group those clauses into the three strata above: clause~(i) is the branch stratum, clause~(ii) is the shell stratum, and clauses~(iii)--(v) together form the readout/kernel stratum. This yields the stated trichotomy. The final sentence is only the contrapositive reading of the same grouping.
\end{proof}

\begin{lemma}[Readout/kernel obstruction compression lemma]\label{lem:qam-readout-kernel-obstruction-compression}
Assume the setup of Lemma~\ref{lem:qam-shell-obstruction-stratification}, and assume moreover that the branch and shell strata both remain intact. Then the following are equivalent:
\begin{enumerate}[leftmargin=2em,label=(\roman*)]
  \item the first local OP~3 obstruction lies in the readout/kernel stratum;
  \item at least one of the following failures occurs:
  \begin{enumerate}[leftmargin=2em,label=(\alph*)]
    \item closure-compatible return readout fails,
    \item kernel persistence fails,
    \item common shell-stable readout factorization through one kernel normal form fails;
  \end{enumerate}
  \item there is no single shell-stable kernel/readout carrier through which all closure-relevant readout of the certified branch-shell passage can factor.
\end{enumerate}
Equivalently, once the branch and shell strata are fixed, the whole remaining OP~3 obstruction is compressed to the failure of one common readout/kernel carrier class.
\end{lemma}

\begin{proof}
Under the standing assumption that the branch and shell strata remain intact, Lemma~\ref{lem:qam-shell-obstruction-stratification} says that any remaining first local obstruction must lie in the readout/kernel stratum. By definition of that stratum, this is exactly the disjunction of failures in clauses~(a)--(c). This proves the equivalence of \emph{(i)} and \emph{(ii)}.

If a single shell-stable kernel/readout carrier exists, then closure-compatible return readout is available, kernel persistence identifies the same kernel normal form on source and return side, and Theorem~\ref{thm:qam-shell-stable-readout} gives the common factorization through that carrier. Hence \emph{(ii)} cannot occur. Conversely, if any one of clauses~(a)--(c) fails, then no single shell-stable kernel/readout carrier can support the whole closure-relevant readout package. This proves the equivalence of \emph{(ii)} and \emph{(iii)}. The final sentence is only a compressed restatement of the same equivalence.
\end{proof}

\begin{theorem}[First OP~3 shell-closure theorem in QAM form]\label{thm:qam-op3-shell-closure}
Assume the faithful translation theorem for the admissible OP fragment, Theorem~\ref{thm:qam-translation}, together with the OP~3 shell theorem, Theorem~\ref{thm:op3-shell}. Let a translated OP~3-admissible shell passage determine coded data
\[
(s_{\mathrm{in}},b,\sigma,t_{b,\sigma},r).
\]
Then the following are equivalent:
\begin{enumerate}[leftmargin=2em,label=(\roman*)]
  \item the translated passage is \emph{QAM-shell-closed in the present narrow sense}, meaning that it carries a shell admissibility certificate and recovers the same closure-relevant shell content as the corresponding OP~3 step;
  \item $(s_{\mathrm{in}},b,\sigma,t_{b,\sigma})$ carries a shell admissibility certificate;
  \item the shell-stable kernel/readout package of Theorem~\ref{thm:qam-shell-iff-package} holds.
\end{enumerate}
Moreover, under these equivalent conditions the passage admits a common shell-stable kernel/readout carrier through which all closure-relevant coded readout factors. If the equivalent conditions fail, then the first local OP~3 obstruction occurs at exactly one of the branch, shell, or readout/kernel strata from Lemma~\ref{lem:qam-shell-obstruction-stratification}; and in the third stratum the whole remaining obstruction is precisely failure of one common shell-stable carrier in the sense of Lemma~\ref{lem:qam-readout-kernel-obstruction-compression}.
\end{theorem}

\begin{proof}
By Theorem~\ref{thm:qam-op3-shell-correspondence}, every translated OP~3-admissible shell passage determines a shell admissibility certificate and recovers the same closure-relevant shell content as the corresponding OP~3 step. Thus \emph{(i)} and \emph{(ii)} are equivalent in the present narrow reading.

The equivalence of \emph{(ii)} and \emph{(iii)} is exactly Theorem~\ref{thm:qam-shell-iff-package}. Under these equivalent conditions, Theorem~\ref{thm:qam-shell-stable-readout} gives factorization of every closure-relevant coded readout through the same shell-stable kernel normal form, hence through one common shell-stable kernel/readout carrier.

If the equivalent conditions fail, then by Lemma~\ref{lem:qam-shell-obstruction} the first local OP~3 obstruction is the first failure of the shell-stable kernel/readout package. Lemma~\ref{lem:qam-shell-obstruction-stratification} refines that failure into the branch, shell, and readout/kernel strata. Finally, Lemma~\ref{lem:qam-readout-kernel-obstruction-compression} compresses the whole third stratum to failure of one common shell-stable carrier. This proves the final sentence.
\end{proof}

\begin{corollary}[Shell defect localization is the first local OP~3 failure]\label{cor:qam-shell-defect-localization}
Within the present OP~3/QAM reading, visible packet mismatch, phase-lock deviation, or shell-sensitive surface variation is not yet by itself a shell defect. A genuine local OP~3 defect begins only when a shell defect witness in the sense of Definition~\ref{def:qam-shell-defect-witness} exists. Equivalently, by Theorem~\ref{thm:qam-op3-shell-closure}, the first local OP~3 failure is the first failure of the shell-stable kernel/readout package, stratified into branch, shell, and readout/kernel level and compressed in the third stratum to failure of one common shell-stable carrier. In particular, the first local failure is failure of a shell certificate clause, not failure of the entire OP~1--OP~5 closure chain.
\end{corollary}

\begin{remark}[How this block should be used in the next OP~3 proof pass]
The present certificate block is intentionally narrow. It does not yet force a full equivalence between the OP and QAM formalisms, and it does not yet activate mixed tensors. Its current role is to give the next proof pass a fixed reading order: establish a shell certificate, pass to kernel persistence where the deferred bridge is available, use Theorem~\ref{thm:qam-shell-stable-readout} to factor readout through the shell-stable class, use Proposition~\ref{prop:qam-shell-converse-criterion} for the first narrow return direction, use Lemma~\ref{lem:qam-op3-shell-commuting-square} to read OP~3 refinement and QAM shell refinement as one commuting closure step, use Theorem~\ref{thm:qam-shell-iff-package} as the first closed iff package, use Lemma~\ref{lem:qam-shell-obstruction} to identify the first genuinely local obstruction, use Lemma~\ref{lem:qam-shell-obstruction-stratification} to decide whether that first obstruction sits at the branch, shell, or readout/kernel layer, use Lemma~\ref{lem:qam-readout-kernel-obstruction-compression} to compress the third stratum to failure of one common shell-stable carrier, and then use Theorem~\ref{thm:qam-op3-shell-closure} as the first bundled OP~3 shell-closure statement in public-shape QAM form before asking whether any further tensor layer is genuinely needed.
\end{remark}


\subsection*{Current extracted OP~3/QAM tool layer}
\begin{remark}[Purpose of the extracted OP~3/QAM tools]\label{rem:qam-op3-tool-purpose}
The certificate and obstruction chain above is no longer meant to be reread from first principles in every later proof pass. Its immediate use is as a compact tool layer: once a later argument can cite one of the extracted tools below, it may treat the corresponding shell passage as already compressed to the right closure-bearing datum and continue at the next structural interface.
\end{remark}

\begin{description}[leftmargin=2.2em,style=nextline]
  \item[Shell-certificate tool.]\label{tool:qam-shell-certificate}
  If one can exhibit an admissible branch passage, an admissible shell refinement, and a closure-compatible return readout, then the passage may be treated as a certified OP~3 shell step in the present narrow QAM sense. This is the entry tool for the whole block.

  \item[Kernel-persistence tool.]\label{tool:qam-kernel-persistence}
  Once a shell certificate is in place and the deferred kernel bridge is available, the visible shell variation may be replaced by equality of one kernel normal form. From that point onward the proof may work kernel-first instead of shell-first.

  \item[Shell-stable readout tool.]\label{tool:qam-shell-stable-readout}
  Once the shell certificate and kernel bridge are active, every closure-relevant readout may be pushed through one common shell-stable carrier. Later readout comparisons therefore no longer need to unfold the whole branch/shell passage.

  \item[Converse reconstruction tool.]\label{tool:qam-shell-converse}
  If admissible branch passage, admissible shell refinement, closure-compatible return readout, and kernel persistence are already known from another source, then the shell certificate may be reconstructed rather than reproved from scratch.

  \item[Commuting-square tool.]\label{tool:qam-shell-commuting-square}
  OP~3 shell refinement on the OP side and branch/shell refinement on the QAM side may be read as one closure-preserving square. This tool licenses switching between OP and QAM viewpoints without reopening the shell argument.

  \item[Obstruction-stratification tool.]\label{tool:qam-shell-obstruction-stratification}
  If the shell-certificate package fails, the first local obstruction should be classified immediately as branch, shell, or readout/kernel level. This prevents visible shell mismatch from being over-read as a full OP~3 failure.

  \item[Carrier-compression tool.]\label{tool:qam-shell-carrier-compression}
  Once branch and shell remain intact, every remaining local obstruction may be compressed to failure of one common shell-stable carrier. This is the shortest diagnostic form of the third obstruction stratum.

  \item[Bundled OP~3 closure tool.]\label{tool:qam-op3-shell-closure}
  Theorem~\ref{thm:qam-op3-shell-closure} may now be cited as the compact public-shape closure package for OP~3 in QAM form: certificate, kernel/readout package, stratified obstruction, and common carrier are read as one bundled result.
\end{description}

\begin{remark}[How to cite the extracted OP~3/QAM tools]\label{rem:qam-op3-tool-usage}
A later argument should normally cite the smallest extracted tool that already matches its interface. Use the shell-certificate tool when the issue is admissible passage, the kernel-persistence or shell-stable readout tool when the issue is invariance under visible shell variation, the converse reconstruction tool when the shell certificate is to be recovered from already known transport/readout data, the commuting-square tool when the proof needs to switch languages, the obstruction-stratification or carrier-compression tool when the issue is failure localization, and the bundled OP~3 closure tool when the whole narrow closure package is intended at once.
\end{remark}


\subsection*{First OP~3 proof run through the extracted QAM tool layer}

\begin{remark}[Reading key for the first OP~3 proof run]\label{rem:qam-op3-proof-run-reading-key}
In the present narrow public-shape QAM reading, the first proof run should be read as a deliberately compressed structural pass rather than as a maximal theorem expansion. The input is a translated OP~3 shell passage. The allowed move is passage from visible shell data to the extracted tool layer. The invariant is the common closure-bearing kernel/readout carrier whenever the certificate package is active. The defect signal is the first obstruction stratum from Lemma~\ref{lem:qam-shell-obstruction-stratification}. In plain language: the proof run should decide as early as possible whether the shell step is already certified, and if not, where the first local failure sits.
\end{remark}

\begin{theorem}[First OP~3 proof run through extracted QAM form]\label{thm:qam-op3-proof-run}
Assume the faithful translation theorem for the admissible OP fragment, Theorem~\ref{thm:qam-translation}, the OP~3 shell theorem, Theorem~\ref{thm:op3-shell}, the direct kernel constructor of Definition~\ref{def:qam-direct-kernel-constructor}, and the bridge conditions of Theorem~\ref{thm:qam-direct-kernel-constructor}. Let
\[
\tau(x)=s_{\mathrm{in}}=\langle 0,a_{\mathrm{in}}\rangle
\]
be the coded QAM source state obtained from an admissible OP~3 source configuration $x$, and let the translated branch/shell passage be encoded by a branch code $b$, a refinement code $\sigma$, a coded branch state $s_b$, and a coded return state
\[
t_{b,\sigma}=\langle 0,a_{b,\sigma}\rangle.
\]
Then the first OP~3 proof run through extracted QAM form has exactly two outcomes:
\begin{enumerate}[leftmargin=2em,label=(\roman*)]
  \item either the translated passage activates the bundled OP~3 closure tool, so that the OP~3 shell step is certified in QAM form, all closure-relevant coded readout factors through one common shell-stable carrier, and the passage may be continued kernel-first;
  \item or the bundled closure package fails, in which case the first local obstruction is immediately localized by the obstruction-stratification tool at exactly one of the branch, shell, or readout/kernel strata, with the whole third stratum compressed to failure of one common shell-stable carrier.
\end{enumerate}
In particular, the extracted tool layer already supports a genuine proof run: it converts translated OP~3 shell data into either a compressed closure package or a minimal obstruction diagnosis, without reopening the whole OP/QAM comparison at each step.
\end{theorem}

\begin{proof}
Translate the admissible OP fragment by Theorem~\ref{thm:qam-translation}. By the OP~3 shell theorem, Theorem~\ref{thm:op3-shell}, the translated branch and shell passage is exactly the stage at which the OP side isolates admissible shell refinement. Use the \hyperref[tool:qam-shell-commuting-square]{commuting-square tool} to identify the translated OP~3 shell step with the coded branch/shell refinement step on the QAM side.

Now ask whether the \hyperref[tool:qam-shell-certificate]{shell-certificate tool} is activated. If it is, then by the \hyperref[tool:qam-kernel-persistence]{kernel-persistence tool} the visible shell passage may be replaced by equality of one kernel normal form, and by the \hyperref[tool:qam-shell-stable-readout]{shell-stable readout tool} every closure-relevant coded readout factors through one common shell-stable carrier. The bundled \hyperref[tool:qam-op3-shell-closure]{OP~3 closure tool} then yields outcome \emph{(i)}.

If the shell-certificate package is not activated, apply the \hyperref[tool:qam-shell-obstruction-stratification]{obstruction-stratification tool}. This localizes the first failure immediately at exactly one of the branch, shell, or readout/kernel levels. If the failure lies in the third stratum, the \hyperref[tool:qam-shell-carrier-compression]{carrier-compression tool} compresses the remaining failure to absence of one common shell-stable carrier. This is outcome \emph{(ii)}.

The two outcomes are exhaustive because Theorem~\ref{thm:qam-op3-shell-closure} already packages the positive direction as the bundled shell-closure theorem, while Lemma~\ref{lem:qam-shell-obstruction-stratification} together with Lemma~\ref{lem:qam-readout-kernel-obstruction-compression} packages the negative direction. Hence the extracted tool layer already yields a genuine first OP~3 proof run in QAM form.
\end{proof}

\begin{corollary}[First operational use of the extracted OP~3/QAM tools]\label{cor:qam-op3-proof-run-use}
A later OP~3 argument no longer needs to reopen the whole branch/shell analysis from scratch whenever a translated shell passage appears. It is enough to test whether the shell-certificate tool is active. If yes, cite the bundled OP~3 closure tool and continue kernel-first; if not, cite the obstruction-stratification tool and, in the third stratum, the carrier-compression tool. Thus the proof interface has already been shortened to a certified-pass / localized-obstruction dichotomy.
\end{corollary}

\begin{proof}
This is only a restatement of Theorem~\ref{thm:qam-op3-proof-run} in proof-interface language.
\end{proof}


\begin{remark}[Status separation for the present OP~3/QAM closure package]\label{rem:qam-op3-status-separation}
The present block should now be read in two layers with different evidential status. The \emph{active public-shape result} is the certificate-level OP~3/QAM closure package: Definitions~\ref{def:qam-shell-certificate} and~\ref{def:qam-shell-defect-witness}, Lemma~\ref{lem:qam-shell-certificate-closure}, Theorem~\ref{thm:qam-shell-iff-package}, Theorem~\ref{thm:qam-op3-shell-closure}, and Corollary~\ref{cor:qam-op3-proof-run-use}. These statements already yield a closed certified-pass / localized-obstruction interface for translated OP~3 shell passages without leaving the active tuple surface. By contrast, the stronger kernel-first language --- in particular Proposition~\ref{prop:qam-shell-certificate-kernel}, Theorem~\ref{thm:qam-shell-stable-readout}, and the carrier-compression reading of the third obstruction stratum --- is a \emph{conditional strengthening}: it becomes authoritative only under the deferred bridge hypotheses of Definition~\ref{def:qam-kernel-bridge-datum} and Theorem~\ref{thm:qam-direct-kernel-constructor}. Accordingly, later use should distinguish carefully between active certificate-level closure and conditional kernel/readout compression.
\end{remark}

\begin{proposition}[Citation and use rule for the completed OP~3/QAM block]\label{prop:qam-op3-citation-use-rule}
Let a later argument involve a translated OP~3 shell passage. Then the following citation discipline is valid.
\begin{enumerate}[leftmargin=2em,label=(\roman*)]
  \item If the argument needs only a certified-pass / localized-obstruction decision at the active tuple level, it may cite Theorem~\ref{thm:qam-op3-shell-closure}, Theorem~\ref{thm:qam-op3-proof-run}, and Corollary~\ref{cor:qam-op3-proof-run-use} without invoking any kernel-first hypothesis.
  \item If the argument further needs invariance under visible shell variation in kernel normal-form language or one-carrier readout factorization, then it may additionally cite Proposition~\ref{prop:qam-shell-certificate-kernel}, Theorem~\ref{thm:qam-shell-stable-readout}, and Lemma~\ref{lem:qam-readout-kernel-obstruction-compression}, but only after the deferred bridge hypotheses have been assumed or discharged.
  \item If neither the certificate package nor the bridge package is available, the argument must fall back to the earlier branch/shell/readout layer and may not treat visible shell data as a closure theorem by themselves.
\end{enumerate}
\end{proposition}

\begin{proof}
Clause~(i) is exactly the active reading packaged by Theorem~\ref{thm:qam-op3-shell-closure}, Theorem~\ref{thm:qam-op3-proof-run}, and Corollary~\ref{cor:qam-op3-proof-run-use}. Clause~(ii) is the additional kernel-first strengthening isolated earlier in Proposition~\ref{prop:qam-shell-certificate-kernel}, Theorem~\ref{thm:qam-shell-stable-readout}, and Lemma~\ref{lem:qam-readout-kernel-obstruction-compression}; those statements explicitly depend on the deferred bridge hypotheses. Clause~(iii) is the corresponding fallback rule: absent both packages, one must continue on the earlier active branch/shell/readout surface.
\end{proof}

\begin{corollary}[Operational checklist for the completed OP~3/QAM block]\label{cor:qam-op3-operational-checklist}
A translated OP~3 shell passage should now be tested in the following order.
\begin{enumerate}[leftmargin=2em]
  \item Check admissible branch passage and branch admissibility.
  \item Check shell refinement and shell admissibility on the chosen branch.
  \item Check existence of a closure-compatible return readout.
  \item If all three checks succeed, activate the shell-certificate package and cite the bundled OP~3 closure tool.
  \item Only if the deferred bridge hypotheses are additionally available, pass to kernel persistence and shell-stable readout factorization.
  \item If one of the certificate checks fails, localize the first obstruction immediately at branch, shell, or readout/kernel level instead of reopening the whole OP~1--OP~5 chain.
\end{enumerate}
In particular, the completed block already provides one stable certified interface: \emph{certificate first, kernel compression second, obstruction localization immediately upon failure}.
\end{corollary}

\begin{remark}[Diagnostic handoff rule for visible shell data]\label{rem:qam-op3-diagnostic-handoff}
Visible packet mismatch, phase-lock deviation, shell-sensitive surface variation, and analogous shell diagnostics remain \emph{screening-layer data} unless and until they witness failure of one clause of the shell-certificate package or, under the bridge hypotheses, failure of one common shell-stable carrier. Thus the present block fixes a handoff rule from pictures and visible diagnostics to theorematic closure language: diagnostics may guide attention to the relevant branch, shell, or readout/kernel stratum, but they do not create a second closure axiom beside the completed OP~3/QAM certificate package. This is the intended local counterpart of the earlier visible/structural separation principle and the correct interface toward the later OP~3 to OP~5 transfer line.
\end{remark}

\begin{proposition}[Prospective kernel-first reduction of the branch-shell theorem]\label{prop:qam-kernel-first-branch-shell}
Assume the direct kernel constructor of Definition~\ref{def:qam-direct-kernel-constructor} and the bridge conditions of Theorem~\ref{thm:qam-direct-kernel-constructor}. Let
\[
s_{\mathrm{in}}=\langle 0,a_{\mathrm{in}}\rangle
\]
be the coded source state, let $b$ be a closure-admissible branch code, let $\sigma$ be a shell-admissible refinement code on $b$, and let
\[
t_{b,\sigma}=\langle 0,a_{b,\sigma}\rangle
\]
be an associated branch-shell coded return realization. Then the conclusion of Theorem~\ref{thm:qam-branch-shell} may be read kernel-first as follows:
\begin{enumerate}
\item one has
\[
\operatorname{nf}(a_{b,\sigma})=\operatorname{nf}(a_{\mathrm{in}}),
\]
\item hence every closure-compatible readout of $s_{\mathrm{in}}$ and $t_{b,\sigma}$ factors through the single canonical kernel class determined by $a_{\mathrm{in}}$,
\item therefore branch choice and shell refinement may be treated as closure-preserving already at exposed-kernel level before returning to the visible QAM state surface.
\end{enumerate}
\end{proposition}

\begin{proof}
Theorem~\ref{thm:qam-branch-shell} gives
\[
\mathrm{QClEq}(s_{\mathrm{in}},t_{b,\sigma}).
\]
Since both states lie on the present state surface, we may write
\[
s_{\mathrm{in}}=\langle 0,a_{\mathrm{in}}\rangle,
\qquad
 t_{b,\sigma}=\langle 0,a_{b,\sigma}\rangle.
\]
By the direct kernel constructor,
\[
\operatorname{ker}_{\mathrm{dir}}(s_{\mathrm{in}})=a_{\mathrm{in}},
\qquad
\operatorname{ker}_{\mathrm{dir}}(t_{b,\sigma})=a_{b,\sigma}.
\]
The bridge clause for closure reflection therefore yields
\[
\operatorname{nf}(a_{b,\sigma})=\operatorname{nf}(a_{\mathrm{in}}).
\]
Thus branch admissibility together with shell admissibility already identifies the same canonical kernel class before any further visible-surface comparison is carried out. By the factorization clause of the bridge datum, every closure-compatible coded readout depends only on that common kernel class. Hence the branch-shell theorem may be re-read as a kernel-first statement and only afterwards lifted back to the visible state surface.
\end{proof}

\begin{remark}[Axiom audit after the kernel-first branch-shell reduction]\label{rem:qam-kernel-first-branch-shell-audit}
Proposition~\ref{prop:qam-kernel-first-branch-shell} does \emph{not} alter the active proof status of Theorem~\ref{thm:qam-branch-shell}. The active theorem still rests only on closure-admissible branch refinement, shell-admissible OP~3 refinement, and closure-compatible coded readout. The kernel-first proposition adds force only under the deferred payload-level bridge hypotheses isolated in Definition~\ref{def:qam-kernel-bridge-datum} and Theorem~\ref{thm:qam-direct-kernel-constructor}. Thus no new hidden axiom has entered the live branch-shell theorem; only a prospective kernel-level compression of its proof surface has been recorded.
\end{remark}

\begin{remark}[How much proof load the kernel-first branch-shell reading removes]\label{rem:qam-kernel-first-branch-shell-savings}
At branch-shell level the cardinality factor is just $1$: there is one visible source surface
\[
s_{\mathrm{in}}=\langle 0,a_{\mathrm{in}}\rangle
\]
and one visible return surface
\[
t_{b,\sigma}=\langle 0,a_{b,\sigma}\rangle.
\]
So the gain is not a reduction in the number of cases, but a reduction in the \emph{size of the compared objects}. On the visible state surface the proof compares two full coded states and then passes a closure-compatible readout across that state-level identification. In the deferred kernel-first reading, that work is replaced by the single normal-form check
\[
\operatorname{nf}(a_{b,\sigma})=\operatorname{nf}(a_{\mathrm{in}})
\]
followed by one readout through the common kernel class. If one measures a visible-state comparison by cost $C_{\mathrm{surf}}$ and a kernel normal-form comparison by cost $C_{\mathrm{ker}}$, then the conceptual load changes from roughly
\[
C_{\mathrm{surf}} + C_{\mathrm{read}}
\]
to
\[
C_{\mathrm{ker}} + C_{\mathrm{class}},
\]
with intended gain $C_{\mathrm{ker}}<C_{\mathrm{surf}}$ and typically $C_{\mathrm{class}}\ll C_{\mathrm{read}}$. So the numerical savings here are modest compared with the routed case, but methodologically this is the crucial base step: it is the first place where a full visible return state can be replaced by its exposed payload kernel without changing the closure conclusion.
\end{remark}

\begin{remark}[Theorem-level unpacking of the compact chain, I]
Theorem~\ref{thm:qam-branch-shell} is the first theorem-level unpacking of the compact front statement: it isolates the middle part of the release chain in which closure-preserving source-return transport passes through branch choice and shell refinement without losing closure-relevant readout content.
\end{remark}

\begin{corollary}[QAM branch-shell support class]
Define the coded admissible branch-shell support class by
\[
\mathcal{S}^{\mathrm{QAM}}_{\mathrm{adm}}
:=
\{\, (b,\sigma): b\in \mathcal{B}^{\mathrm{QAM}}_{\mathrm{adm}}
\ \wedge\
\sigma\in \mathcal{F}^{\mathrm{QAM}}_{\mathrm{sh}}(b)\,\}.
\]
Then every coded return realization supported on
\[
\mathcal{S}^{\mathrm{QAM}}_{\mathrm{adm}}
\]
is closure-preserving at the coded QAM level.
\end{corollary}

\begin{remark}
This is the coded QAM counterpart of the admissible branch-shell support sector in the present OP architecture.
\end{remark}

\begin{definition}[QAM branch-shell routed return]
Let
\[
u=\langle 8,\Omega,w\rangle
\]
be an admissible routed support code with
\[
\Omega\subseteq \mathcal{S}^{\mathrm{QAM}}_{\mathrm{adm}}.
\]
A \emph{QAM branch-shell routed return} is a coded return aggregation associated with the support family $\Omega$ and the weight assignment $w$.
\end{definition}

\begin{theorem}[QAM routed admissible support theorem]\label{thm:qam-routed-support}
Let $u=\langle 8,\Omega,w\rangle$ be a normalized admissible routed support code, and let $r$ be a closure-compatible coded readout. Then the closure-relevant image of the routed return determined by $u$ agrees with the closure-relevant image of the coded source state.
\end{theorem}

\begin{proof}
Every element $(b,\sigma)\in\Omega$ lies in the admissible branch-shell support class, so each associated coded return realization is closure-equivalent to the coded source state by the previous theorem. Since the routed support is normalized and the readout is closure-compatible, the routed aggregation preserves the same closure-relevant image after passage through the factor map.
\end{proof}

\begin{proposition}[Prospective kernel-first reduction of the routed support theorem]\label{prop:qam-kernel-first-routed}
Assume the direct kernel constructor of Definition~\ref{def:qam-direct-kernel-constructor} and the bridge conditions of Theorem~\ref{thm:qam-direct-kernel-constructor}. Let
\[
s_{\mathrm{in}}=\langle 0,a_{\mathrm{in}}\rangle
\]
be the coded source state, and let
\[
u=\langle 8,\Omega,w\rangle
\]
be a normalized admissible routed support code. For each $(b,\sigma)\in\Omega$, let
\[
t_{b,\sigma}=\langle 0,a_{b,\sigma}\rangle
\]
be an associated branch-shell coded return realization. Then the conclusion of Theorem~\ref{thm:qam-routed-support} may be read kernel-first as follows:
\begin{enumerate}
\item for every $(b,\sigma)\in\Omega$ one has
\[
\operatorname{nf}(a_{b,\sigma})=\operatorname{nf}(a_{\mathrm{in}}),
\]
\item hence every closure-compatible routed readout factors through the single canonical kernel class of $a_{\mathrm{in}}$,
\item therefore the routed aggregation determined by $u$ has the same closure-relevant image as the coded source state without needing to compare the full family of visible state surfaces directly.
\end{enumerate}
\end{proposition}

\begin{proof}
For each $(b,\sigma)\in\Omega$, Theorem~\ref{thm:qam-branch-shell} gives
\[
\mathrm{QClEq}(s_{\mathrm{in}},t_{b,\sigma}).
\]
Since both states lie on the present state surface, we may write
\[
s_{\mathrm{in}}=\langle 0,a_{\mathrm{in}}\rangle,
\qquad
 t_{b,\sigma}=\langle 0,a_{b,\sigma}\rangle.
\]
By the direct kernel constructor, the kernel representatives are just the exposed payload kernels:
\[
\operatorname{ker}_{\mathrm{dir}}(s_{\mathrm{in}})=a_{\mathrm{in}},
\qquad
\operatorname{ker}_{\mathrm{dir}}(t_{b,\sigma})=a_{b,\sigma}.
\]
The bridge clause for closure reflection therefore yields
\[
\operatorname{nf}(a_{b,\sigma})=\operatorname{nf}(a_{\mathrm{in}})
\]
for every $(b,\sigma)\in\Omega$. Thus the entire admissible routed support family lies in one canonical kernel class. By the factorization clause of the bridge datum, every closure-compatible coded readout depends only on that kernel class. Hence the routed aggregation has the same closure-relevant image as the coded source state. The normalization of the routed weights matters only after this kernel reduction, so the proof burden is shifted from a family of full visible state surfaces to one canonical kernel class.
\end{proof}

\begin{remark}[Why the routed support theorem is the best stress test]
Among the currently marked candidates, Theorem~\ref{thm:qam-routed-support} is the heaviest set-theoretic test case for the deferred zero-kernel language: it simultaneously carries a support family $\Omega$, a coded weight assignment $w$, branch admissibility, shell admissibility, routed aggregation, and closure-compatible readout. Exactly there the kernel-first reading is strongest, because the visible family of state surfaces is replaced by one canonical kernel class before the routed aggregation is evaluated.
\end{remark}


\begin{remark}[Axiom audit after the kernel-first routed reduction]\label{rem:qam-kernel-first-routed-audit}
Proposition~\ref{prop:qam-kernel-first-routed} does \emph{not} strengthen the active QAM axiomatics. Its extra force comes only from deferred hypotheses already isolated earlier: the direct kernel constructor of Definition~\ref{def:qam-direct-kernel-constructor} and the four payload-level bridge conditions of Theorem~\ref{thm:qam-direct-kernel-constructor}. In particular, the active theorem chain still proves Theorem~\ref{thm:qam-routed-support} without mentioning kernels at all. The new proposition merely says that \emph{if} those deferred bridge conditions are later discharged, then the routed-support proof may be re-read kernel-first. Thus no silent new axiom has entered the then-present public line; only a prospective compression interface has been exhibited.
\end{remark}

\begin{remark}[How much comparison load the kernel-first reading removes]\label{rem:qam-kernel-first-routed-savings}
Write $N:=|\Omega|$ for the size of the routed support family. On the visible state surface, the routed-support proof conceptually tracks $N$ branch-shell return states
\[
 t_{b,\sigma}=\langle 0,a_{b,\sigma}\rangle
\]
and compares each of them with the coded source state at full state-surface level before the routed aggregation is passed through a closure-compatible readout. In the deferred kernel-first reading, those $N$ state-surface comparisons are replaced by the $N$ normal-form checks
\[
\operatorname{nf}(a_{b,\sigma})=\operatorname{nf}(a_{\mathrm{in}}),
\qquad (b,\sigma)\in\Omega,
\]
followed by a \emph{single} readout through the common canonical kernel class of $a_{\mathrm{in}}$.
Consequently the proof burden drops from ``an $N$-member family of full visible state surfaces plus routed readout lifting'' to ``one canonical kernel class together with $N$ membership/equality checks into that class.''
If one measures a full visible-state comparison by cost $C_{\mathrm{surf}}$ and a kernel normal-form comparison by cost $C_{\mathrm{ker}}$, then the conceptual comparison load changes from roughly
\[
N\,C_{\mathrm{surf}} + C_{\mathrm{route}}
\]
to
\[
N\,C_{\mathrm{ker}} + C_{\mathrm{class}},
\]
where typically $C_{\mathrm{class}}\ll C_{\mathrm{route}}$ and the intended gain is precisely that $C_{\mathrm{ker}}<C_{\mathrm{surf}}$. In proof-shape terms, nothing miraculous happens to the factor $N$ itself, but the \emph{object size} being compared is reduced from full state surfaces to exposed kernels, and the routed readout is evaluated once on a single kernel class instead of being conceptually re-established across the whole visible family.
\end{remark}

\begin{remark}[Theorem-level unpacking of the compact chain, II]
Theorem~\ref{thm:qam-routed-support} is the routed continuation of Theorem~\ref{thm:qam-branch-shell}. Together they unpack the compressed public-shape chain
\[
s_{\mathrm{in}} \longrightarrow (b,\sigma) \longrightarrow u \longrightarrow r
\]
of Theorem~\ref{thm:compact-public-shape-preservation} into theorem-level branch/shell preservation followed by theorem-level routed aggregation.
\end{remark}

\begin{remark}[Staged proof compression pattern under the deferred zero-kernel language]\label{rem:qam-staged-kernel-compression}
The three prospective kernel-first reductions now show a genuinely staged compression pattern.
At the first stage, Proposition~\ref{prop:qam-kernel-first-transport} replaces the visible six-state transport chain by equality of six exposed kernel normal forms, so that admissible transport is already read as kernel-class preservation before any branch or shell data enter.
At the second stage, Proposition~\ref{prop:qam-kernel-first-branch-shell} replaces one visible return-surface comparison
\[
\langle 0,a_{b,\sigma}\rangle \sim \langle 0,a_{\mathrm{in}}\rangle
\]
by the single kernel normal-form comparison
\[
\operatorname{nf}(a_{b,\sigma})=\operatorname{nf}(a_{\mathrm{in}}).
\]
At the third stage, Proposition~\ref{prop:qam-kernel-first-routed} repeats that reduced comparison only at the level of exposed kernels across the whole routed family $\Omega$ and then evaluates the closure-compatible readout once on the common kernel class. Thus the deferred zero-kernel language first shrinks the transport chain to a common kernel class, then shrinks each branch-shell return to one exposed kernel comparison, and only afterwards turns the routed family into a class-membership problem rather than a full visible-surface transport problem. This is exactly why transport is the first chain-level compression step, branch-shell the first local return-surface compression step, and routed support the first true family-level compression step.
\end{remark}

\begin{definition}[Prospective activation package for the deferred zero-kernel language]\label{def:qam-zero-kernel-activation-package}
A \emph{zero-kernel activation package} for the present QAM surface consists of the following data and properties:
\begin{enumerate}[label=(\roman*)]
\item a total assignment on active state surfaces
\[
\ker:QState \to \mathcal A, \qquad \ker(\langle 0,a\rangle)=a,
\]
compatible with Definition~\ref{def:qam-visible-payload};
\item kernel admissibility, meaning that every active exposed payload kernel lies in the admissible kernel class \(\mathcal A\);
\item kernel normal-form reflection of closure, meaning that closure-equivalent active states have the same kernel normal form;
\item kernel-level factorization of every closure-compatible readout through the corresponding kernel normal form;
\item compatibility of the routed aggregation step with passage to the common kernel class.
\end{enumerate}
\end{definition}

\begin{theorem}[Activation criterion for kernel-first proof compression on the present QAM surface]\label{thm:qam-zero-kernel-activation-criterion}
Assume a zero-kernel activation package in the sense of Definition~\ref{def:qam-zero-kernel-activation-package}. Then the deferred zero-kernel language may be used as a genuine proof-compression layer for the present QAM surface in the following precise sense:
\begin{enumerate}[label=(\alph*)]
\item every active state-surface argument may be reduced to a kernel normal-form argument without loss of closure-relevant content;
\item the first transport theorem may be read kernel-first by proving preservation of one common kernel normal-form class along the admissible transport chain;
\item the branch-shell theorem may be read kernel-first by proving normal-form preservation of the exposed return kernels;
\item the routed support theorem may be read kernel-first by proving membership of all exposed routed kernels in a single common normal-form class and evaluating the closure-compatible readout once on that class.
\end{enumerate}
In particular, the deferred zero-kernel language ceases to be merely heuristic at exactly the point where closure, readout, and routed aggregation all factor through kernel normal forms. The criterion identifies \emph{what} must hold; the later staged-roadmap remarks then specify \emph{in which order} those requirements should be established on the present QAM surface.
\end{theorem}

\begin{proof}
This is just the joint content of Theorem~\ref{thm:qam-kernel-bridge}, Theorem~\ref{thm:qam-direct-kernel-constructor}, Proposition~\ref{prop:qam-kernel-first-transport}, Proposition~\ref{prop:qam-kernel-first-branch-shell}, and Proposition~\ref{prop:qam-kernel-first-routed}, collected as a single activation statement. Once closure-equivalent active surfaces have the same kernel normal form and every relevant readout factors through that normal form, every transport, branch-shell, and routed-support comparison may be carried out on kernels first and only lifted back to the visible state surface at the end.
\end{proof}

\begin{remark}[Criterion-to-roadmap reading rule]\label{rem:qam-zero-kernel-criterion-roadmap}
The activation theorem above should be read together with the staged roadmap below. The theorem gives the logical threshold for activation, whereas the roadmap orders the internal work needed to reach that threshold:
\[
\begin{aligned}
&\text{transport-class preservation}\\[0.15em]
&\Longrightarrow\;\text{local branch-shell kernel reduction}\\[0.15em]
&\Longrightarrow\;\text{routed family compression}.
\end{aligned}
\]
In this combined reading, the criterion says \emph{when} kernel-first proof compression becomes legitimate, while the roadmap says \emph{how} the present document should approach that point without introducing premature changes at visible tuple level.
\end{remark}

\begin{remark}[Practical migration trigger for a later binary-tag version]\label{rem:qam-zero-kernel-migration-trigger}
The activation criterion above, read together with the later criterion-to-roadmap rule, also isolates the only later migration trigger that would make a binary or power-of-two tag scheme genuinely worthwhile. Such a migration would become structurally justified only once the present state surface is already being used kernel-first in proofs, so that the visible tag layer has become a comparatively thin presentation layer while closure and readout are actually controlled at kernel-normal-form level. Before that point, a binary tag migration would mainly rename visible surfaces; after that point, it could act as a real simplification of the presentation layer.
\end{remark}

\begin{remark}[What still has to be shown before activation]\label{rem:qam-zero-kernel-activation-gap}
At the present stage, the deferred zero-kernel language still stops short of activation for one simple reason: the document now contains the bridge data, the direct kernel candidate, and three staged stress tests, but it does not yet contain a fully internal proof that the active closure relation and every closure-compatible readout already factor through the exposed payload normal forms on the entire current QAM surface. The current status is therefore stronger than a mere metaphor but weaker than an active rewrite rule.
\end{remark}


\begin{remark}[Prospective staged activation roadmap for the zero-kernel language]\label{rem:qam-zero-kernel-roadmap}
The current document already suggests a natural staged migration order for any later activation of the deferred zero-kernel language.
Stage~I is the \emph{transport stage}: one first proves that admissible source-return transport preserves a single common kernel normal-form class all along the visible active chain.
Stage~II is the \emph{local return stage}: one then proves that every admissible branch-shell return surface may be replaced by one exposed kernel comparison against the incoming kernel class.
Stage~III is the \emph{family stage}: one finally proves that the whole routed family collapses to membership in that common kernel class, so that closure-compatible readout is evaluated once at class level rather than re-established on each visible routed surface.
Only after these three stages are internally available would a later binary or power-of-two visible presentation acquire real proof-theoretic leverage, because only then would the visible tuple tags have become a thin presentation shell over an already kernel-first proof flow.
\end{remark}

\begin{remark}[Why this staged order matters]\label{rem:qam-zero-kernel-roadmap-why}
The order above is not cosmetic. If routed family compression were attempted before transport-class preservation and local branch-shell reduction were available, then the family step would still depend on full visible surfaces and would not yet be a genuine kernel-first simplification. The staged order therefore identifies the first point at which the deferred zero-kernel language would cease to be merely suggestive and would begin to remove real set-theoretic comparison load in a controlled way. Together with the criterion-to-roadmap rule, it turns the activation block from a static admissibility test into a concrete migration sequence.
\end{remark}

\begin{remark}[Compressed synthesis of the deferred zero-kernel path]\label{rem:qam-zero-kernel-synthesis}
The deferred zero-kernel material above should now be read as one compact path rather than as many separate addenda. The present visible state surface remains unchanged. A later kernel language would first expose the exposed payload kernel of a coded state as a candidate kernel, then compare closure only through kernel normal forms, then compress transport, local branch-shell return, and finally routed families in that order, and only afterwards consider whether a binary visible presentation would simplify the already kernel-first proof flow. In that reading, the deferred block does not introduce a second theory beside the current QAM package. It isolates one possible later proof-compression pipeline inside the same public branch/shell/routed line.
\end{remark}

\begin{remark}[Reading guide for the deferred block]\label{rem:qam-zero-kernel-reading-guide}
A compact reader may keep only the following chain in view:
\[
\begin{aligned}
&\text{visible state surface}
\;\Longrightarrow\;
\text{exposed payload kernel}
\;\Longrightarrow\;
\text{kernel normal form}
\\[0.2em]
&\Longrightarrow\;
\text{transport/local/family compression}
\;\Longrightarrow\;
\text{possible later binary presentation}.
\end{aligned}
\]
The detailed deferred notes are included only to justify each arrow in this chain and to mark exactly which theorem-level bridges would still need to be proved before activation.
\end{remark}


\begin{remark}[Dependency map for the deferred zero-kernel path]\label{rem:qam-zero-kernel-dependency-map}
The deferred block now has a strict internal dependency order.
First, Definitions~\ref{def:qam-zero-kernel}--\ref{def:qam-kernel-key} fix the later kernel vocabulary: carriers, filters, primitive zero-objects, recursive reuse, aspect projections, normal forms, and internal search keys.
Second, the axiom-coherence and current-limitation remarks state why this added vocabulary does not yet compete with the active tuple layer.
Third, Definition~\ref{def:qam-kernel-bridge-datum}, Theorem~\ref{thm:qam-kernel-bridge}, Definition~\ref{def:qam-visible-payload}, and Theorem~\ref{thm:qam-direct-kernel-constructor} isolate the exact interface through which active coded states could later be re-read kernel-first.
Fourth, Propositions~\ref{prop:qam-kernel-first-transport}, \ref{prop:qam-kernel-first-branch-shell}, and \ref{prop:qam-kernel-first-routed} test that interface on the three currently most relevant proof stages.
Finally, Definition~\ref{def:qam-zero-kernel-activation-package}, Theorem~\ref{thm:qam-zero-kernel-activation-criterion}, and the later activation-roadmap, synthesis, and reading-guide remarks convert those tests into an activation criterion and a migration order.
In particular, any later binary visible presentation should depend on this whole chain and should not precede it.
\end{remark}

\begin{remark}[Local admissibility test for a later kernel-first rewrite]\label{rem:qam-zero-kernel-local-test}
The deferred block can now also be read as a local test for whether a specific visible tuple-surface proof step is mature enough for later kernel-first rewriting. Given a proof step whose visible data consist of active coded states
\[
s_i=\langle 0,a_i\rangle,
\]
one should ask four questions in order.
\begin{enumerate}
\item \emph{Exposure:} do all exposed payload kernels $a_i$ already lie in the intended admissible kernel domain $\mathcal A$?
\item \emph{Reduction target:} can the visible conclusion be reformulated as equality, preservation, or common membership of the corresponding kernel normal forms $\operatorname{nf}(a_i)$?
\item \emph{Factorization:} does every closure-relevant readout or routed aggregation used in that proof step depend only on the resulting kernel normal-form class?
\item \emph{Lift back:} once the kernel-class statement is proved, does the original visible state-surface conclusion follow again without adding new data?
\end{enumerate}
If the answer is yes in all four places, then the proof step is a genuine candidate for later kernel-first compression. If one of the four answers is still no, the corresponding proof step should remain on the visible state surface for the time being.
\end{remark}

\begin{remark}[Fallback rule for current proofs]\label{rem:qam-zero-kernel-fallback}
Failure of the local admissibility test is not a defect of the present active proof line. It simply means that the relevant argument should continue to cite the active visible-surface theorem rather than its deferred kernel-first restatement. This fallback rule prevents premature theorem drift: the deferred block may prepare later compression, but it does not lower the evidential status required for current public proofs.
\end{remark}

\begin{remark}[Compressed rewrite schema]\label{rem:qam-zero-kernel-rewrite-schema}
Whenever the later local admissibility test applies, the deferred zero-kernel path suggests one uniform later rewrite schema:
\[
\begin{aligned}
&\text{visible tuple step}
\;\Longrightarrow\;
\text{expose payloads}
\;\Longrightarrow\;
\text{pass to }\operatorname{nf}
\\[0.2em]
&\Longrightarrow\;
\text{prove closure/readout at class level}
\;\Longrightarrow\;
\text{lift back to the visible surface}.
\end{aligned}
\]
This does not yet activate any new theorem in the present document. It simply records, in one line, the exact proof pattern that the three staged stress tests have already isolated separately for transport, local branch/shell return, and routed families.
\end{remark}

\begin{remark}
This is the coded QAM analogue of the combined branch-shell transport theorem together with normalized routed admissible preservation.
\end{remark}

\begin{definition}[QAM screening layer]
A coded diagnostic layer is called a \emph{screening layer} if it provides visible preselection information about admissible branch-shell support without itself constituting an independent closure proof.
\end{definition}

\begin{remark}
In this sense, packet-, phase-lock-, and epsilon-type coded diagnostics should be treated in QAM as screening-layer data rather than as primary closure criteria. Their structural relevance begins only when they track loss of branch admissibility or loss of shell admissibility.
\end{remark}

\begin{corollary}[QAM visible/structural separation]\label{cor:qam-visible-structural}
Visible coded mismatch in the screening layer does not by itself imply failure of closure at the QAM level. Closure failure begins only when coded branch admissibility or coded shell admissibility is lost.
\end{corollary}

\begin{remark}
This is the QAM-over-ZF shadow of the visible-mismatch-versus-structural-mismatch principle already established in the current Companion line.
\end{remark}


\subsection*{QAM tensor admissibility and tensor closure rules}

We now strengthen the tensor-coded part of the QAM language. At the present stage, tensor composition is still not treated as a full tensor algebra. Instead, it is introduced as a controlled structural constructor equipped with admissibility, closure, transport, and readout compatibility rules.

\begin{definition}[Tensor-coded formation]
If $p$ and $q$ are QAM-coded objects, then
\[
p\otimes q := \langle 7,p,q\rangle
\]
is their tensor-coded composite. We write
\[
\mathrm{QTensor}(p\otimes q)
\]
for the statement that the code has the tensor tag.
\end{definition}

\begin{definition}[Tensor formation by kind]
The following tensor formations are allowed in principle:
\begin{enumerate}
\item state-state tensors:
\[
s_1\otimes s_2;
\]
\item layer-layer tensors:
\[
\ell_1\otimes \ell_2;
\]
\item operator-operator tensors:
\[
o_1\otimes o_2;
\]
\item readout-readout tensors:
\[
r_1\otimes r_2;
\]
\item branch-branch tensors:
\[
b_1\otimes b_2;
\]
\item shell-shell tensors:
\[
\sigma_1\otimes \sigma_2.
\]
\end{enumerate}
A tensor-coded composite is called \emph{well-typed} if both entries belong to one of the admissible kind-pairs listed above.
\end{definition}

\begin{remark}
At this stage, mixed tensors between different kinds are not taken as primitive. They may be introduced later through derived coding conventions if needed.
\end{remark}

\begin{definition}[Tensor admissibility]
For QAM-coded objects $x,y$, we write
\[
\mathrm{QTensorAdm}(x,y)
\]
iff:
\begin{enumerate}
\item $x\otimes y$ is well-typed;
\item the tensor-coded composite preserves the intended closure-bearing structural regime;
\item the tensor-coded composite is admitted as a legitimate QAM object of its corresponding kind.
\end{enumerate}
\end{definition}

\begin{remark}
Thus tensor admissibility is not automatic from mere code formation. It is a structural compatibility condition.
\end{remark}

\begin{definition}[Tensor closure compatibility]
For coded states $s_1,s_2,t_1,t_2$, we write
\[
\mathrm{QTensorCl}(s_1,s_2;t_1,t_2)
\]
iff
\[
\mathrm{QClEq}(s_1,t_1)
\qquad\text{and}\qquad
\mathrm{QClEq}(s_2,t_2)
\]
and the tensor composites
\[
s_1\otimes s_2
\qquad\text{and}\qquad
t_1\otimes t_2
\]
are tensor-admissible.
\end{definition}

\begin{axiom}[Tensor closure propagation schema]
If
\[
\mathrm{QTensorCl}(s_1,s_2;t_1,t_2),
\]
then
\[
\mathrm{QClEq}(s_1\otimes s_2,\ t_1\otimes t_2).
\]
\end{axiom}

\begin{remark}
This is the first basic tensor closure principle: closure-equivalent components yield closure-equivalent tensor composites whenever tensor admissibility holds.
\end{remark}

\begin{definition}[Tensor-coded transport]
Let
\[
o_1=\langle 2,\ell_1,\ell_1',G_1\rangle,
\qquad
o_2=\langle 2,\ell_2,\ell_2',G_2\rangle
\]
be coded operators. Their tensor-coded transport is
\[
o_1\otimes o_2.
\]
We say that
\[
\mathrm{QTensorTransport}(o_1,o_2;s_1,s_2;t_1,t_2)
\]
holds iff:
\begin{enumerate}
\item $\mathrm{QAdmTransport}(o_1,s_1,t_1)$;
\item $\mathrm{QAdmTransport}(o_2,s_2,t_2)$;
\item $\mathrm{QTensorAdm}(s_1,s_2)$ and $\mathrm{QTensorAdm}(t_1,t_2)$;
\item $\mathrm{QTensorAdm}(o_1,o_2)$.
\end{enumerate}
\end{definition}

\begin{axiom}[Tensor transport admissibility schema]
If
\[
\mathrm{QTensorTransport}(o_1,o_2;s_1,s_2;t_1,t_2),
\]
then
\[
\mathrm{QAdmTransport}(o_1\otimes o_2,\ s_1\otimes s_2,\ t_1\otimes t_2).
\]
\end{axiom}

\begin{remark}
This is the tensorial lift of admissible transport: admissible componentwise transport induces admissible tensor transport whenever the tensor compatibility conditions are satisfied.
\end{remark}

\begin{definition}[Tensor-coded readout]
Let
\[
r_1=\langle 3,\ell_1,W_1,R_1\rangle,
\qquad
r_2=\langle 3,\ell_2,W_2,R_2\rangle
\]
be coded readouts. Their tensor-coded readout is
\[
r_1\otimes r_2.
\]
We say that $r_1\otimes r_2$ is \emph{tensor closure-compatible} if:
\begin{enumerate}
\item $r_1$ and $r_2$ are closure-compatible;
\item $\mathrm{QTensorAdm}(r_1,r_2)$;
\item the closure-relevant readout factorization of $r_1\otimes r_2$ depends only on the tensor closure class of the source state.
\end{enumerate}
\end{definition}

\begin{remark}
This is the tensor analogue of closure-compatible readout. It allows tensor-composite visible outputs to differ at raw level while preserving closure content after factorization.
\end{remark}

\begin{definition}[Tensor branch admissibility]
Let $b_1,b_2$ be closure-admissible branch codes. Their tensor composite
\[
b_1\otimes b_2
\]
is called \emph{tensor branch-admissible} if:
\begin{enumerate}
\item $\mathrm{QTensorAdm}(b_1,b_2)$;
\item componentwise branch refinement preserves closure class;
\item the induced tensor branch refinement preserves closure class at the tensor-composite level.
\end{enumerate}
\end{definition}

\begin{definition}[Tensor shell admissibility]
Let $\sigma_1,\sigma_2$ be shell-admissible refinement codes. Their tensor composite
\[
\sigma_1\otimes \sigma_2
\]
is called \emph{tensor shell-admissible} if:
\begin{enumerate}
\item $\mathrm{QTensorAdm}(\sigma_1,\sigma_2)$;
\item componentwise shell refinement preserves closure class;
\item the induced tensor shell refinement preserves closure class at the tensor-composite level.
\end{enumerate}
\end{definition}

\begin{remark}
Tensor branch admissibility and tensor shell admissibility are the coded QAM precursors of coupled branch/shell transport across multiple structural sectors.
\end{remark}

\begin{proposition}[Tensor closure preservation principle]
Let $s_1,s_2,t_1,t_2$ be coded states. Assume
\[
\mathrm{QClEq}(s_1,t_1),
\qquad
\mathrm{QClEq}(s_2,t_2),
\qquad
\mathrm{QTensorAdm}(s_1,s_2),
\qquad
\mathrm{QTensorAdm}(t_1,t_2).
\]
Then
\[
\mathrm{QClEq}(s_1\otimes s_2,\ t_1\otimes t_2).
\]
\end{proposition}

\begin{proof}
This is exactly the tensor closure propagation schema.
\end{proof}

\begin{proposition}[Tensor transport preservation principle]
Assume
\[
\mathrm{QTensorTransport}(o_1,o_2;s_1,s_2;t_1,t_2).
\]
Then
\[
\mathrm{QClEq}(s_1\otimes s_2,\ t_1\otimes t_2).
\]
\end{proposition}

\begin{proof}
By tensor transport admissibility,
\[
\mathrm{QAdmTransport}(o_1\otimes o_2,\ s_1\otimes s_2,\ t_1\otimes t_2).
\]
Admissible transport implies closure preservation, hence
\[
\mathrm{QClEq}(s_1\otimes s_2,\ t_1\otimes t_2).
\]
\end{proof}

\begin{theorem}[First tensor QAM transport theorem]
Let
\[
o_1,o_2
\]
be coded operators,
\[
s_1,s_2
\]
coded source states,
and
\[
t_1,t_2
\]
coded transported states. Assume
\[
\mathrm{QTensorTransport}(o_1,o_2;s_1,s_2;t_1,t_2),
\]
and let
\[
r_1\otimes r_2
\]
be a tensor closure-compatible readout. Then the tensor-composite source state and the tensor-composite transported state lie in the same coded closure class:
\[
\mathrm{QClEq}(s_1\otimes s_2,\ t_1\otimes t_2),
\]
and therefore their closure-relevant tensor readout content agrees after factorization through the tensor class-target map.
\end{theorem}

\begin{proof}
By the previous proposition,
\[
\mathrm{QClEq}(s_1\otimes s_2,\ t_1\otimes t_2).
\]
Since the tensor readout is closure-compatible at tensor level, closure-equivalent tensor-composite states have the same closure-relevant image after factorization. This proves the claim.
\end{proof}

\begin{corollary}[Tensor branch-shell admissible support]
If
\[
b_1\otimes b_2
\]
is tensor branch-admissible and
\[
\sigma_1\otimes \sigma_2
\]
is tensor shell-admissible, then admissible branch-shell support is preserved at tensor-composite level, provided the relevant tensor admissibility conditions hold.
\end{corollary}

\begin{remark}
This is the first indication that the current branch-shell architecture of the Companion can be lifted into a tensor-coded QAM language without changing its closure-first logic. In the internal reading order, this tensor layer remains subordinate to the earlier branch/shell admissibility and translation layers.
\end{remark}

\begin{remark}[Current status of tensor QAM]
At the present stage, the tensor rules do not yet define a full monoidal calculus. They only establish the first closure-preserving, transport-preserving, and readout-compatible tensor schemata needed for a structure-first QAM development over ZF.
\end{remark}


\subsection*{QAM translation theorem for the admissible OP fragment}

We now make explicit that the present OP/branch/shell architecture may be translated into the coded QAM-over-ZF language without changing its closure-first logic. In v3.16.0 this translation block should be read as the theorem-level mirror of the front extraction theorem (Theorem~\ref{thm:compact-public-shape-preservation}): the front block states the compressed public-shape chain, while the present subsection states why that same chain is a faithful coded image of the already stabilized OP architecture. At the current stage, this is not yet claimed as a full equivalence of entire formal systems. It is claimed only as a faithful translation on the admissible fragment of the OP transport architecture. Accordingly, the present subsection is ordered to match the front reading guide: coded source state and closure-compatible coded readout first, closure-admissible branch code second, shell-admissible refinement code third, admissible routed support fourth, and closure-relevant readout content fifth.

\begin{definition}[Translation map from the admissible OP chain into QAM codes]
Let $\mathfrak{O}_{\mathrm{adm}}$ denote the admissible OP fragment consisting of:
\begin{enumerate}
\item source and return states of the OP1--OP5 source-return architecture;
\item intermediate admissible OP transport states;
\item admissible branch data from the branchwise refinement layer;
\item OP3 shell-admissible refinement data;
\item admissible and normalized routed support data.
\end{enumerate}
A \emph{QAM translation map} is a coding assignment
\[
\tau:\mathfrak{O}_{\mathrm{adm}}\to \mathfrak{Q}_{\mathrm{adm}}
\]
from admissible OP objects into admissible QAM-coded objects such that:
\begin{enumerate}
\item source/readout objects are sent to coded source states and coded closure-compatible readouts of the same front-facing type used in the public-shape extraction;
\item OP transport maps are sent to coded operator maps preserving the admissible transport role between coded source, branch, shell, routed-support, and readout stages;
\item admissible branches are sent to closure-admissible branch codes;
\item OP3 shell-admissible refinements are sent to shell-admissible refinement codes;
\item admissible and normalized routed support is sent to admissible and normalized routed support codes.
\end{enumerate}
\end{definition}

\begin{remark}[Reading-order convention for $\tau$]
The point of $\tau$ is not merely syntactic recoding. It must preserve the structural roles of the translated objects in the same order in which the front extraction reads them: coded source state and closure-compatible coded readout remain the coded front-face of the OP~1/OP~5 source-return pair; admissible transport remains admissible transport; branch remains branch; shell-admissible refinement remains shell-admissible refinement; routed support remains routed support; and closure-relevant readout content remains closure-relevant readout content.
\end{remark}

\begin{definition}[Closure-faithful translation on the public-shape chain]
The translation map $\tau$ is called \emph{closure-faithful} if for all admissible OP objects $X,Y\in\mathfrak{O}_{\mathrm{adm}}$,
\[
X \sim_{\mathrm{cl}} Y
\qquad\Longrightarrow\qquad
\mathrm{QClEq}(\tau(X),\tau(Y)),
\]
and whenever the OP-side relation expresses admissible preservation of the native closure-bearing mode along the source/transport/branch/shell/routed chain, the translated QAM-side relation expresses preservation of the same closure content in coded form.
\end{definition}

\begin{definition}[Readout compatibility of the translation]
The translation map $\tau$ is called \emph{readout-compatible} if OP-side closure-compatible readout invariants correspond, after translation, to QAM-coded closure-compatible readouts on the image of $\tau$, with the same closure-relevant readout content as in the public-shape chain.
\end{definition}

Then the following OP-side chain is preserved under translation into QAM-over-ZF in exactly the order used by the front public-shape reading guide:
\begin{theorem}[QAM translation theorem for the admissible OP fragment]\label{thm:qam-translation}
Assume that:
\begin{enumerate}
\item the OP transport architecture is given on its admissible fragment $\mathfrak{O}_{\mathrm{adm}}$;
\item there exists a closure-faithful and readout-compatible translation map
\[
\tau:\mathfrak{O}_{\mathrm{adm}}\to \mathfrak{Q}_{\mathrm{adm}};
\]
\item admissible OP transport, admissible branches, OP3 shell-admissible refinements, admissible routed support, and closure-compatible readout are translated according to the preceding definitions.
\end{enumerate}
Then the following OP-side chain is preserved under translation into QAM-over-ZF in exactly the order used by the front public-shape reading guide:
\[
\begin{aligned}
\text{coded source state / OP1--OP5 readout face}
&\longrightarrow \text{closure-admissible branch code}\\
&\longrightarrow \text{shell-admissible refinement code}\\
&\longrightarrow \text{admissible routed support code}\\
&\longrightarrow \text{closure-relevant coded readout content}.
\end{aligned}
\]
Equivalently, after naming the translated objects in public-shape order, one obtains the same compressed chain from Theorem~\ref{thm:compact-public-shape-preservation}:
\[
\begin{aligned}
s_{\mathrm{in}}
&\longrightarrow (b,\sigma)\\
&\longrightarrow u\\
&\longrightarrow r.
\end{aligned}
\]
More precisely:

\begin{enumerate}
\item \textbf{OP closure preservation.}  
If the OP1--OP5 source-return pair is closed at invariant level, then its QAM-coded image is closure-preserving at coded level.

\item \textbf{Intermediate transport preservation.}  
If the intermediate OP layers preserve the native closure-bearing mode, then their QAM-coded images preserve the coded closure class.

\item \textbf{Branch admissibility preservation.}  
If a branch is closure-admissible on the OP side, then its image under $\tau$ is a closure-admissible branch code on the QAM side.

\item \textbf{Shell admissibility preservation.}  
If an OP3 refinement is shell-admissible on the OP side, then its image under $\tau$ is a shell-admissible refinement code on the QAM side.

\item \textbf{Routed admissible support preservation.}  
If a routed support sector is admissible and normalized on the OP side, then its image under $\tau$ is an admissible and normalized routed support code on the QAM side, and its closure-relevant content is preserved.

\item \textbf{Front-facing readout compatibility.}  
If the translated routed support sector is evaluated by a closure-compatible OP readout, then the resulting coded readout is exactly of the public-shape type used in Theorem~\ref{thm:compact-public-shape-preservation}; in particular it remains a coded readout of preserved closure content and not a new primitive closure test.
\end{enumerate}
\end{theorem}

\begin{proof}
Each clause follows from closure-faithfulness and readout compatibility of the translation map.

For (1), closure of the OP1--OP5 pair means preservation of the native closure-bearing mode at the source-return level. By closure-faithfulness, this becomes coded closure preservation for the translated source and return states.

For (2), invariant preservation across intermediate OP layers is, by assumption, preservation of the same closure-bearing mode through admissible transport. Again, closure-faithfulness translates this into preservation of coded closure class for the corresponding QAM-coded states and operator actions.

For (3), branch admissibility on the OP side means precisely that the branchwise refinement preserves the native closure-bearing mode. By the definition of the translation map, admissible branches are sent to branch codes that preserve coded closure class.

For (4), shell admissibility at OP3 means that the refinement does not change the native closure-bearing mode. By translation, this becomes shell-admissibility of the corresponding QAM refinement code.

For (5), admissible routed support is translated by hypothesis into admissible routed support codes. Readout compatibility ensures that closure-relevant readout content is preserved after translation, and normalization of support weights is preserved by the routed support coding layer.

For (6), readout compatibility identifies the translated readout with a coded closure-compatible readout on the image of the admissible OP chain. Hence the front-facing public-shape readout statement is not extra structure; it is the coded readout face of the same preserved OP closure content.
\end{proof}


\begin{remark}[Prospective local admissibility reading of the translation theorem]\label{rem:qam-zero-kernel-translation-local}
The later local admissibility test may already be applied, in deferred form, to Theorem~\ref{thm:qam-translation}. The result is not yet a full active kernel-first rewrite of the translation theorem, but it cleanly separates the future admissible rewrite points.
\begin{enumerate}
\item Clauses~(1)--(2) are locally admissible exactly at those translated source/transport surfaces that already lie on the present visible state surface $\langle 0,a\rangle$; there the payload can be exposed directly and compared by kernel normal form.
\item Clauses~(3)--(5) are locally admissible once the translated branch, shell, and routed-support codes are shown to inherit the same exposed-kernel discipline used in Propositions~\ref{prop:qam-kernel-first-branch-shell} and~\ref{prop:qam-kernel-first-routed}.
\item Clause~(6) becomes locally admissible precisely when translated closure-compatible readout factors through the same kernel normal-form class, so that the coded public-shape readout is recovered from class level rather than re-established on the full visible translated chain.
\end{enumerate}
Accordingly, the translation theorem is already organized in the right order for later kernel-first compression, but it becomes genuinely rewrite-ready only stage by stage, in the same order isolated by the staged roadmap and the criterion-to-roadmap rule.
\end{remark}

\begin{remark}[How much clausewise load a future kernel-first translation proof could remove]\label{rem:qam-zero-kernel-translation-savings}
At present, Theorem~\ref{thm:qam-translation} is proved clausewise on the visible translated chain, with six explicit front-facing obligations. A later kernel-first proof would not remove the need to show transport, branch, shell, routed-support, and readout preservation; however, it could compress the visible proof surface from six separate clause-level checks to two kernel-level obligations:
\[
\begin{aligned}
&\text{(i) one common kernel-normal-form preservation statement}\\[-0.05em]
&\text{\phantom{(i)} along the translated admissible chain,}\\[0.25em]
&\text{(ii) one class-level readout factorization statement}\\[-0.05em]
&\text{\phantom{(ii)} recovering the coded public-shape readout.}
\end{aligned}
\]
In that sense, the prospective savings at theorem level are best read not as a reduction in mathematical content, but as a reduction from \emph{six visible surface checks} to \emph{one kernel-class propagation argument plus one class-level readout argument}. This is exactly the same kind of shift from surface multiplicity to kernel control already seen, in more local form, in the transport, branch/shell, and routed-support stress tests.
\end{remark}

\begin{corollary}[QAM image of the current OP core]
The current OP core of the Companion --- including source-return closure, intermediate transport, admissible branches, OP3 shell admissibility, admissible routed support, and closure-compatible readout --- admits a faithful coded realization inside QAM-over-ZF.
\end{corollary}

\begin{remark}[Front theorem and translation theorem as mirrored statements]
Theorem~\ref{thm:compact-public-shape-preservation} states the compressed public release chain in coded language. Theorem~\ref{thm:qam-translation} states that the same coded chain is the faithful image of the stabilized OP chain. Read together with Theorem~\ref{thm:qam-branch-shell}, Theorem~\ref{thm:qam-routed-support}, and Corollary~\ref{cor:qam-visible-structural}, the front block and the later QAM block are therefore intended to read almost as two projections of the same closure statement: first in compressed public-shape form, then in theorem-level translation form. In particular, the front reading order
\[
s_{\mathrm{in}}\longrightarrow (b,\sigma)\longrightarrow u\longrightarrow r
\]
is now the same order in which $\tau$, the closure-faithfulness condition, the readout-compatibility condition, and the partial decoding map $\delta$ are meant to be understood.
\end{remark}

\begin{remark}
This corollary is the first formal justification for introducing QAM as a structure-first coded language over ZF. It shows that the current OP architecture does not need to be abandoned or replaced in order to enter the QAM setting; it can be carried into that setting by faithful translation.
\end{remark}

\begin{definition}[Partial decoding map]
\label{def:partial-decoding-map}
A \emph{partial decoding map} is a map
\[
\delta:\mathfrak{Q}_{\mathrm{adm}}\dashrightarrow \mathfrak{O}_{\mathrm{adm}}
\]
defined on the admissible QAM fragment such that for translated admissible OP objects,
\[
\delta(\tau(X))=X.
\]
\end{definition}

\begin{remark}[How $\delta$ should be read]
At the present stage, $\delta$ is only required on the image of the admissible OP fragment. No full decoding of arbitrary QAM-coded expressions is assumed yet. When it exists, $\delta$ should be read as the reverse reading map for the same public-shape chain: it sends coded source/readout, branch, shell, and routed-support objects back into the already stabilized OP architecture without introducing a second closure criterion.
\end{remark}

\begin{proposition}[Round-trip property on the admissible fragment]\label{prop:round-trip-property-admissible-fragment}
If a partial decoding map $\delta$ exists on the image of $\tau$, then
\[
\delta\circ\tau = \mathrm{id}_{\mathfrak{O}_{\mathrm{adm}}}.
\]
Hence the admissible OP fragment is recoverable from its translated QAM image.
\end{proposition}

\begin{proof}
This is immediate from the definition of the partial decoding map.
\end{proof}

\begin{remark}
The round-trip property is intentionally weaker than full two-way interpretability. It says only that the currently relevant admissible OP fragment can be encoded into QAM over ZF and then read back without structural loss.
\end{remark}

\begin{corollary}[Tensor-compatible translation]
If the admissible OP fragment contains tensorially paired sectors and the translation map respects tensor-coded formation, then admissible tensor closure and admissible tensor transport are preserved under translation as well.
\end{corollary}

\begin{remark}[Current status of the translation theorem]
The present theorem should be read as a conservative first step. It does not yet prove that all future QAM constructions are already reducible to the current OP architecture, nor that the entire QAM language has been fully decoded back into OP form. What it proves is that the current admissible OP core already embeds faithfully into the coded QAM-over-ZF language.
\end{remark}

